\section{Statistik und Statistical Learning}

\subsection{Skalenniveaus und Charakterisierung von Merkmalen}

\textbf{Datentypen:}
\begin{itemize}[leftmargin=1.5em]
    \item \textbf{Strukturierte Daten:} Daten mit vordefinierter, strukturierter Form (Tabellen mit Zeilen und Spalten).
    \item \textbf{Unstrukturierte Daten:} Daten ohne identifizierbare formalisierte Struktur (Textdateien, Websites, Sensordaten, Bilder, Videos, Audiodateien, Emails, Social-Media-Daten).
\end{itemize}

\textbf{Herausforderungen der Data Science:}
\begin{itemize}[leftmargin=1.5em]
    \item Überführung der Flut von Rohdaten in verwertbare Informationen.
    \item Aufbereitung von unstrukturierten Rohdaten in strukturierte Form.
    \item Zusammenführung von Daten aus unterschiedlichen Quellen.
\end{itemize}

\definition{\textbf{Strukturierte Datenstruktur:}}
\begin{itemize}[leftmargin=1.5em]
    \item \textbf{Merkmalsträger (Fälle, Observations, Erhebungseinheiten)} $i = 1, \ldots, n$: Objekte, an denen Größen beobachtet oder erfasst werden (Zeilen eines Datensatzes), z.\,B.\ eine Person.
    \item \textbf{Merkmale (Variablen, Features)} $j = 1, \ldots, p$: An den Objekten betrachtete Charakteristika (Spalten eines Datensatzes), z.\,B.\ Körpergröße, Geschlecht.
    \item \textbf{Merkmalsausprägungen:} Werte eines Merkmalsträgers bezüglich eines Merkmals, z.\,B.\ $175$ cm, männlich.
    \item \textbf{Datenmatrix:} $n \times p$-Matrix, wobei Zeilen Merkmalsträger und Spalten Merkmale sind (in R: Data Frame).
\end{itemize}

\textbf{Klassifikation und Charakterisierung von Merkmalen:} Wesentlich für Datenanalyse, visuelle Darstellung und Wahl des statistischen Modells.

\definition{\textbf{Qualitative vs.\ Quantitative Variablen:}}
\begin{itemize}[leftmargin=1.5em]
    \item \textbf{Qualitativ (kategorial):} Begrenzte Anzahl möglicher Kategorien, z.\,B.\ Augenfarbe.
    \item \textbf{Quantitativ (numerisch):} Werte werden durch Zahlen dargestellt, z.\,B.\ Alter.
\end{itemize}

\definition{\textbf{Diskrete vs.\ Stetige Variablen:}}
\begin{itemize}[leftmargin=1.5em]
    \item \textbf{Diskret:} Endlich oder abzählbar unendlich viele Ausprägungen, z.\,B.\ Anzahl von Personen.
    \item \textbf{Stetig:} Unendlich viele Messwerte eines Intervalls (reelle Zahlen), z.\,B.\ Körpergröße.
    \item \textbf{Quasi-stetig:} Diskrete Merkmale mit vielen Ausprägungen (z.\,B.\ Einkommen) oder diskretisierte stetige Variablen (Klassenbildung).
\end{itemize}

\definition{\textbf{Manifeste vs.\ Latente Variablen:}}
\begin{itemize}[leftmargin=1.5em]
    \item \textbf{Manifest:} Direkt beobachtbar/messbar, z.\,B.\ Körpergröße.
    \item \textbf{Latent:} Nicht direkt beobachtbar; durch Indikatoren messbar gemacht (Konstrukte), z.\,B.\ Arbeitszufriedenheit, Kreativität.
\end{itemize}

\definition{\textbf{Unabhängige vs.\ Abhängige Variablen:}}
\begin{itemize}[leftmargin=1.5em]
    \item \textbf{Unabhängig:} Einflussgröße, Prädiktor.
    \item \textbf{Abhängig:} Zielgröße, Kriterium (z.\,B.\ in Regressionsanalysen).
\end{itemize}

\definition{\textbf{Skalenniveaus (Messniveaus):} Hierarchische Klassifikation, die definiert, welche mathematischen Operationen zulässig sind:}

\begin{enumerate}[leftmargin=2em]
    \item \textbf{Nominalskala:} Reine Klassifikation ohne Rangordnung. Zulässige Operationen: $=$, $\neq$. Beispiele: Geschlecht, Farbe, Nationalität. Lagemaß: Modus.
    
    \item \textbf{Ordinalskala:} Kategorien mit Rangordnung, aber kein einheitlicher Abstand zwischen Kategorien. Zulässige Operationen: $=$, $\neq$, $<$, $>$. Beispiele: Schulnoten, Likert-Skalen. Lagemaße: Modus, Median.
    
    \item \textbf{Intervallskala:} Ordnung mit einheitlichen Abständen, aber kein wahrer Nullpunkt (konventionell festgelegt). Zulässige Operationen: $=$, $\neq$, $<$, $>$, $+$, $-$. Beispiele: IQ, Temperatur (°C oder °F). Lagemaße: Modus, Median, Mittelwert.
    
    \item \textbf{Verhältnisskala:} Wie Intervallskala, aber mit wahrem Nullpunkt. Zulässige Operationen: $=$, $\neq$, $<$, $>$, $+$, $-$, $\cdot$, $\div$. Beispiele: Längenmaße, Preise, Temperatur (Kelvin). Lagemaße: Modus, Median, Mittelwert; Variationsmessungen vergleichbar.
    
    \item \textbf{Absolutskala:} Wahrer Nullpunkt mit natürlicher Maßeinheit (Stückzahl, Häufigkeiten). Alle mathematischen Operationen erlaubt. Beispiele: Anzahl von Personen, Stückzahl.
\end{enumerate}

\newpage

\subsection{Lage-, Streuungs- und Assoziationsmaße}

\textbf{Ziel:} Erste Eindrücke über die Verteilung der Daten durch Kennzahlen (Maßzahlen):


\textbf{Grafische Darstellung nach Skalenniveau und Analyseziel:}

\definition{\textbf{Grafische Darstellung von Häufigkeiten:}}
\begin{itemize}[leftmargin=1.5em]
    \item \textbf{Balkendiagramm (Bar Chart):} Für kategoriale (nominal, ordinal) und metrische Daten. Zeigt absolute oder relative Häufigkeiten. Variante: Gruppiertes oder gestapeltes Balkendiagramm für Vergleich mehrerer Kategorien.
    \item \textbf{Kreis- oder Ringdiagramm (Pie Chart):} Für kategoriale Daten mit kleiner Anzahl von Kategorien. Zeigt Anteile als Kreissegmente. Einsatz begrenzt wegen visueller Schwierigkeiten bei Vergleich von Winkeln.
    \item \textbf{Histogramm:} Für kontinuierliche metrische Daten. Zeigt Häufigkeitsverteilung in Klassen. Höhe der Balken proportional zur Häufigkeit/Dichte.
\end{itemize}

\definition{\textbf{Grafische Darstellung von Lage, Streuung und Form einer Verteilung:}}
\begin{itemize}[leftmargin=1.5em]
    \item \textbf{Boxplot (Box-and-Whisker Plot):} Zeigt Median, Quartile, Range und Ausreißer. Effizient zur visuellen Darstellung von Streuung und Schiefe. Einsatz bei ordinalen und metrischen Daten.
    \item \textbf{Violin-Plot:} Erweiterte Variante des Boxplots mit Kernel-Density-Estimation. Zeigt vollständige Verteilungsform zusätzlich zu Lagemaßen.
    \item \textbf{Histogramm mit Dichtekurve:} Kombination von Histogramm und geschätzter Dichtefunktion zur Visualisierung der Verteilungsform.
\end{itemize}

\definition{\textbf{Grafische Darstellung von Zusammenhängen (bivariate Analyse):}}
\begin{itemize}[leftmargin=1.5em]
    \item \textbf{Streudiagramm (Scatterplot):} Zwei metrische Variablen: Ein Punkt für jede Beobachtung $(x_i, y_i)$. Zeigt Korrelation und Trends. Erweiterung: Bubble Chart mit dritter Variable als Bubble-Größe.
    \item \textbf{Gruppiertes/Gestapeltes Balkendiagramm:} Zwei kategoriale Variablen oder eine kategoriale und eine metrische. Vergleicht Verteilungen zwischen Gruppen.
    \item \textbf{Boxplot nach Gruppen:} Metrische Variable versus kategoriale Gruppierungsvariable. Vergleicht Lagemaße und Streuung über Gruppen hinweg.
    \item \textbf{Liniendiagramm (Line Chart):} Für Zeitreihen oder ordinale Daten mit Trend. Zeigt Entwicklung über Zeit oder Kategorien.
    \item \textbf{Heatmap:} Zweidimensionale Darstellung von Intensitäten (z.\,B.\ Korrelationsmatrix). Farben kodieren Wertebereiche.
\end{itemize}

\textbf{Wahl des Diagrammtyps:} Hängt von Art der Variablen (kategorial/metrisch), Anzahl der Variablen und Analyseziel (Häufigkeiten, Verteilung, Zusammenhang) ab. Grundsatz: Maximale Informationsdichte mit minimaler visueller Komplexität.

\newpage

\subsection{Ähnlichkeitsmaße und Distanzmaße}

\definition{\textbf{Grundkonzepte:}}
\begin{itemize}[leftmargin=1.5em]
    \item \textbf{Distanz (Unähnlichkeit):} Maß für die Unähnlichkeit zweier Objekte. Hohe Distanz = niedrige Ähnlichkeit.
    \item \textbf{Ähnlichkeit (Similarity):} Maß für die Ähnlichkeit zweier Objekte. Hohe Ähnlichkeit = niedrige Distanz.
    \item Distanz und Ähnlichkeit sind inverse Größen: $s = 1 - d$ oder $d = 1 - s$ (normalisiert).
\end{itemize}

\definition{\textbf{Euklidische Distanz:}}
Für zwei $p$-dimensionale Punkte $\mathbf{x}_A$ und $\mathbf{x}_B$: $$d_{\text{Euclid}} = \sqrt{\sum_{j=1}^{p} (x_{B,j} - x_{A,j})^2} = \|\mathbf{x}_B - \mathbf{x}_A\|_2$$ Geometrische Interpretation: Luftlinienabstand im euklidischen Raum. \textbf{Nachteil:} Empfindlich gegenüber Ausreißern und unterschiedlichen Skalierungen -- Standardisierung der Variablen empfohlen.

\definition{\textbf{Manhattan-Distanz (City Block):}}
$$d_{\text{Manhattan}} = \sum_{j=1}^{p} |x_{B,j} - x_{A,j}|$$
Interpretation: Summe der Koordinatendifferenzen (wie Fortbewegung im Straßengitter). Robuster gegenüber Ausreißern als die euklidische Distanz.

\definition{\textbf{Minkowski-Distanz:}}
Verallgemeinerung mit Parameter $k$:
$$d_k(\mathbf{x}_A, \mathbf{x}_B) = \left(\sum_{j=1}^{p} |x_{B,j} - x_{A,j}|^k\right)^{1/k}$$ Spezialfälle: $k=1$ Manhattan; $k=2$ Euklidisch; $k \to \infty$ Chebyshev-Distanz ($\max_j |x_{A,j} - x_{B,j}|$).

\definition{\textbf{Mahalanobis-Distanz:}}
$$d_{\text{Mahal}} = \sqrt{(\mathbf{x}_B - \mathbf{x}_A)^T \hat{\boldsymbol{\Sigma}}^{-1} (\mathbf{x}_B - \mathbf{x}_A)}$$
wobei $\hat{\boldsymbol{\Sigma}}$ die aus den Daten \emph{geschätzte} Kovarianzmatrix ist. \textbf{Vorteil:} Berücksichtigt Korrelationen zwischen Variablen und unterschiedliche Skalierungen -- keine separate Standardisierung nötig. \textbf{Nachteil:} Erfordert invertierbare Kovarianzmatrix; instabil bei $p > n$.

\definition{\textbf{Kosinus-Ähnlichkeit:}}
$$\cos(\theta) = \frac{\mathbf{x}_A \cdot \mathbf{x}_B}{\|\mathbf{x}_A\|_2 \|\mathbf{x}_B\|_2}$$ Wertebereich: $[-1, 1]$; höhere Werte = höhere Ähnlichkeit. Misst den Winkel zwischen zwei Vektoren, \textbf{unabhängig von ihrer Länge}. Verwendung: Text- und Bildvektoren (z.\,B.\ TF-IDF-Vektoren, Word Embeddings).

\definition{\textbf{Tanimoto-Koeffizient (Jaccard-Ähnlichkeit):}}
Für \textbf{binäre} Vektoren oder Mengen $A, B \subseteq \{1,\ldots,p\}$: $$S_{\text{Tanimoto}} = \frac{|A \cap B|}{|A \cup B|} = \frac{\sum_{j} \mathbf{1}[x_{A,j}=1 \wedge x_{B,j}=1]}{\sum_{j} \mathbf{1}[x_{A,j}=1 \vee x_{B,j}=1]}$$ Wertebereich: $[0, 1]$. \textit{Hinweis:} Diese Formel gilt exakt für Binärvektoren. Für kontinuierliche Vektoren existiert eine verallgemeinerte Tanimoto-Formel ($\frac{\mathbf{x}_A \cdot \mathbf{x}_B}{\|\mathbf{x}_A\|^2 + \|\mathbf{x}_B\|^2 - \mathbf{x}_A \cdot \mathbf{x}_B}$), die jedoch seltener verwendet wird. Anwendung: Genetik (Markerpräsenz), Chemoinformatik (Molekülstruktur-Vergleich).

\definition{\textbf{Dice-Koeffizient:}}
$$S_{\text{Dice}} = \frac{2|A \cap B|}{|A| + |B|}$$ Ähnlich wie Tanimoto, gewichtet die Überschneidung stärker.
Zusammenhang: $S_{\text{Dice}} = \frac{2 S_{\text{Tanimoto}}}{1 + S_{\text{Tanimoto}}}$. Häufig in Ökologie, Biologie und NLP (Token-Überlappung).

\definition{\textbf{Korrelationsbasierte Distanzmaße:}}
\begin{itemize}[leftmargin=1.5em]
    \item \textbf{Pearson-Distanz:} $d = 1 - r$, wobei $r$ die Pearson-Korrelation ist. Misst \emph{lineare} Abhängigkeit; invariant gegenüber linearen Transformationen.
    \item \textbf{Spearman-Distanz:} $d = 1 - r_S$, wobei $r_S$ die Spearman-Rangkorrelation ist. Robuster gegenüber Ausreißern und nicht-linearen monotonen Zusammenhängen.
\end{itemize}

\textbf{Wahl des Distanzmaßes:}
\begin{table}[h]
    \centering
    \small
    \begin{tabular}{|l|l|}
        \hline
        \textbf{Situation} & \textbf{Empfohlenes Maß} \\
        \hline
        Normalisierte numerische Daten & Euklidisch \\
        \hline
        Korrelierte Variablen & Mahalanobis \\
        \hline
        Hochdimensionale Text-/Bildvektoren & Kosinus \\
        \hline
        Binäre Merkmalsvektoren & Tanimoto/Dice \\
        \hline
        Rangbasierte oder robuste Analyse & Spearman-Distanz \\
        \hline
        Ausreißer vorhanden & Manhattan oder Mahalanobis \\
        \hline
    \end{tabular}
\end{table}

\subsection{Diagrammtypen und Einsatzgebiete}
\newpage

\subsection{Wahrscheinlichkeiten und Wahrscheinlichkeitsräume}

\definition{\textbf{Zufallsexperiment (Random Experiment):} Ein Prozess, dessen Ausgang nicht im Voraus bekannt ist, aber alle möglichen Ausgänge sind bekannt. Genau eines dieser Ergebnisse (ein sogenanntes Elementarereignis $\omega$) wird mit Sicherheit eintreten.}

\textbf{Wahrscheinlichkeitsinterpretationen:}
\begin{itemize}[leftmargin=1.5em]
    \item \textbf{Frequentistische Sicht:} Das Zufallsexperiment wird sehr häufig (tatsächlich oder gedanklich) unter identischen Bedingungen wiederholt. Die relative Häufigkeit, mit der ein Ereignis auftritt, ist seine Wahrscheinlichkeit. Diese Definition ist beobacherunabhängig (objektiv).
    \item \textbf{Bayessche Sicht:} Wahrscheinlichkeit als Grad der Gewissheit oder subjektives Vertrauen basierend auf verfügbarem Wissen. Beispiele: ``Ich werde das Examen wahrscheinlich bestehen'' oder ``Es besteht eine 99\%-Wahrscheinlichkeit, dass der Fingerabdruck von Person X stammt''.
\end{itemize}

\definition{\textbf{Stichprobenraum (Ergebnisraum):} Die Menge aller möglichen Ausgänge eines Zufallsexperiments:
\formel{\Omega = \{\omega \,|\, \omega \text{ ist ein Elementarereignis}\}}
$\Omega$ kann endlich, abzählbar unendlich oder überabzählbar unendlich sein.}

\definition{\textbf{Ereignis:} Eine Teilmenge $E \subseteq \Omega$. Ein Ereignis tritt ein, wenn der tatsächliche Ausgang des Zufallsexperiments ein Element von $E$ ist.}

\definition{\textbf{Ereignisraum ($\sigma$-Algebra):} Die Menge $\mathcal{A}$ aller Ereignisse, denen wir Wahrscheinlichkeiten zuordnen möchten:}
\formel{\mathcal{A} = \{E \,|\, E \subseteq \Omega \text{ ist ein Ereignis}\}}
Die $\sigma$-Algebra muss folgende Eigenschaften erfüllen:
\begin{enumerate}[leftmargin=1.5em]
    \item Sie enthält das sichere Ereignis $\Omega$.
    \item Für jedes Ereignis $E \in \mathcal{A}$ enthält sie auch das komplementäre Ereignis $E^c = \Omega \setminus E$.
    \item Für jede abzählbare Folge von Ereignissen $E_1, E_2, \ldots \in \mathcal{A}$ enthält sie auch deren Vereinigung: $\bigcup_{i=1}^\infty E_i \in \mathcal{A}$.
\end{enumerate}

\definition{\textbf{Wahrscheinlichkeitsfunktion (Wahrscheinlichkeitsmaß):} Eine Funktion $P: \mathcal{A} \to [0, 1]$, die jedem Ereignis $E \in \mathcal{A}$ eine Wahrscheinlichkeit $P(E)$ zuordnet.}

\textbf{Kolmogorov-Axiome:}
\begin{enumerate}[leftmargin=1.5em]
    \item $P(\Omega) = 1$ (das sichere Ereignis hat Wahrscheinlichkeit 1).
    \item Für alle $E \in \mathcal{A}$ gilt $P(E) \geq 0$ (Wahrscheinlichkeiten sind nicht-negativ).
    \item \textbf{Additivität für disjunkte Ereignisse:} Für zwei disjunkte Ereignisse $E_1 \cap E_2 = \emptyset$ gilt:
    \formel{P(E_1 \cup E_2) = P(E_1) + P(E_2)}
    Diese Eigenschaft gilt auch für abzählbar unendlich viele paarweise disjunkte Ereignisse ($E_i \cap E_j = \emptyset$ für $i \neq j$):
    \formel{P\left(\bigcup_{i=1}^\infty E_i\right) = \sum_{i=1}^\infty P(E_i)}
\end{enumerate}

\definition{\textbf{Wahrscheinlichkeitsraum:} Das Tripel $(\Omega, \mathcal{A}, P)$ bestehend aus dem Stichprobenraum $\Omega$, einer $\sigma$-Algebra $\mathcal{A}$ und einer Wahrscheinlichkeitsfunktion $P$ wird als Wahrscheinlichkeitsraum bezeichnet.}

\definition{\textbf{Laplace'scher Wahrscheinlichkeitsraum:} Ein endlicher Wahrscheinlichkeitsraum, in dem alle Elementarereignisse gleiche Wahrscheinlichkeit haben (Indifferenzprinzip).

Für einen Laplace-Raum mit $m$ Elementarereignissen gilt:
\formel{P(\{\omega\}) = \frac{1}{m} \quad \text{für alle } \omega \in \Omega}

Für ein beliebiges Ereignis $E$ mit $g$ Elementen:
\formel{P(E) = \frac{g}{m} = \frac{\text{Anzahl der günstigen Fälle}}{\text{Anzahl der möglichen Fälle}}}

\textbf{Beispiel (Würfeln):} Beim Werfen eines fairen Würfels ist $\Omega = \{1, 2, 3, 4, 5, 6\}$, $m = 6$ und für das Ereignis ``gerade Zahl'' $E = \{2, 4, 6\}$ mit $g = 3$ erhalten wir $P(E) = \frac{3}{6} = \frac{1}{2}$.}

\textbf{Abgeleitete Wahrscheinlichkeitsregeln:}
\begin{itemize}[leftmargin=1.5em]
    \item \textbf{Komplementärregel:} $P(E^c) = 1 - P(E)$
    \item \textbf{Additionsregel:} $P(E_1 \cup E_2) = P(E_1) + P(E_2) - P(E_1 \cap E_2)$
    \item \textbf{Monotonie:} Wenn $E_1 \subseteq E_2$, dann $P(E_1) \leq P(E_2)$.
\end{itemize}

\newpage





\subsection{Kombinatorik}

\definition{\textbf{Kombinatorik:} Die Lehre vom Abzählen endlicher Mengen und ihren Teilmengen. Sie ist fundamental für Wahrscheinlichkeitsberechnungen in \textbf{Laplace-Räumen}: Da $P(E) = \frac{\text{günstige Fälle}}{\text{mögliche Fälle}}$, muss man oft zählen, wie viele Anordnungen oder Auswahlen möglich sind -- genau das leistet die Kombinatorik.}

\definition{\textbf{Permutation ohne Wiederholung:} Die Anzahl der Möglichkeiten, $n$ \emph{unterschiedliche} Objekte in eine Reihenfolge zu bringen: $$P(n) = n! = n \cdot (n-1) \cdot \ldots \cdot 2 \cdot 1$$ \textbf{Beispiel:} 3 Personen können auf $3! = 6$ Arten angeordnet werden.}

\definition{\textbf{Permutation mit Wiederholung:} Wenn unter $n$ Objekten $k$ Gruppen mit $n_1, n_2, \ldots, n_k$ identischen Objekten existieren ($n_1 + \ldots + n_k = n$): $$P(n;\, n_1, n_2, \ldots, n_k) = \frac{n!}{n_1!\cdot n_2!\cdot\ldots\cdot n_k!}$$ \textbf{Beispiel:} Die Buchstaben in ``AABBCC'' können auf $\frac{6!}{2!\cdot 2!\cdot 2!} = 90$ Arten angeordnet werden.}

\definition{\textbf{Variation ohne Wiederholung:} Die Anzahl der Möglichkeiten, aus $n$ unterschiedlichen Objekten genau $k$ auszuwählen \emph{und} in eine Reihenfolge zu bringen ($k \leq n$): $$V(n, k) = \frac{n!}{(n-k)!} = n \cdot (n-1) \cdot \ldots \cdot (n-k+1)$$ \textit{Abgrenzung zur Permutation:} Eine Permutation ist der Spezialfall $k = n$; eine Variation wählt nur $k < n$ Elemente aus und ordnet diese an.\\ $n$ unterschiedlichen Objekten genau $k$ auszuwählen \emph{und} in eine Reihenfolge zu bringen ($k \leq n$): $$V(n, k) = \frac{n!}{(n-k)!} = n \cdot (n-1) \cdot \ldots \cdot (n-k+1)$$ \textit{Abgrenzung zur Permutation:} Eine Permutation ist der Spezialfall $k = n$; eine Variation wählt nur $k < n$ Elemente aus und ordnet diese an.\\ \textbf{Beispiel:} Aus 10 Pferden die ersten 3 Plätze bei einem Rennen: $$V_w(n, k) = n^k$$ \textbf{Beispiel:} 3-stellige Zahlencodes aus Ziffern 0--9:
$V_w(10,3) = 10^3 = 1000$ Möglichkeiten.}

\definition{\textbf{Kombination ohne Wiederholung (Binomialkoeffizient):} Die Anza Positionen aus einem Repertoire von $n$ verschiedenen Elementen zu füllen, wobei Wiederholungen erlaubt sind: $$V_w(n, k) = n^k$$ \textbf{Beispiel:} 3-stellige Zahlencodes aus Ziffern 0--9: 

\begin{itemize}[leftmargin=1.5em]
    \item $\binom{n}{k} = \binom{n}{n-k}$ (Symmetrie)
    \item $\binom{n}{0} = \binom{n}{n} = 1$
    \item $\binom{n}{k} + \binom{n}{k+1} = \binom{n+1}{k+1}$ (Pascal'sche Identität)
\end{itemize}
\textbf{Beispiel:} 5-Karten-Hände aus einem 52-Karten-Deck: $\binom{52}{5} = 2\,598\,960$ Möglichkeiten.}

\definition{\textbf{Kombination mit Wiederholung:} Die Anzahl der Möglichkeiten, $k$ Elemente aus $n$ verschiedenen Typen auszuwählen, wobei jeder Typ beliebig oft verwendet werden kann und die Reihenfolge nicht zählt: $$C_w(n, k) = \binom{n+k-1}{k} = \frac{(n+k-1)!}{k!\cdot(n-1)!}$$ \textbf{Beispiel:} Auswahl von 3 Früchten aus 5 Sorten mit Wiederholung: $\binom{7}{3} = 35$ Möglichkeiten.}

\textbf{Zusammenfassende Übersicht:}
\begin{table}[h]
    \centering
    \small
    \begin{tabular}{|l|c|c|}
        \hline
        \textbf{Art} & \textbf{Ohne Wiederholung} & \textbf{Mit Wiederholung} \\
        \hline
        \textbf{Permutation} (alle $n$ Elemente anordnen) &
        $n!$ & $\dfrac{n!}{n_1!\cdots n_k!}$ \\[6pt]
        \hline
        \textbf{Variation} ($k$ aus $n$, Reihenfolge zählt) &
        $\dfrac{n!}{(n-k)!}$ & $n^k$ \\[6pt]
        \hline
        \textbf{Kombination} ($k$ aus $n$, Reihenfolge zählt nicht) &
        $\dbinom{n}{k}$ & $\dbinom{n+k-1}{k}$ \\[6pt]
        \hline
    \end{tabular}
\end{table}

\textbf{Entscheidungsbaum zur Formelwahl:}
\begin{itemize}[leftmargin=1.5em]
    \item Werden \emph{alle} $n$ Objekte angeordnet? $\to$ \textbf{Permutation}
    \item Werden nur $k < n$ Objekte ausgewählt?
    \begin{itemize}[leftmargin=1.5em]
        \item Zählt die Reihenfolge? $\to$ \textbf{Variation}
        \item Zählt die Reihenfolge \emph{nicht}? $\to$ \textbf{Kombination}
    \end{itemize}
    \item Sind Wiederholungen erlaubt? $\to$ jeweilige ``mit Wiederholung''-Formel
\end{itemize}

\subsection{Diskrete und stetige Zufallsvariablen}
\newpage

\subsection{Statistische Tests}

\definition{\textbf{Statistischer Test (Hypothesentest):} Ein Verfahren, mit dem anhand von Stichprobendaten eine Entscheidung darüber getroffen wird, ob eine Vermutung (Hypothese) über einen Parameter oder eine Verteilung einer Grundgesamtheit haltbar ist.}

\textbf{Grundkonzept der Hypothesen:}

\definition{\textbf{Nullhypothese ($H_0$):} Die ``Standardannahme'' oder konservative Behauptung, die wir zunächst für wahr halten. Oft besagt sie, dass kein Effekt, keine Veränderung oder keine Abweichung vorhanden ist.

\textbf{Alternativhypothese ($H_1$ oder $H_a$):} Die Gegenhypothese, die das Gegenteil von $H_0$ behauptet und auf deren Gültigkeit wir testen möchten.}

\beispiel{
\begin{enumerate}[leftmargin=2em]
    \item \textbf{Qualitätskontrolle:} Ein Lieferant verspricht maximal 1,5\% fehlerhafte Komponenten. Bei Stichprobe $n=50$ mit 1 Fehler testen wir:
    \begin{itemize}
        \item $H_0: p \leq 0{,}015$ (Lieferant hält sein Versprechen)
        \item $H_1: p > 0{,}015$ (Lieferant bricht sein Versprechen) — einseitiger Test
    \end{itemize}
    
    \item \textbf{Produktqualität:} Ein Hersteller behauptet, sein Produkt enthält durchschnittlich höchstens 10g Zucker pro 100g. Mit Stichprobe $n=10$: $\bar{x}=11{,}4$, $s=1{,}2$g testen wir:
    \begin{itemize}
        \item $H_0: \mu \leq 10$ (Hersteller hat recht)
        \item $H_1: \mu > 10$ (Hersteller täuscht) — einseitiger Test
    \end{itemize}
\end{enumerate}
}

\textbf{Einteilung von Hypothesen:}
\begin{itemize}[leftmargin=1.5em]
    \item \textbf{Einfache Hypothese:} Spezifiziert genau einen Parameterwert, z.\,B.\ $H_0: \mu = 10$.
    \item \textbf{Zusammengesetzte Hypothese:} Spezifiziert einen Wertebereich, z.\,B.\ $H_0: \mu \leq 10$.
    \item \textbf{Einseitiger Test:} $H_1$ liegt auf einer Seite von $H_0$, z.\,B.\ $H_1: \mu > 10$ oder $H_1: \mu < 10$.
    \item \textbf{Zweiseitiger Test:} $H_1$ liegt auf beiden Seiten, z.\,B.\ $H_1: \mu \neq 10$.
\end{itemize}

\definition{\textbf{Teststatistik ($T$):} Eine nach einer festgelegten Vorschrift aus den Stichprobendaten berechnete Zufallsvariable, deren Verteilung unter $H_0$ bekannt ist. Unterschiedliche Tests nutzen unterschiedliche Teststatistiken (z.\,B.\ $z$-, $t$-, $\chi^2$-, $F$-Statistik).}

\definition{\textbf{Signifikanzniveau ($\alpha$):} Die vorher festgelegte maximale Wahrscheinlichkeit, einen Fehler 1.\ Art zu begehen. Typische Werte: $\alpha \in \{0{,}01, 0{,}05, 0{,}10\}$.}

\textbf{Ablauf eines klassischen Hypothesentests:}

\begin{enumerate}[leftmargin=1.5em]
    \item \textbf{Hypothesen formulieren:} $H_0$ und $H_1$ sowie Signifikanzniveau $\alpha$ festlegen.
    \item \textbf{Teststatistik berechnen:} Aus der Stichprobe den beobachteten Wert $t_{\text{obs}}$ der Teststatistik berechnen.
    \item \textbf{Kritischer Bereich bestimmen:} Basierend auf $\alpha$ und der Verteilung von $T$ unter $H_0$ den kritischen/Ablehnungsbereich bestimmen. Außerhalb liegt der Annahmebereich.
    \item \textbf{Testentscheidung treffen:} 
    \begin{itemize}
        \item Liegt $t_{\text{obs}}$ im Ablehnungsbereich $\Rightarrow$ $H_0$ ablehnen, $H_1$ annehmen.
        \item Liegt $t_{\text{obs}}$ im Annahmebereich $\Rightarrow$ $H_0$ nicht ablehnen.
    \end{itemize}
\end{enumerate}

\definition{\textbf{p-Wert (p-value):} Die bedingte Wahrscheinlichkeit, unter $H_0$ einen Wert für die Teststatistik zu beobachten, der mindestens so extrem ist wie der beobachtete (in Richtung $H_1$):
\formel{p = P(T \geq t_{\text{obs}} \mid H_0 \text{ ist wahr})}
Der p-Wert ist eine ``überschreitungswahrscheinlichkeit''. 

\textbf{Interpretation:} 
\begin{itemize}[leftmargin=1.5em]
    \item Ist $p < \alpha$ $\Rightarrow$ Ablehnung von $H_0$ (Ergebnis ist statistisch signifikant).
    \item Ist $p \geq \alpha$ $\Rightarrow$ $H_0$ nicht ablehnen (kein Beweis gegen $H_0$).
\end{itemize}

\textbf{Wichtig:} Der p-Wert ist \textbf{nicht} die Wahrscheinlichkeit, dass $H_0$ wahr ist! Sondern: Wie wahrscheinlich ist das beobachtete (oder extremere) Ergebnis, wenn $H_0$ stimmt?}

\definition{\textbf{Fehler beim Hypothesentest:}

\begin{center}
\begin{tabular}{|l|c|c|}
    \hline
    & $H_0$ tatsächlich wahr & $H_0$ tatsächlich falsch \\
    \hline
    $H_0$ nicht ablehnen & Korrekte Entscheidung (Konfidenz: $1-\alpha$) & \textbf{Fehler II.\ Art} ($\beta$) \\
    \hline
    $H_0$ ablehnen & \textbf{Fehler I.\ Art} ($\alpha$, Signifikanzniveau) & Korrekte Entscheidung (Power: $1-\beta$) \\
    \hline
\end{tabular}
\end{center}

\begin{itemize}[leftmargin=1.5em]
    \item \textbf{Fehler I.\ Art (Fehler erster Art, Type I Error):} $H_0$ wird abgelehnt, obwohl $H_0$ tatsächlich wahr ist (``Fehlalarm''). Wahrscheinlichkeit: $P(\text{Fehler I}) = \alpha$.
    
    \item \textbf{Fehler II.\ Art (Fehler zweiter Art, Type II Error):} $H_0$ wird nicht abgelehnt, obwohl $H_1$ tatsächlich wahr ist (``falsch negativ''). Wahrscheinlichkeit: $P(\text{Fehler II}) = \beta$.
    
    \item \textbf{Teststärke (Power):} Die Wahrscheinlichkeit, die Nullhypothese richtig abzulehnen, wenn $H_1$ wahr ist: $\text{Power} = 1 - \beta$. Eine hohe Power ist wünschenswert.
\end{itemize}

\textbf{Zusammenhang:} $\alpha$ und $\beta$ sind negativ korreliert. Für feste Stichprobengröße: Je kleiner $\alpha$, desto größer $\beta$. Um beide zu reduzieren, muss die Stichprobengröße erhöht werden.}

\textbf{Auswahl zwischen ein- und zweiseitigen Tests:} Ein einseitiger Test ist nur rechtfertigt, wenn die Richtung des möglichen Effekts vor der Datenerhebung klar definiert ist. Ein zweiseitiger Test ist konservativer und allgemeiner einsetzbar.

\newpage

\subsection{Population, Stichproben und Konfidenzintervalle}

\definition{\textbf{Grundgesamtheit (Population):} Die Menge aller statistischen Einheiten (Objekte), über die man Aussagen gewinnen möchte. Sie folgt einer zugrundeliegenden, aber unbekannten Verteilung (z.\,B.\ alle Wähler eines Landes, alle produzierten Komponenten). \textbf{Stichprobe:} Die tatsächlich untersuchte Teilmenge der Grundgesamtheit, gewöhnlich als Zufallsstichprobe gezogen. Daten aus der Stichprobe ermöglichen es, auf Charakteristika der Grundgesamtheit zu schließen.}

\textbf{Inferenzprinzip:} Von der empirischen Verteilung der Stichprobe auf die wahre, unbekannte Verteilung der Grundgesamtheit schließen.

\definition{\textbf{Punktschätzer:} Eine Statistik, die einen Parameter der Grundgesamtheit mit einem einzelnen Zahlenwert schätzt. Beispiele: Stichprobenmittelwert $\bar{X} = \frac{1}{n}\sum_{i=1}^n X_i$ schätzt den Populationsmittelwert $\mu$; Stichprobenvarianz $S^2 = \frac{1}{n-1}\sum_{i=1}^n (X_i - \bar{X})^2$ schätzt die Populationsvarianz $\sigma^2$; Stichprobenanteil $\hat{p} = \frac{\text{Anzahl Erfolge}}{n}$ schätzt den Populationsanteil $p$. \textbf{Problem:} Da die Stichprobe zufällig ist, ist die Wahrscheinlichkeit, den wahren Parameterwert exakt zu treffen, im kontinuierlichen Fall gleich Null.}

\definition{\textbf{Konfidenzintervall (Vertrauensintervall):} Ein Bereich (typischerweise ein Intervall), in dem der unbekannte Parameter der Grundgesamtheit mit hoher Sicherheit liegt. Formal: Ein $(1-\alpha)$-Konfidenzintervall für einen Parameter $\theta$ ist ein zufälliges Intervall $[\theta_l, \theta_u]$ mit: $P(\theta_l \leq \theta \leq \theta_u) = 1 - \alpha$. \textbf{Interpretation:} Wenn wir das Experiment (Ziehen einer Stichprobe und Konstruktion eines KI) sehr häufig wiederholen würden, würde der wahre Parameter in $(1-\alpha) \cdot 100\%$ aller Konfidenzintervalle liegen. \textbf{Wichtig:} Das Konfidenzintervall ist zufällig (hängt von der Stichprobe ab), der Parameter $\theta$ ist fest (wenn auch unbekannt). Der Parameter liegt entweder im Intervall oder nicht — die Wahrscheinlichkeit $1-\alpha$ bezieht sich auf wiederholte Stichproben. \textbf{Konfidenzniveau:} Üblicherweise wird ein Konfidenzniveau von $1-\alpha = 95\%$ (oder $90\%$, $99\%$) gewählt, entsprechend $\alpha = 0{,}05$ (oder $0{,}10$, $0{,}01$).}

\textbf{Arten von Konfidenzintervallen:}

\definition{\textbf{Zweiseitiges Konfidenzintervall:} Der Parameter $\theta$ kann außerhalb des Intervalls auf beiden Seiten liegen, mit gleicher Wahrscheinlichkeit:
\formel{P(\theta < \theta_l) = P(\theta > \theta_u) = \frac{\alpha}{2}}

Die Grenzen werden mit den Quantilen für $\frac{\alpha}{2}$ und $1 - \frac{\alpha}{2}$ bestimmt.}

\definition{\textbf{Einseitige Konfidenzintervalle:}}
\begin{itemize}[leftmargin=1.5em]
    \item \textbf{Untere einseitige KI:} $[\theta_l, +\infty)$ mit $P(\theta_l \leq \theta) = 1 - \alpha$. Obere Grenze ist unendlich. Quantil für $1-\alpha$ erforderlich.
    \item \textbf{Obere einseitige KI:} $(-\infty, \theta_u]$ mit $P(\theta \leq \theta_u) = 1 - \alpha$. Untere Grenze ist $-\infty$. Quantil für $1-\alpha$ erforderlich.
\end{itemize}

\definition{\textbf{Konfidenzintervall für den Erwartungswert $\mu$ (bekannte Varianz $\sigma_0^2$):} \textbf{Voraussetzungen:} Entweder: Population ist normalverteilt, oder $n$ ist groß genug ($n \geq 30$), dass der Zentrale Grenzwertsatz angewendet werden kann. Nach dem Zentralen Grenzwertsatz ist: $\bar{X} \sim N\left(\mu, \frac{\sigma_0^2}{n}\right) \quad \Rightarrow \quad Z = \frac{\bar{X} - \mu}{\sigma_0/\sqrt{n}} \sim N(0,1)$. Das $(1-\alpha)$-Konfidenzintervall für $\mu$: $\left[\bar{X} - z_{1-\alpha/2} \cdot \frac{\sigma_0}{\sqrt{n}}, \quad \bar{X} + z_{1-\alpha/2} \cdot \frac{\sigma_0}{\sqrt{n}}\right]$ wobei $z_{1-\alpha/2}$ das $(1-\alpha/2)$-Quantil der Standardnormalverteilung ist (für $\alpha=0{,}05$: $z_{0{,}975} \approx 1{,}96$).}

\definition{\textbf{Konfidenzintervall für den Erwartungswert $\mu$ (unbekannte Varianz):} Wenn $\sigma^2$ unbekannt ist, verwenden wir die Stichprobenvarianz $S^2 = \frac{1}{n-1}\sum_{i=1}^n (X_i - \bar{X})^2$ und die $t$-Verteilung mit $n-1$ Freiheitsgraden: $T = \frac{\bar{X} - \mu}{S/\sqrt{n}} \sim t_{n-1}$. Das $(1-\alpha)$-Konfidenzintervall für $\mu$: $\left[\bar{X} - t_{n-1, 1-\alpha/2} \cdot \frac{S}{\sqrt{n}}, \quad \bar{X} + t_{n-1, 1-\alpha/2} \cdot \frac{S}{\sqrt{n}}\right]$ wobei $t_{n-1, 1-\alpha/2}$ das $(1-\alpha/2)$-Quantil der $t$-Verteilung mit $n-1$ Freiheitsgraden ist. \textbf{Beachte:} Die $t$-Verteilung hat schwerere Schwänze als die Normalverteilung, besonders für kleine $n$, was zu breiteren Konfidenzintervallen führt.}

\definition{\textbf{Konfidenzintervall für die Varianz $\sigma^2$ (normalverteilte Population):}

Die Teststatistik:
\formel{Q = \frac{(n-1)S^2}{\sigma^2} \sim \chi^2_{n-1}}

ist $\chi^2$-verteilt mit $n-1$ Freiheitsgraden.

Das $(1-\alpha)$-Konfidenzintervall für $\sigma^2$:
\formel{\left[\frac{(n-1)S^2}{\chi^2_{n-1, 1-\alpha/2}}, \quad \frac{(n-1)S^2}{\chi^2_{n-1, \alpha/2}}\right]}

wobei $\chi^2_{n-1, p}$ das $p$-Quantil der $\chi^2$-Verteilung mit $n-1$ Freiheitsgraden ist.} Die Teststatistik: $Q = \frac{(n-1)S^2}{\sigma^2} \sim \chi^2_{n-1}$ ist $\chi^2$-verteilt mit $n-1$ Freiheitsgraden. Das $(1-\alpha)$-Konfidenzintervall für $\sigma^2$: $\left[\frac{(n-1)S^2}{\chi^2_{n-1, 1-\alpha/2}}, \quad \frac{(n-1)S^2}{\chi^2_{n-1, \alpha/2}}\right]$ 

\subsection{Zentraler Grenzwertsatz und Gesetz der großen Zahlen}

\textbf{Zwei fundamentale Theoreme der Wahrscheinlichkeitstheorie:}

\definition{\textbf{Unabhängig und identisch verteilt (u.i.v., iid):} Wenn ein Zufallsexperiment — dargestellt durch eine Zufallsvariable $X$ mit Erwartungswert $\mu$, Varianz $\sigma^2$ und Verteilungsfunktion $F(x)$ — unabhängig $n$-mal wiederholt wird, so bezeichnen wir das Ergebnis beim $i$-ten Versuch mit $X_i$ für $i = 1, \ldots, n$. Die Zufallsvariablen $X_1, \ldots, X_n$ sind: unabhängig voneinander; identisch verteilt (haben dieselbe Verteilung wie $X$); haben alle denselben Erwartungswert $\mu$ und dieselbe Varianz $\sigma^2$. Man kann $X_1, \ldots, X_n$ als unabhängige ``Kopien'' von $X$ auffassen. Die konkreten Zahlenwerte $x_1, \ldots, x_n$ heißen Realisierungen.}

\textbf{\textbf{---} Ungleichungen mit praktischer Bedeutung \textbf{---}}

\definition{\textbf{Markov-Ungleichung:} Für eine nicht-negative Zufallsvariable $X$ (diskret oder stetig) und alle $a > 0$ gilt: $P(X \geq a) \leq \frac{\mathbb{E}[X]}{a}$. \textbf{Interpretation:} Die Wahrscheinlichkeit, dass eine nicht-negative Zufallsvariable einen Wert $\geq a$ annimmt, ist höchstens das Verhältnis ihres Erwartungswerts zu $a$. Dies gibt eine modellunabhängige obere Schranke. \textbf{Beispiel:} Wenn eine Zufallsvariable $X$ (z.\,B.\ Einkommen) einen Erwartungswert von 30.000 EUR hat, ist die Wahrscheinlichkeit, dass $X \geq 90.000$ EUR, höchstens $\frac{30.000}{90.000} = \frac{1}{3}$.}

\definition{\textbf{Tschebyscheff-Ungleichung (Chebyshev's Inequality):} Für eine Zufallsvariable $X$ mit Erwartungswert $\mu = \mathbb{E}[X]$ und Varianz $\sigma^2 = \text{Var}(X)$ gilt für alle $c > 0$: Variante 1: $P(|X - \mu| \geq c) \leq \frac{\sigma^2}{c^2}$. Variante 2: $P(|X - \mu| < c) \geq 1 - \frac{\sigma^2}{c^2}$. \textbf{Interpretation:} Die Wahrscheinlichkeit, dass eine Zufallsvariable mehr als $c$ Einheiten von ihrem Mittelwert abweicht, ist durch das Verhältnis ihrer Varianz zu $c^2$ beschränkt — \textbf{unabhängig von der konkreten Verteilung} von $X$. \textbf{Vereinfachte Form (mit Standard-Abweichungen):} Wenn $c = k \cdot \sigma$ (Abweichung in Einheiten der Standardabweichung): $P(|X - \mu| \geq k \cdot \sigma) \leq \frac{1}{k^2}$ oder $P(|X - \mu| < k \cdot \sigma) \geq 1 - \frac{1}{k^2}$. \textbf{Beispiele:} $k=2$: Mindestens $75\%$ der Daten liegen im Intervall $[\mu - 2\sigma, \mu + 2\sigma]$. $k=3$: Mindestens $88{,}9\%$ der Daten liegen im Intervall $[\mu - 3\sigma, \mu + 3\sigma]$. \textbf{Beweis-Skizze:} Definiere eine Hilfsvariable $Y = 1$ wenn $|X - \mu| \geq c$ und $Y = 0$ sonst. Dann $\mathbb{E}[Y] = P(|X - \mu| \geq c)$ und $Y \leq \frac{(X-\mu)^2}{c^2}$, daher $\mathbb{E}[Y] \leq \mathbb{E}\left[\frac{(X-\mu)^2}{c^2}\right] = \frac{\sigma^2}{c^2}$.}

\textbf{\textbf{---} Die zwei großen Grenzwertsätze \textbf{---}}

\definition{\textbf{Gesetz der großen Zahlen (Law of Large Numbers, LLN):} \textbf{Theorem:} Für u.i.v.\ Zufallsvariablen $X_1, X_2, \ldots, X_n$ mit endlichem Erwartungswert $\mu$ und endlicher Varianz $\sigma^2$ gilt: $P\left(\left|\bar{X}_n - \mu\right| < c\right) \to 1 \quad \text{für } n \to \infty$ für jedes beliebige $c > 0$. Equivalente Formulierung: Der Stichprobenmittelwert $\bar{X}_n = \frac{1}{n}(X_1 + X_2 + \cdots + X_n)$ \textbf{konvergiert in Wahrscheinlichkeit} gegen den Populationsmittelwert $\mu$: $\bar{X}_n \xrightarrow{P} \mu$. \textbf{Beweis (via Tschebyscheff):} $\mathbb{E}[\bar{X}_n] = \mu$ (unverzerrter Schätzer); $\text{Var}(\bar{X}_n) = \frac{\sigma^2}{n}$ (Varianz nimmt mit $n$ ab); Nach Tschebyscheff: $P\left(\left|\bar{X}_n - \mu\right| \geq c\right) \leq \frac{\sigma^2}{nc^2} \to 0$ für $n \to \infty$. \textbf{Anschauliche Interpretation:} Mit genügend großem Stichprobenumfang ist die Wahrscheinlichkeit beliebig groß, dass der Stichprobenmittelwert nahe beim wahren Populationsmittelwert liegt. Das ist die mathematische Rechtfertigung für den klassischen Ansatz: ``Große Stichproben liefern verlässliche Schätzungen''. \textbf{Bernoulli's Theorem (Spezialfall):} Wenn wir ein Ereignis $A$ (mit Wahrscheinlichkeit $p$) $n$-mal unabhängig durchführen und $H_n$ die Anzahl der Erfolge ist, dann: $P\left(\left|\frac{H_n}{n} - p\right| < c\right) \to 1 \quad \text{für } n \to \infty$. Die relative Häufigkeit konvergiert gegen die theoretische Wahrscheinlichkeit.}

\definition{\textbf{Zentraler Grenzwertsatz (Central Limit Theorem, CLT):}

\textbf{Theorem:} Wenn $X_1, X_2, \ldots, X_n$ eine Stichprobe von u.i.v.\ Zufallsvariablen mit Erwartungswert $\mu$ und Varianz $\sigma^2 > 0$ sind, dann ist ihre Summe

\formel{S_n = X_1 + X_2 + \cdots + X_n}

im Limes für $n \to \infty$ normalverteilt:

\formel{S_n \xrightarrow{d} N(n\mu, n\sigma^2)}

oder äquivalent: Der Stichprobenmittelwert ist asymptotisch normalverteilt:

\formel{\bar{X}_n = \frac{S_n}{n} \xrightarrow{d} N\left(\mu, \frac{\sigma^2}{n}\right)}

Die \textbf{standardisierte Form}:

\formel{Z_n = \frac{S_n - n\mu}{\sqrt{n}\sigma} = \frac{\bar{X}_n - \mu}{\sigma/\sqrt{n}} \xrightarrow{d} N(0,1)}

\textbf{Das Besondere:} Die Summe (oder der Mittelwert) ist \textbf{näherungsweise normalverteilt, unabhängig von der Verteilung} der ursprünglichen Zufallsvariablen $X_i$! Dies gilt selbst wenn die $X_i$ uniform, exponentialverteilt oder beliebig anders verteilt sind.

\textbf{Praktische Relevanz:} Viele reale Phänomene sind Summen vieler unabhängiger Einflussfaktoren (z.\,B.\ Messfehler, Körpergröße, Prüfungsergebnisse) und sind daher näherungsweise normalverteilt.

\textbf{Beispiel (Würfel):} 
\begin{itemize}[leftmargin=1.5em]
    \item Ein Würfel hat $\mu = 3{,}5$ und $\sigma^2 = 2{,}9$.
    \item Die Summe von $n=100$ Würfeln: $S_{100} \sim N(350, 290)$ (näherungsweise).
    \item Der Mittelwert von 100 Würfeln: $\bar{X}_{100} \sim N(3{,}5, 0{,}029)$ (näherungsweise).
\end{itemize}

\textbf{Praktische Anwendung in der Statistik:}
\begin{itemize}[leftmargin=1.5em]
    \item Konstruktion von Konfidenzintervallen für den Mittelwert (auch bei nicht-normalverteilten Populationen)
    \item Begründung für die Verwendung von $z$-Tests auch bei großen Stichproben
    \item Erklärung, warum viele Tests und Schätzverfahren asymptotisch normalverteilt sind
\end{itemize}

\textbf{Wichtige Faustregel:} Der CLT wirkt schon bei moderaten Stichprobengrößen ($n \approx 30$), besonders wenn die ursprüngliche Verteilung nicht zu asymmetrisch ist. Bei kleinen $n$ oder sehr schiefen Verteilungen sollte man größer sein.}

\newpage

\subsection{Mehrdimensionale Zufallsvariablen}

\definition{\textbf{Mehrdimensionale (Bivariate) Zufallsvariablen:} Ein zufälliger Vektor oder eine gemeinsame Zufallsvariable mit $n$ Komponenten:
\formel{\mathbf{X} = \begin{pmatrix} X_1 \\ X_2 \\ \vdots \\ X_n \end{pmatrix} \quad \text{oder im bivariaten Fall:} \quad (X, Y) = \begin{pmatrix} X \\ Y \end{pmatrix}}
wobei jede Komponente selbst eine diskrete oder stetige Zufallsvariable ist. Der Wertemenge ist eine Teilmenge von $\mathbb{R}^n$.

\textbf{Beispiel:} Bei einer zufällig ausgewählten Person interessieren wir uns für $(X, Y)$ mit $X$ = Körpergröße, $Y$ = Körpergewicht.}

\definition{\textbf{Gemeinsame Verteilung (Joint Distribution):} 

\textbf{Im diskreten Fall:} Die gemeinsame Wahrscheinlichkeitsfunktion (joint PMF) ist:
\formel{f_{X,Y}(x, y) = P(X = x, Y = y)}
mit der Normierungsbedingung: $\sum_x \sum_y f_{X,Y}(x, y) = 1$.

\textbf{Im stetigen Fall:} Die gemeinsame Dichtefunktion (joint PDF) ist:
\formel{f_{X,Y}(x, y) \geq 0 \quad \text{mit} \quad \int_{-\infty}^{\infty} \int_{-\infty}^{\infty} f_{X,Y}(x, y) \, dx \, dy = 1}

Die gemeinsame Verteilungsfunktion (CDF) für beide Fälle:
\formel{F_{X,Y}(x, y) = P(X \leq x, Y \leq y)}

Die gemeinsame Verteilung enthält vollständige Information über das Verhalten beider Variablen und deren Abhängigkeitsstruktur.}

\definition{\textbf{Randverteilungen (Marginal Distributions):} Die Verteilung einer einzelnen Komponente, wenn wir die Information über die anderen ``heraussummieren'':

\textbf{Im diskreten Fall:}
\formel{f_X(x) = \sum_y f_{X,Y}(x, y) \quad \text{und} \quad f_Y(y) = \sum_x f_{X,Y}(x, y)}

\textbf{Im stetigen Fall:}
\formel{f_X(x) = \int_{-\infty}^{\infty} f_{X,Y}(x, y) \, dy \quad \text{und} \quad f_Y(y) = \int_{-\infty}^{\infty} f_{X,Y}(x, y) \, dx}

Die Randverteilungen sind wieder Verteilungen eindimensionaler Zufallsvariablen. Sie enthalten weniger Information als die gemeinsame Verteilung.

\textbf{Beispiel (Kontingenztafel):} Bei diskreten Variablen sind die Randverteilungen die Zeilen- und Spaltensummen der Kontingenztafel.}

\definition{\textbf{Bedingte Verteilungen (Conditional Distributions):} Die Verteilung einer Variablen unter der Bedingung, dass eine andere einen bestimmten Wert hat.

\textbf{Im diskreten Fall:}
\formel{P(X = x \mid Y = y) = \frac{P(X = x, Y = y)}{P(Y = y)} = \frac{f_{X,Y}(x, y)}{f_Y(y)}}
wenn $P(Y = y) > 0$.

\textbf{Im stetigen Fall:}
\formel{f_{X \mid Y}(x \mid y) = \frac{f_{X,Y}(x, y)}{f_Y(y)}}
wenn $f_Y(y) > 0$.

Die bedingte Dichtefunktion erfüllt: $\int_{-\infty}^{\infty} f_{X \mid Y}(x \mid y) \, dx = 1$ für jedes feste $y$.}

\definition{\textbf{Unabhängigkeit von Zufallsvariablen:} Zwei Zufallsvariablen $X$ und $Y$ sind unabhängig, wenn ihre gemeinsame Verteilung das Produkt der Randverteilungen ist:

\textbf{Im diskreten Fall:}
\formel{f_{X,Y}(x, y) = f_X(x) \cdot f_Y(y) \quad \text{für alle } x, y}

\textbf{Im stetigen Fall:}
\formel{f_{X,Y}(x, y) = f_X(x) \cdot f_Y(y) \quad \text{für alle } x, y}

\textbf{Äquivalente Bedingung (für bedingte Wahrscheinlichkeiten):} Die bedingte Verteilung gleicht der Randverteilung:
\formel{P(X = x \mid Y = y) = P(X = x) \quad \text{und} \quad P(Y = y \mid X = x) = P(Y = y)}

\textbf{Intuition:} Die Information, dass $Y$ einen bestimmten Wert hat, ändert unsere Wahrscheinlichkeitsaussagen über $X$ nicht.

\textbf{Wichtig:} Unkorreliertheit $\Rightarrow$ NICHT notwendig unabhängig, aber Unabhängigkeit $\Rightarrow$ Unkorreliertheit.}

\textbf{Zusammenfassung (Schema):}
\begin{table}[h]
    \centering
    \small
    \begin{tabular}{|l|l|l|}
        \hline
        \textbf{Konzept} & \textbf{Diskret} & \textbf{Stetig} \\
        \hline
        Gemeinsam & $f_{X,Y}(x, y) = P(X=x, Y=y)$ & $f_{X,Y}(x, y)$ Dichtefunktion \\
        \hline
        Rand ($X$) & $f_X(x) = \sum_y f_{X,Y}(x, y)$ & $f_X(x) = \int f_{X,Y}(x, y) dy$ \\
        \hline
        Bedingt & $P(X=x \mid Y=y) = \frac{f_{X,Y}(x,y)}{f_Y(y)}$ & $f_{X \mid Y}(x \mid y) = \frac{f_{X,Y}(x,y)}{f_Y(y)}$ \\
        \hline
        Unabhängig & $f_{X,Y}(x,y) = f_X(x) \cdot f_Y(y)$ & $f_{X,Y}(x,y) = f_X(x) \cdot f_Y(y)$ \\
        \hline
    \end{tabular}
\end{table}

\newpage

\subsection{Kovarianz und Korrelation}

\definition{\textbf{Kovarianz (Covariance):} Ein Maß für den linearen Zusammenhang zwischen zwei Zufallsvariablen $X$ und $Y$. Sie misst, wie zwei Variablen gemeinsam von ihren Mittelwerten abweichen.}

\textbf{Definition:}
\formel{\text{Cov}(X, Y) = \sigma_{XY} = \mathbb{E}[(X - \mu_X)(Y - \mu_Y)] = \mathbb{E}[XY] - \mathbb{E}[X] \mathbb{E}[Y]}

\textbf{Im Stichprobenfall (empirische Kovarianz):}
\formel{s_{XY} = \frac{1}{n-1} \sum_{i=1}^n (x_i - \bar{x})(y_i - \bar{y})}

\textbf{Interpretation:}
\begin{itemize}[leftmargin=1.5em]
    \item $\text{Cov}(X, Y) > 0$: Positive Assoziation (wenn $X$ größer als $\mu_X$, dann tendenziell $Y > \mu_Y$).
    \item $\text{Cov}(X, Y) < 0$: Negative Assoziation (wenn $X$ größer als $\mu_X$, dann tendenziell $Y < \mu_Y$).
    \item $\text{Cov}(X, Y) = 0$: Keine lineare Assoziation (unkorreliert).
\end{itemize}

\textbf{Probleme:} Die Kovarianz hängt von der Skalierung der Variablen ab — zwei äquivalente Variablen mit unterschiedlichen Einheiten haben unterschiedliche Kovarianzen. Dies erschwert Vergleiche.

\definition{\textbf{Korrelationskoeffizient (Pearson's Correlation Coefficient):} Das standardisierte Maß für den linearen Zusammenhang zwischen $X$ und $Y$:}

\formel{\rho_{X,Y} = \frac{\text{Cov}(X, Y)}{\sigma_X \cdot \sigma_Y} \in [-1, 1]}

\textbf{Im Stichprobenfall (empirischer Korrelationskoeffizient):}
\formel{r = \frac{\sum_{i=1}^n (x_i - \bar{x})(y_i - \bar{y})}{\sqrt{\sum_{i=1}^n (x_i - \bar{x})^2} \cdot \sqrt{\sum_{i=1}^n (y_i - \bar{y})^2}} = \frac{s_{XY}}{s_X \cdot s_Y}}

\textbf{Interpretation:}
\begin{itemize}[leftmargin=1.5em]
    \item $\rho = 1$: Perfekte positive lineare Beziehung ($Y = a + bX$ mit $b > 0$).
    \item $\rho = -1$: Perfekte negative lineare Beziehung ($Y = a + bX$ mit $b < 0$).
    \item $\rho = 0$: Keine lineare Beziehung (kann aber nichtlineare Beziehung existieren!).
    \item $|\rho|$ nähe 0: Schwache lineare Beziehung; $|\rho|$ nahe 1: Starke lineare Beziehung.
\end{itemize}

\textbf{Faustregel für Stärke:} $|\rho| > 0{,}7$ = Stark; $0{,}4 < |\rho| \leq 0{,}7$ = Moderat; $|\rho| \leq 0{,}4$ = Schwach (variiert je nach Kontext).

\definition{\textbf{Eigenschaften der Korrelation:}}
\begin{itemize}[leftmargin=1.5em]
    \item \textbf{Dimensionslos:} Im Gegensatz zur Kovarianz ist die Korrelation unabhängig von der Skalierung.
    \item \textbf{Symmetrisch:} $\rho_{X,Y} = \rho_{Y,X}$.
    \item \textbf{Invariant unter linearen Transformationen:} Wenn $X' = aX + b$ und $Y' = cY + d$ mit $a, c > 0$, dann $\text{Corr}(X', Y') = \text{Corr}(X, Y)$.
    \item \textbf{Nur bei gemeinsamer Normalverteilung äquivalent zu Unabhängigkeit:} $\rho = 0$ $\Leftrightarrow$ Unabhängigkeit nur für $(X,Y) \sim N(\mu_X, \mu_Y, \sigma_X^2, \sigma_Y^2, \rho)$. In anderen Fällen: Unkorreliertheit $\not\Rightarrow$ Unabhängigkeit.
\end{itemize}

\textbf{Korrelationsmatrix:} Für mehrdimensionale Vektoren $(X_1, \ldots, X_p)$ die $p \times p$-Matrix:
\formel{\mathbf{R} = \begin{pmatrix}
1 & \rho_{12} & \cdots & \rho_{1p} \\
\rho_{21} & 1 & \cdots & \rho_{2p} \\
\vdots & \vdots & \ddots & \vdots \\
\rho_{p1} & \rho_{p2} & \cdots & 1
\end{pmatrix}}
mit $\rho_{ij} = \rho_{ji}$ und $\rho_{ii} = 1$.

\textbf{Heatmap-Visualisierung:} Korrelationsmatrizen werden oft als Heatmaps gezeigt, um Muster schnell zu erkennen.

\definition{\textbf{Kovarianzmatrix und Varianz-Kovarianz-Matrix:} Für einen Vektor $\mathbf{X} = (X_1, \ldots, X_p)^T$ mit Mittelwert $\boldsymbol{\mu} = (\mu_1, \ldots, \mu_p)^T$:

\formel{\boldsymbol{\Sigma} = \text{Cov}(\mathbf{X}) = \mathbb{E}[(\mathbf{X} - \boldsymbol{\mu})(\mathbf{X} - \boldsymbol{\mu})^T] = \begin{pmatrix}
\sigma_1^2 & \sigma_{12} & \cdots & \sigma_{1p} \\
\sigma_{21} & \sigma_2^2 & \cdots & \sigma_{2p} \\
\vdots & \vdots & \ddots & \vdots \\
\sigma_{p1} & \sigma_{p2} & \cdots & \sigma_p^2
\end{pmatrix}}

Die Diagonalelemente sind Varianzen, die Off-Diagonalelemente sind Kovarianzen. Diese Matrix ist symmetrisch und positiv semi-definit.}

\textbf{Korrelation vs. Kausalität — das wichtigste Konzept:}

\textbf{Fundamentale Warnung:} ``Korrelation impliziert NICHT Kausalität!'' Eine hohe Korrelation $\rho \approx 1$ zwischen zwei Variablen bedeutet nicht, dass eine die andere verursacht. Es gibt mehrere Szenarien:

\begin{itemize}[leftmargin=1.5em]
    \item \textbf{Kausalität:} $X$ beeinflusst $Y$ (z.\,B.\ Rauchen $\to$ Lungenkrebs).
    \item \textbf{Umgekehrte Kausalität:} $Y$ beeinflusst $X$.
    \item \textbf{Gemeinsame Ursache:} Beide werden durch eine dritte Variable $Z$ beeinflusst (Confounding). Z.\,B.\ Eiscreme-Konsum und Ertranken — beide steigen mit der Temperatur.
    \item \textbf{Zufälliger Zusammenhang:} Spurious Correlation, besonders bei vielen Variablen (wenn $p$ groß ist, findet man zufällig hohe Korrelationen).
\end{itemize}

\textbf{Beispiel (Simpson's Paradoxon):} Gesamt-Korrelation kann unterschiedliches Vorzeichen haben als die Korrelation in Subgruppen!

\textbf{Fazit:} Für Kausalitätsaussagen braucht man experimentelle Designs oder Kausalitäts-Modellierung (Causal Inference), nicht nur Korrelationen.

\newpage

\subsection{Regressionen}

\definition{\textbf{Regressionsanalyse:} Ein statistisches Verfahren zur Modellierung des Zusammenhangs zwischen einer oder mehreren unabhängigen Variablen (\textbf{Prädiktoren, Features}) $x_1, \ldots, x_p$ und einer abhängigen Variable (\textbf{Response}) $y$.}

\textbf{Ziele:}
\begin{itemize}[leftmargin=1.5em]
    \item \textbf{Prognose (Prediction):} Vorhersage von $y$ für neue, unbekannte Prädiktoren.
    \item \textbf{Inferenz (Inference):} Verständnis und Quantifizierung des Einflusses einzelner Prädiktoren auf $y$.
    \item \textbf{Modellierung:} Formulierung des funktionalen Zusammenhangs $y = f(\mathbf{x}) + \varepsilon$.
\end{itemize}

\subsubsection{Einfache Lineare Regression (SLR)}

\definition{\textbf{Modell der Einfachen Linearen Regression:} Ein Regressionsmodell mit genau \textbf{einem} Prädiktor $x$ und einer stetigen Response $y$:

\formel{y_i = \beta_0 + \beta_1 x_i + \varepsilon_i, \quad i = 1, \ldots, n}

wobei:
\begin{itemize}[leftmargin=1.5em]
    \item $\beta_0$ = Intercept (Achsenabschnitt), der y-Wert bei $x=0$.
    \item $\beta_1$ = Slope (Steigung), die Änderung von $y$ pro Einheit Änderung in $x$.
    \item $\varepsilon_i$ = Fehlerterm (zufällig, unbeobachtbar).
    \item Die Beziehung $E[y \mid x] = \beta_0 + \beta_1 x$ ist linear in $x$.
\end{itemize}}

\definition{\textbf{Annahmen der SLR:}
\begin{itemize}[leftmargin=1.5em]
    \item \textbf{Linearität:} Der Mittelwert von $y$ ist eine lineare Funktion von $x$.
    \item \textbf{Unabhängigkeit:} Fehler $\varepsilon_i$ und $\varepsilon_j$ sind unabhängig für $i \neq j$.
    \item \textbf{Homoskedastizität:} Die Varianz $\text{Var}(\varepsilon) = \sigma^2$ ist konstant (hängt nicht von $x$ ab).
    \item \textbf{Normalverteilung der Fehler:} $\varepsilon_i \sim N(0, \sigma^2)$ (nötig für Inferenz, nicht für Punktschätzung).
\end{itemize}}

\textbf{Schätzung via Kleinste-Quadrate (Ordinary Least Squares, OLS):}

\textbf{Herleitung (Kurz):} Minimiere die Summe der quadrierten Residuen:
\formel{S(\beta_0, \beta_1) = \sum_{i=1}^n (y_i - \beta_0 - \beta_1 x_i)^2}

Partielle Ableitungen nach $\beta_0, \beta_1$ setzen gleich Null:
\formel{\frac{\partial S}{\partial \beta_0} = -2 \sum_{i=1}^n (y_i - \hat{\beta}_0 - \hat{\beta}_1 x_i) = 0}
\formel{\frac{\partial S}{\partial \beta_1} = -2 \sum_{i=1}^n x_i(y_i - \hat{\beta}_0 - \hat{\beta}_1 x_i) = 0}

Dies führt zu den \textbf{Normalgleichungen} mit den Lösungen:
\formel{\hat{\beta}_1 = \frac{\sum_{i=1}^n (x_i - \bar{x})(y_i - \bar{y})}{\sum_{i=1}^n (x_i - \bar{x})^2} = \frac{s_{xy}}{s_x^2} = r_{xy} \frac{s_y}{s_x}}

\formel{\hat{\beta}_0 = \bar{y} - \hat{\beta}_1 \bar{x}}

Die geschätzte Regressionsgerade ist: $\hat{y} = \hat{\beta}_0 + \hat{\beta}_1 x$.

\definition{\textbf{Residuen und Varianzzerlegung:} 
\begin{itemize}[leftmargin=1.5em]
    \item \textbf{Residuum:} $e_i = y_i - \hat{y}_i = y_i - \hat{\beta}_0 - \hat{\beta}_1 x_i$ (beobachtete minus vorhergesagte Werte).
    \item \textbf{Residualsumme der Quadrate (RSS):} $\text{RSS} = \sum_i e_i^2$, minimiert durch OLS.
    \item \textbf{Gesamtsumme der Quadrate (TSS):} $\text{TSS} = \sum_i (y_i - \bar{y})^2$.
    \item \textbf{Erklärte Summe der Quadrate (ESS):} $\text{ESS} = \sum_i (\hat{y}_i - \bar{y})^2 = \text{TSS} - \text{RSS}$.
\end{itemize}

Varianzzerlegung: $\text{TSS} = \text{ESS} + \text{RSS}$.}

\definition{\textbf{Bestimmtheitsmaß ($R^2$, Coefficient of Determination):}
\formel{R^2 = \frac{\text{ESS}}{\text{TSS}} = 1 - \frac{\text{RSS}}{\text{TSS}} \in [0, 1]}

$R^2$ ist der Anteil der Variabilität in $y$, der durch das Modell erklärt wird.
\begin{itemize}[leftmargin=1.5em]
    \item $R^2 = 1$: Perfekte Anpassung (alle Residuen sind Null).
    \item $R^2 = 0$: Das Modell erklärt keine Variabilität (die Regressionslinie ist horizontal).
    \item Für die SLR: $R^2 = r_{xy}^2$, d.h.\ das Bestimmtheitsmaß ist das Quadrat des Korrelationskoeffizienten.
\end{itemize}}

\subsubsection{Multiple Lineare Regression (MLR)}

\definition{\textbf{Modell der Multiplen Linearen Regression:} Erweiterung auf $p > 1$ Prädiktoren:

\formel{y_i = \beta_0 + \beta_1 x_{i1} + \beta_2 x_{i2} + \cdots + \beta_p x_{ip} + \varepsilon_i}

\textbf{Matrixnotation:}
\formel{\mathbf{y} = \mathbf{X} \boldsymbol{\beta} + \boldsymbol{\varepsilon}}

mit:
\begin{itemize}[leftmargin=1.5em]
    \item $\mathbf{y} = (y_1, \ldots, y_n)^T$ (Response-Vektor).
    \item $\mathbf{X} = \begin{pmatrix} 1 & x_{11} & \cdots & x_{1p} \\ \vdots & \vdots & \ddots & \vdots \\ 1 & x_{n1} & \cdots & x_{np} \end{pmatrix}$ ($n \times (p+1)$ Design-Matrix mit Spalte der 1er für Intercept).
    \item $\boldsymbol{\beta} = (\beta_0, \beta_1, \ldots, \beta_p)^T$ (Koeffizientenvektor).
    \item $\boldsymbol{\varepsilon} = (\varepsilon_1, \ldots, \varepsilon_n)^T$ (Fehlervektor).
\end{itemize}

\textbf{OLS-Schätzer:} Der OLS-Schätzer minimiert $\|\mathbf{y} - \mathbf{X}\boldsymbol{\beta}\|^2$. Die Lösung (falls $\mathbf{X}^T\mathbf{X}$ invertierbar):
\formel{\hat{\boldsymbol{\beta}} = (\mathbf{X}^T \mathbf{X})^{-1} \mathbf{X}^T \mathbf{y}}

\textbf{Vorhersage:} $\hat{\mathbf{y}} = \mathbf{X} \hat{\boldsymbol{\beta}}$, mit Residuen $\mathbf{e} = \mathbf{y} - \hat{\mathbf{y}}$.}

\definition{\textbf{Anpassungsgüte in MLR:}
\begin{itemize}[leftmargin=1.5em]
    \item \textbf{Multiple $R^2$:} Wie bei SLR, aber mit mehreren Prädiktoren. Steigt immer, wenn Prädiktoren hinzugefügt werden (auch unwichtige!).
    \item \textbf{Adjusted $R^2$ (R$^2_{\text{adj}}$):} Bestraft zusätzliche Prädiktoren:
    \formel{R^2_{\text{adj}} = 1 - \frac{(1 - R^2)(n-1)}{n - p - 1}}
    ist strenger und wird für Modellvergleich bevorzugt.
    \item \textbf{Residualstandardfehler (RSE):} $\text{RSE} = \sqrt{\frac{\text{RSS}}{n - p - 1}}$, Schätzung von $\sigma$.
\end{itemize}}

\definition{\textbf{Statistische Tests in MLR:}
\begin{itemize}[leftmargin=1.5em]
    \item \textbf{t-Test für einzelne Koeffizienten:} Test ob $\beta_j = 0$ (kein Effekt von $x_j$, gegeben andere Prädiktoren).
    \item \textbf{F-Test für Gesamtmodell:} Test ob mindestens ein $\beta_j \neq 0$. Alternativ: Vergleich verschachtelter Modelle via F-Test.
\end{itemize}}

\subsubsection{Vergleich: Einfach (SLR) vs. Multipel (MLR) vs. Multivariat}

\textbf{Multivariate Regression:} In der Statistik bedeutet ``multivariat'' typischerweise \textbf{mehrere Response-Variablen} $y_1, \ldots, y_k$ (nicht mehrere Prädiktoren!). Obwohl Regression meist mit \textit{einer} Response erfolgt, kann man auch $k$ separate (oder gemeinsame) Regressionen fitten:
\formel{y_{ij} = \beta_{0j} + \beta_{1j} x_i + \cdots + \beta_{pj} x_{ip} + \varepsilon_{ij} \quad \text{für } j = 1, \ldots, k}

In Matrixform: $\mathbf{Y} = \mathbf{X} \mathbf{B} + \mathbf{E}$, wobei $\mathbf{Y}$ ist $(n \times k)$, $\mathbf{B}$ ist $(p+1) \times k$.

\begin{table}[h]
    \centering
    \small
    \begin{tabular}{|l|c|c|c|}
        \hline
        \textbf{Modell} & \textbf{Prädiktoren} & \textbf{Response} & \textbf{Koeffizienten} \\
        \hline
        Einfach (SLR) & 1 $(x)$ & 1 $(y)$ & $\beta_0, \beta_1$ \\
        \hline
        Multipel (MLR) & $p > 1$ $(x_1, \ldots, x_p)$ & 1 $(y)$ & $\beta_0, \beta_1, \ldots, \beta_p$ \\
        \hline
        Multivariat & $p \geq 1$ & $k > 1$ $(y_1, \ldots, y_k)$ & Koeff.-Matrix $\mathbf{B}$ \\
        \hline
    \end{tabular}
\end{table}

\subsubsection{Logistische Regression}

\definition{\textbf{Logistische Regression:} Regression für \textbf{binäre Response} $y \in \{0, 1\}$ (z.\,B.\ Klassifikation: Ja/Nein, Krankheit/Gesund).

\textbf{Motivation:} Lineare Regression auf binäre Ziele ist problematisch (vorhergesagte Werte können außerhalb $[0, 1]$ liegen, die Fehlervarianz ist nicht konstant). Logistische Regression modelliert stattdessen die \textbf{Wahrscheinlichkeit}:
\formel{p_i = P(y_i = 1 \mid \mathbf{x}_i)}

Das Modell:
\formel{p_i = \frac{e^{\beta_0 + \beta_1 x_{i1} + \cdots + \beta_p x_{ip}}}{1 + e^{\beta_0 + \beta_1 x_{i1} + \cdots + \beta_p x_{ip}}} = \frac{1}{1 + e^{-(\beta_0 + \boldsymbol{\beta}^T \mathbf{x}_i)}}}

Das ist die \textbf{logistische Funktion} (Sigmoid), die Werte in $[0, 1]$ erzwingt.

\textbf{Log-Odds (Logits):} Die logit-Transformation ist die Umkehrung:
\formel{\text{logit}(p_i) = \log\left(\frac{p_i}{1 - p_i}\right) = \beta_0 + \beta_1 x_{i1} + \cdots + \beta_p x_{ip}}

Dies ist linear in den Prädiktoren — daher der Name ``logistische \textit{Regression}''.}

\textbf{Interpretation der Koeffizienten:} Ein Anstieg um 1 Einheit in $x_j$ multipliziert die \textbf{Odds} mit Faktor $e^{\beta_j}$.

\textbf{Schätzung:} OLS funktioniert nicht; stattdessen \textbf{Maximum Likelihood Estimation (MLE)}. Die Log-Likelihood wird iterativ maximiert (z.\,B.\ via Newton-Raphson).

\textbf{Modellvergütung in Logistischer Regression:}
\begin{itemize}[leftmargin=1.5em]
    \item \textbf{Deviance:} Ähnlich wie RSS in linearer Regression.
    \item \textbf{McFadden $R^2$} oder \textbf{Pseudo $R^2$}: Nicht direkt vergleichbar mit $R^2$ der linearen Regression.
    \item \textbf{Confusion Matrix:} True Positives, False Positives, etc. zur Berechnung von Accuracy, Sensitivity, Specificity.
    \item \textbf{ROC-Kurve und AUC:} Visualisiert Klassifikationsleistung über verschiedene Klassifikationsschwellen.
\end{itemize}

\newpage




\subsection{Polynomregression, Splines und additive Modelle}

\definition{\textbf{Polynomregression:} Erweiterung der linearen Regression auf \textbf{polynomiale} Funktionen des Prädiktors $x$:

\formel{y_i = \beta_0 + \beta_1 x_i + \beta_2 x_i^2 + \cdots + \beta_d x_i^d + \varepsilon_i}

wobei $d$ der Polynomgrad ist (z.\,B.\ $d = 2$ für quadratisch, $d = 3$ für kubisch).

\textbf{Reformulierung als lineare Regression:} Definiere neue Prädiktoren $x_{i1} = x_i$, $x_{i2} = x_i^2$, \ldots, $x_{id} = x_i^d$. Dann ist das Modell wieder \textbf{linear in den Koeffizienten}:
\formel{y_i = \beta_0 + \beta_1 x_{i1} + \beta_2 x_{i2} + \cdots + \beta_d x_{id} + \varepsilon_i}

Daher kann Polynomregression \textbf{als spezielle Multipel-Lineare Regression} behandelt werden und mit OLS geschätzt werden.}

\textbf{Modellkomplexität und Overfitting:}
\begin{itemize}[leftmargin=1.5em]
    \item Höhere Polynomgrade ($d$ groß) können zu \textbf{Overfitting} führen — das Modell passt sich zu stark an die Trainingsdaten an, generalisiert aber schlecht auf neue Daten.
    \item \textbf{Bias-Variance Tradeoff}: Niedriges $d$ = hoher Bias (Underfitting), hohes $d$ = hohe Varianz (Overfitting).
    \item Wahl von $d$ via \textbf{Cross-Validation}, AIC, BIC oder Visualisierung der Residuen.
\end{itemize}

\definition{\textbf{Splines (Glättungssplines):} Eine Verallgemeinerung von Polynomregression, die \textbf{lokal} polynomiale Funktionen kombiniert.

\textbf{Idee:} Teile den Definitionsbereich von $x$ in $K$ Intervalle auf, in jedem Intervall sei die Funktion polynomial (üblicherweise kubisch), mit Stetigkeit und Glattheits-Bedingungen an den Grenzpunkten (\textbf{Knots}).

\textbf{Cubic Spline:} Ein weit verbreitetes Modell:
\begin{itemize}[leftmargin=1.5em]
    \item In jedem Intervall: Kubisches Polynom (Grad 3).
    \item An den Knots: Funktion, erste und zweite Ableitung sind stetig.
\end{itemize}

\textbf{Anzahl der Freiheitsgrade:} Hängt von der Anzahl der Knots $K$ ab. Mit $K$ Knots hat ein kubischer Spline etwa $K + 3$ Freiheitsgrade (näherungsweise).}

\textbf{Mathematische Formulierung via Basis-Funktionen:} Splines können als Linearkombination von Basis-Funktionen dargestellt werden:
\formel{f(x) = \sum_{j=1}^M \beta_j B_j(x)}

wobei $B_j(x)$ sind Basis-Funktionen (z.\,B.\ B-Spline-Basis) und $M$ ist die Anzahl der Basis-Funktionen. Das Modell ist wieder linear in den Koeffizienten und kann mit OLS geschätzt werden!

\textbf{Vergleich: Polynomregression vs.\ Splines:}
\begin{table}[h]
    \centering
    \small
    \begin{tabular}{|l|l|l|}
        \hline
        \textbf{Aspekt} & \textbf{Polynomregression} & \textbf{Splines} \\
        \hline
        Flexibilität & Global (hängt von $d$ ab) & Lokal (variabel in Intervallen) \\
        \hline
        Overfitting-Risiko & Höher bei großem $d$ & Besser kontrollierbar \\
        \hline
        Glattheit & Automatisch Grad $d-1$ & Kontrollierbar (typisch C2) \\
        \hline
        Extrapolation & Polynomial explodiert & Oft linear außerhalb \\
        \hline
        Interpretation & Einfacher & Komplexer \\
        \hline
    \end{tabular}
\end{table}

\definition{\textbf{Smoothing Splines:} Eine regularisierte Variante von Splines, die einen Kompromiss zwischen Anpassung und Glattheit findet.

\textbf{Ziel-Funktion:} Minimiere:
\formel{\sum_{i=1}^n (y_i - f(x_i))^2 + \lambda \int (f''(x))^2 dx}

wobei:
\begin{itemize}[leftmargin=1.5em]
    \item Erster Term: Anpassung an Daten (wie OLS).
    \item Zweiter Term: Glattheits-Strafterm (je rauer $f$, desto höher der Strafterm).
    \item $\lambda \geq 0$: Regularisierungsparameter (Glättungskonstante).
\end{itemize}

Bei $\lambda = 0$: Perfekte Interpolation (sehr wiggly). Bei $\lambda \to \infty$: Lineare Regressionsgerade. Die optimale $\lambda$ wird oft via \textbf{Cross-Validation} gewählt.}

\subsubsection{Additive Modelle (GAM)}

\definition{\textbf{Generalized Additive Models (GAM):} Ein flexibles Regressions-Framework, das additive Struktur nutzt:

\formel{y_i = \beta_0 + f_1(x_{i1}) + f_2(x_{i2}) + \cdots + f_p(x_{ip}) + \varepsilon_i}

wobei $f_j$ sind \textbf{glatte Funktionen} (nicht unbedingt linear). Der große Vorteil: Man schätzt jede Funktion $f_j$ \textbf{separat}, eine Variable nach der anderen.}

\textbf{Flexibilität vs.\ Interpretierbarkeit:}
\begin{itemize}[leftmargin=1.5em]
    \item \textbf{Lineares Modell:} $f_j(x_j) = \beta_j x_j$ (starr, aber interpretierbar).
    \item \textbf{Additive Modelle:} $f_j$ sind glatte Funktionen (flexibel, aber weniger interpretierbar).
    \item \textbf{Allgemeines nichtlineares Modell:} $y = f(\mathbf{x}) + \varepsilon$ (höchste Flexibilität, schwer zu verstehen).
\end{itemize}

\definition{\textbf{Schätzung in GAM via Backfitting:} Ein iterativer Algorithmus:

\begin{enumerate}[leftmargin=2em]
    \item Initialisiere $\hat{\beta}_0 = \bar{y}$.
    \item Für jede Variable $j = 1, \ldots, p$:
    \begin{itemize}
        \item Berechne das Residuum, ``abzüglich'' aller anderen Funktionen: $\mathbf{r}_j = \mathbf{y} - \hat{\beta}_0 - \sum_{k \neq j} \hat{f}_k(\mathbf{x}_k)$.
        \item Schätze $\hat{f}_j$ durch Glätten von $\mathbf{r}_j$ gegen $\mathbf{x}_j$ (z.\,B.\ via Smoothing Spline).
    \end{itemize}
    \item Wiederhole, bis Konvergenz.
\end{enumerate}

Dies ist ein \textbf{greedy} Algorithmus, der schnell ist, aber nicht garantiert ein globales Optimum findet.}

\textbf{Basis-Funktionen in GAM:} Jede $f_j$ wird als Linearkombination von Basis-Funktionen dargestellt:
\formel{f_j(x_j) = \sum_{k=1}^{K_j} \beta_{jk} B_{jk}(x_j)}

Beliebte Wahlen: B-Splines, natural cubic splines, smooth functions. Mit dieser Darstellung wird das GAM-Optimization wieder zu einem \textbf{regularisierten Linearen Problem}.

\definition{\textbf{Variablen-Auswahl in GAM:} 

\begin{itemize}[leftmargin=1.5em]
    \item \textbf{Partial Dependence Plot:} Plot von $\hat{f}_j(x_j)$ gegen $x_j$, zeigt die isolierte Effekt von $x_j$.
    \item \textbf{Effectivity:} ANOVA-ähnliche Tests um zu prüfen ob $f_j$ signifikant ist (gegen $f_j = \text{const}$).
    \item \textbf{Regularisierung:} Penalty auf Rauheit zwingt unwichtige Variablen zu $f_j(x_j) = \text{const}$ (automatische Selection).
\end{itemize}}

\textbf{GAM vs.\ Lineare Regression:}
\begin{table}[h]
    \centering
    \small
    \begin{tabular}{|l|l|l|}
        \hline
        \textbf{Aspekt} & \textbf{Lineare Regression} & \textbf{GAM} \\
        \hline
        Modell & $y = \beta_0 + \sum \beta_j x_j + \varepsilon$ & $y = \beta_0 + \sum f_j(x_j) + \varepsilon$ \\
        \hline
        Flexibilität & Niedrig (linear) & Höher (nichtlinear) \\
        \hline
        Interpretierbarkeit & Hoch (Koeff. klar) & Moderat (Plots nötig) \\
        \hline
        Rechenkomplexität & Niedrig (geschlossen) & Höher (iterativ) \\
        \hline
        Overfitting & Geringer & Höher (erfordert Regularisierung) \\
        \hline
    \end{tabular}
\end{table}

\definition{\textbf{Multivariate GAM:} Erweiterung auf \textbf{mehrere Response-Variablen}:

\formel{y_{ij} = \beta_{0j} + f_{1j}(x_{i1}) + \cdots + f_{pj}(x_{ip}) + \varepsilon_{ij} \quad (j = 1, \ldots, k)}

oder mit geteilten Basis-Funktionen für \textbf{simultane} Schätzung aller Responses. Dies ist besonders wertvoll wenn die Responses korreliert sind.}

\textbf{Praktische Anwendungen:}
\begin{itemize}[leftmargin=1.5em]
    \item \textbf{Polynomregression:} Wenn die Beziehung quadratisch oder kubisch vermutet wird; einfach zu interpretieren.
    \item \textbf{Splines:} Für glatte, nichtlineare Beziehungen mit lokal variierender Komplexität.
    \item \textbf{GAM:} Für viele Prädiktoren mit komplexen Abhängigkeiten, wenn Interpretierbarkeit wichtig ist (Compromise zwischen linearer Regression und Black-Box-Modellen).
\end{itemize}

\newpage

\subsection{PCA, PCR und MDS}

\subsubsection{Principal Component Analysis (PCA)}

\definition{\textbf{Principal Component Analysis (PCA):} Ein unsupervisierter Algorithmus zur \textbf{Dimensionsreduktion} und Datenexploration, der eine neue Koordinatensystem findet, in welchem die höchste Varianz (= Information) in den ersten wenigen Achsen konzentriert ist.

\textbf{Kernidee:} 
\begin{itemize}[leftmargin=1.5em]
    \item Gegeben: Datenmatrix $\mathbf{X}$ (Größe $n \times p$) mit $n$ Beobachtungen und $p$ (möglicherweise hochdimensionalen) Variablen.
    \item Gesucht: Neue orthogonale Koordinaten (Principal Components, PCs), die eine maximale Varianz erklären.
    \item Die neuen Achsen sind \textbf{orthogonal} (unkorreliert) und linear kombiniert die ursprünglichen Variablen.
\end{itemize}

\textbf{Motivation:}
\begin{itemize}[leftmargin=1.5em]
    \item \textbf{Visualisierung:} Bei $p = 20$ Variablen gibt es $\binom{20}{2} = 190$ paarweise Scatterplots. Mit PCA können wir auf 2 oder 3 Dimensionen reduzieren und ein einziges Scatterplot verwenden.
    \item \textbf{Datenreduktion:} Oft ist nur ein Bruchteil der Variabilität in den ersten wenigen PCs konzentriert; der Rest ist Rauschen. Diese können ignoriert werden.
    \item \textbf{Multikollinearität in MLR:} Wenn Prädiktoren stark korreliert sind, kann PCR (Regression auf PCs) verwendet werden, um das Problem zu beheben.
    \item \textbf{Explorative Datenanalyse:} Identifikation von Clustern, Ausreißern, oder Struktur in Hochdimensionalen Daten.
\end{itemize}}

\definition{\textbf{Mathematische Formulierung:}

Die PCA-Zerlegung ist:
\formel{\mathbf{X} = \mathbf{T} \mathbf{P}^T + \mathbf{E}}

wobei:
\begin{itemize}[leftmargin=1.5em]
    \item $\mathbf{T}$ ($n \times a$) = \textbf{Score-Matrix}: Die Koordinaten der Beobachtungen in der neuen Koordinatensystem.
    \item $\mathbf{P}$ ($p \times a$) = \textbf{Loading-Matrix}: Die Spalten sind die Eigenvektoren der Kovarianzmatrix (die neuen Achsen-Richtungen).
    \item $\mathbf{E}$ ($n \times p$) = Fehlermatrix (wird meist ignoriert).
    \item $a = \min(n, p)$ wenn alle Komponenten extrahiert werden.
\end{itemize}

Die Spalten von $\mathbf{P}$ sind \textbf{orthonormal}: $\mathbf{P}^T \mathbf{P} = \mathbf{I}_a$. Die Scores in $\mathbf{T}$ sind \textbf{unkorreliert}: $\mathbf{T}^T \mathbf{T} = \text{diag}(\lambda_1, \ldots, \lambda_a)$ (diagonal).}

\definition{\textbf{Algorithmus (Herangehensweise):}

\begin{enumerate}[leftmargin=2em]
    \item \textbf{Standardisierung:} Zentriere und skaliere $\mathbf{X}$ so, dass jede Spalte Mittelwert 0 und Standardabweichung 1 hat (wichtig für fair comparision, wenn Variablen unterschiedliche Einheiten haben).
    \item \textbf{Kovarianzmatrix:} Berechne die Kovarianzmatrix $\mathbf{C} = \frac{1}{n-1} \mathbf{X}^T \mathbf{X}$ (oder Korrelationsmatrix).
    \item \textbf{Eigenwertzerlegung:} Berechne Eigenwerte $\lambda_1 \geq \lambda_2 \geq \cdots \geq \lambda_p \geq 0$ und entsprechende Eigenvektoren $\mathbf{p}_1, \mathbf{p}_2, \ldots, \mathbf{p}_p$ von $\mathbf{C}$.
    \item \textbf{Komponenten extrahieren:} Wähle die ersten $a$ Eigenvektoren als Spalten von $\mathbf{P}$. Diese entsprechen den $a$ größten Eigenwerten.
    \item \textbf{Scores berechnen:} $\mathbf{T} = \mathbf{X} \mathbf{P}$.
\end{enumerate}}

\definition{\textbf{Varianzzerlegung:}

Der Eigenwert $\lambda_j$ quantifiziert die Varianz, die durch die $j$-te Komponente erklärt wird:
\formel{\text{Varianz von PC}_j = \lambda_j}

Der Anteil der Varianz erklärt durch die erste $a$ Komponenten:
\formel{\text{Cumulative Explained Variance} = \frac{\sum_{j=1}^a \lambda_j}{\sum_{j=1}^p \lambda_j} \in [0, 1]}

Ein \textbf{Scree-Plot} zeigt die Eigenwerte gegen die PC-Nummer. Man sucht oft einen ``Ellbogen'' (elbow), um die optimale Anzahl $a$ zu wählen.}

\definition{\textbf{Score-Plots und Loading-Plots:}

\begin{itemize}[leftmargin=1.5em]
    \item \textbf{Score-Plot:} Scatterplot der Scores $\mathbf{T}$, üblicherweise PC1 vs. PC2. Visualisiert die Beobachtungen im neuen Raum. Können Cluster, Ausreißer, oder Gruppierungen zeigen.
    \item \textbf{Loading-Plot:} Scatterplot der Loadings $\mathbf{P}$. Zeigt welche \textbf{Originalvariablen} stark auf jede PC laden (beeinflussen). Variablen nah bei Ursprung haben niedrige Loadings; Variable weit weg have hohe Loadings auf mindestens eine PC.
    \item \textbf{Biplot:} Kombiniert Score- und Loading-Plots in einem Diagramm. Erlaubt gleichzeitige Interpretation von Beobachtungen und Variablen.
\end{itemize}}

\newpage

\subsubsection{Principal Component Regression (PCR)}

\definition{\textbf{Problem, das PCR löst:} In der multiplen linearen Regression $\mathbf{y} = \mathbf{X} \boldsymbol{\beta} + \boldsymbol{\varepsilon}$ kann es problematisch sein wenn:
\begin{itemize}[leftmargin=1.5em]
    \item Die Anzahl der Prädiktoren groß ist (high $p$), oder sogar $p > n$ (mehr Prädiktoren als Beobachtungen).
    \item Die Prädiktoren \textbf{hochgradig korreliert} sind (\textbf{Multikollinearität}).
\end{itemize}

In diesen Fällen ist die OLS-Lösung $\hat{\boldsymbol{\beta}} = (\mathbf{X}^T \mathbf{X})^{-1} \mathbf{X}^T \mathbf{y}$ \textbf{instabil} oder \textbf{nicht definiert}.}

\definition{\textbf{PCR-Ansatz:} Statt mit den ursprünglichen Prädiktoren $\mathbf{X}$ zu arbeiten, verwende die \textbf{PCA-Scores} $\mathbf{T}$ als neue Prädiktoren:

\formel{\mathbf{y} = \mathbf{T} \boldsymbol{\vartheta} + \boldsymbol{\varepsilon}}

wobei $\boldsymbol{\vartheta} = (\vartheta_1, \ldots, \vartheta_a)^T$ sind die Regressionskoeffizienten bezüglich der Scores.}

\textbf{Vorteile:}
\begin{itemize}[leftmargin=1.5em]
    \item Die Scores $\mathbf{T}$ sind \textbf{orthogonal} (unkorreliert), daher keine Multikollinearität.
    \item Die OLS-Schätzung wird stabil: $\hat{\boldsymbol{\vartheta}} = (\mathbf{T}^T \mathbf{T})^{-1} \mathbf{T}^T \mathbf{y}$ ist wohldefiniert (eine Diagonalmatrix).
    \item Durch Verwendung von nur $a < p$ Komponenten wird Overfitting reduziert.
\end{itemize}

\definition{\textbf{PCR-Algorithmus:}

\begin{enumerate}[leftmargin=2em]
    \item Führe PCA auf $\mathbf{X}$ durch, extrahiere die ersten $a$ Principal Components.
    \item Berechne die Score-Matrix $\mathbf{T}$ der Größe $n \times a$.
    \item Führe ordinäre lineare Regression durch: $\mathbf{y} = \mathbf{T} \boldsymbol{\vartheta} + \boldsymbol{\varepsilon}$.
    \item Berechne $\hat{\boldsymbol{\vartheta}} = (\mathbf{T}^T \mathbf{T})^{-1} \mathbf{T}^T \mathbf{y}$.
    \item Transformiere zurück zu den ursprünglichen Koeffizienten: $\hat{\boldsymbol{\beta}}_{\text{PCR}} = \mathbf{P} \hat{\boldsymbol{\vartheta}}$.
    \item Die Wahl von $a$ wird üblicherweise via \textbf{Cross-Validation} bestimmt.
\end{enumerate}}

\textbf{Vergleich: MLR vs. PCR vs. Ridge/Lasso:}
\begin{table}[h]
    \centering
    \small
    \begin{tabular}{|l|l|l|l|}
        \hline
        \textbf{Aspekt} & \textbf{MLR} & \textbf{PCR} & \textbf{Ridge/Lasso} \\
        \hline
        Multikollinearität & Problem & Gelöst & Eingedämmt \\
        \hline
        $p > n$ & Nicht möglich & Möglich & Möglich \\
        \hline
        Interpretierbarkeit & Hoch (Originale) & Mittel (PCs) & Mittel (Shrinkage) \\
        \hline
        Rechenaufwand & Niedrig & Mittel (PCA + Reg) & Mittel \\
        \hline
    \end{tabular}
\end{table}

\subsubsection{Multidimensionale Skalierung (MDS)}

\definition{\textbf{Multidimensional Scaling (MDS):} Ein Unsupervisiertes Verfahren zur Dimensionsreduktion, das von einer \textbf{Distanz-/Unähnlichkeitsmatrix} ausgeht und versucht, eine Konfiguration von Objekten in niedrigdimensionalen Raum (typisch 2D oder 3D) zu finden, die diese Abstände möglichst gut erhält.

\textbf{Startpunkt:} Keine Rohdaten $\mathbf{X}$, sondern nur eine $n \times n$ Distanzmatrix $\mathbf{D}$ mit Elementen $d_{ij}$ = Distanz zwischen Objekten $i$ und $j$.

\textbf{Ziel:} Finde niedrigdimensionale Koordinaten $\mathbf{Z}$ (Größe $n \times a$ mit $a \ll p$) so, dass die Euklidischen Abstände in $\mathbf{Z}$ den ursprünglichen Abständen entsprechen.}

\textbf{Motivation:}
\begin{itemize}[leftmargin=1.5em]
    \item \textbf{Oft sind nur Distanzen bekannt:} Z.\,B.\ wenn Experten Paare von Objekten bewerten (ähnlich oder nicht), oder bei Umfrageabständen.
    \item \textbf{Reverse-Engineering:} Kann man aus Abständen die zugrundeliegende Struktur rekonstruieren?
    \item \textbf{Visualisierung:} Besonders wertvoll wenn die Rohdaten sehr hochdimensional sind oder komplexe Abstandsmetriken haben.
\end{itemize}

\definition{\textbf{Metric vs. Non-Metric MDS:}

\begin{itemize}[leftmargin=1.5em]
    \item \textbf{Metric (Classical) MDS:} Versucht, die \textbf{absoluten Distanzen} zu reproduzieren.
    \begin{itemize}
        \item Minimiere Stress-Funktion: $\text{Stress} = \sqrt{\frac{\sum_{i,j} (d_{ij} - \hat{d}_{ij})^2}{\sum_{i,j} d_{ij}^2}}$.
        \item $d_{ij}$ = ursprüngliche Distanz, $\hat{d}_{ij}$ = Distanz in niedrigdimensionaler Rekonstruktion.
    \end{itemize}
    \item \textbf{Non-Metric MDS:} Versucht, nur die \textbf{Ordnung} der Distanzen zu erhalten (kleinste Distanz bleibt klein, größte bleibt groß, aber exakte Werte sind unwichtig).
    \begin{itemize}
        \item Robuster gegen Messfehler oder nicht-euklidische Abstände.
        \item Verwendet Rank-korrelation statt absoluter Fehler.
    \end{itemize}
\end{itemize}}

\definition{\textbf{Classical MDS-Algorithmus:}

\begin{enumerate}[leftmargin=2em]
    \item \textbf{Double Centering:} Zentriere die Distanzmatrix $\mathbf{D}$ so, dass eine Gram-Matrix $\mathbf{G}$ entsteht (Innenprodukten der zentrierten Punkte).
    \item \textbf{Eigenwertzerlegung:} Berechne Eigenwerte $\lambda_1 \geq \lambda_2 \geq \cdots$ und Eigenvektoren von $\mathbf{G}$.
    \item \textbf{Koordinaten:} Die ersten $a$ Koordinaten eines Objekts sind $z_i = \sqrt{\lambda_k} \mathbf{v}_k$ (wobei $\mathbf{v}_k$ die Eigenvektoren).
    \item \textbf{Rekonstruktion:} Überprüfe ob die rekonstruierten Distanzen den ursprünglichen nahe kommen.
\end{enumerate}}

\textbf{Unterschied zu PCA:}
\begin{table}[h]
    \centering
    \small
    \begin{tabular}{|l|l|l|}
        \hline
        \textbf{Aspekt} & \textbf{PCA} & \textbf{MDS} \\
        \hline
        Input & Datenmatrix $\mathbf{X}$ & Distanzmatrix $\mathbf{D}$ \\
        \hline
        Methode & Maximiere Varianz & Reproduziere Distanzen \\
        \hline
        Dimensionen & Oft 2-3 & Typisch 2-3 \\
        \hline
        Interpretation & Varianz-Richtungen & Objekt-Ähnlichkeiten \\
        \hline
        Mathematik & Eigenwertzerlegung von Kovarianz & Eigenwertzerlegung von Gram-Matrix \\
        \hline
    \end{tabular}
\end{table}

\textbf{Praktische Anwendungen:}
\begin{itemize}[leftmargin=1.5em]
    \item \textbf{PCA:} Bioinformatik (Gen-Expression), Bildkompression, Feature Engineering für maschinelles Lernen.
    \item \textbf{PCR:} Hochdimensionale Prädiktoren mit Multikollinearität (z.\,B.\ Spektroskopie-Daten, Gene).
    \item \textbf{MDS:} Psychometrie (Subjektive Ähnlichkeitsurteile), Genetische Distanzen, Chemische Ähnlichkeitsmaße.
\end{itemize}

\newpage

\subsection{Lasso und Ridge-Regression}

\definition{\textbf{Regularisierung und Shrinkage-Methoden:} Techniken, die die Größe der Regressionskoeffizienten beschränken (reduzieren), um Overfitting zu vermeiden und stabilere Schätzungen zu erhalten. Dies ist besonders wichtig bei hochdimensionalen Daten oder Multikollinearität.}

\subsubsection{Ridge Regression}

\definition{\textbf{Motivation für Ridge Regression:}

\begin{itemize}[leftmargin=1.5em]
    \item In der klassischen OLS-Regression $\hat{\boldsymbol{\beta}} = (\mathbf{X}^T \mathbf{X})^{-1} \mathbf{X}^T \mathbf{y}$ ist die Invertierung der Matrix $\mathbf{X}^T \mathbf{X}$ numerisch instabil, wenn die Spalten von $\mathbf{X}$ (Prädiktoren) stark korreliert sind (\textbf{Multikollinearität}), oder wenn Eigenwerte von $\mathbf{X}^T \mathbf{X}$ nahe Null sind.
    \item Dies führt zu großen und instabilen Schätzungen $\hat{\boldsymbol{\beta}}$.
    \item \textbf{Ridge Regression:} Addiere eine kleine Konstante $\lambda > 0$ zu den Diagonalelementen von $\mathbf{X}^T \mathbf{X}$:
    \formel{\hat{\boldsymbol{\beta}}_{\text{ridge}} = (\mathbf{X}^T \mathbf{X} + \lambda \mathbf{I}_p)^{-1} \mathbf{X}^T \mathbf{y}}
    
    wobei $\mathbf{I}_p$ ist die $p \times p$ Identitätsmatrix und $\lambda \geq 0$ ist der \textbf{Regularisierungsparameter}.
\end{itemize}}

\definition{\textbf{Äquivalente Formulierungen der Ridge Regression:}

\textbf{1. Optimierungsperspektive:} Minimiere die Objective Function:
\formel{L_{\text{ridge}}(\boldsymbol{\beta}, \lambda) = \sum_{i=1}^n (y_i - \mathbf{x}_i^T \boldsymbol{\beta})^2 + \lambda \sum_{j=1}^p \beta_j^2 = \|\mathbf{y} - \mathbf{X} \boldsymbol{\beta}\|^2_2 + \lambda \|\boldsymbol{\beta}\|^2_2}

Der erste Term ist RSS (Anpassung an Daten), der zweite Term ist die \textbf{L2-Penalty} (Ridge-Penalty).

\textbf{2. Konstrained Optimization:} Äquivalent dazu ist zu minimieren:
\formel{\|\mathbf{y} - \mathbf{X} \boldsymbol{\beta}\|^2_2 \quad \text{unter Bedingung} \quad \sum_{j=1}^p \beta_j^2 \leq t}

für eine Schwelle $t \geq 0$. Die Beziehung zwischen $\lambda$ und $t$ ist inversiv.}

\definition{\textbf{Ridge als Shrinkage-Methode:} Für orthonormale Prädiktoren (= unabhängig, Einheitsvarianz) gilt:
\formel{\hat{\boldsymbol{\beta}}_{\text{ridge}} = \frac{1}{1 + \lambda} \hat{\boldsymbol{\beta}}_{\text{OLS}}}

\textbf{Interpretation:} Jeder Koeffizient wird um den Faktor $\frac{1}{1+\lambda}$ \textbf{gegen Null geschrumpft}:
\begin{itemize}[leftmargin=1.5em]
    \item Wenn $\lambda = 0$: Keine Shrinkage, identisch zu OLS.
    \item Wenn $\lambda \to \infty$: Maximale Shrinkage, $\hat{\boldsymbol{\beta}}_{\text{ridge}} \to \mathbf{0}$ (nur Intercept bleibt).
    \item Für allgemeine (nicht orthonormale) Prädiktoren: Richtungen mit kleineren Eigenwerten werden stärker geschrumpft.
\end{itemize}}

\definition{\textbf{Wahl des Regularisierungsparameters $\lambda$:} Der Parameter $\lambda$ kontrolliert den Bias-Variance Tradeoff:
\begin{itemize}[leftmargin=1.5em]
    \item \textbf{Kleine $\lambda$:} Niedriger Bias, höhere Varianz (ähnlich wie OLS).
    \item \textbf{Große $\lambda$:} Höherer Bias, niedrigere Varianz (stärkere Regularisierung).
\end{itemize}

Standard-Ansatz: \textbf{Cross-Validation} (z.\,B.\ 5-fold oder 10-fold CV) über ein Gitter von $\lambda$-Werten. Wähle das $\lambda$, das den kleinsten CV-Fehler liefert.}

\subsubsection{Lasso Regression}

\definition{\textbf{The Lasso (Least Absolute Shrinkage and Selection Operator):} Alternative zu Ridge mit einer anderen Penalty.

\textbf{Objective Function:}
\formel{L_{\text{lasso}}(\boldsymbol{\beta}, \lambda) = \sum_{i=1}^n (y_i - \mathbf{x}_i^T \boldsymbol{\beta})^2 + \lambda \sum_{j=1}^p |\beta_j| = \|\mathbf{y} - \mathbf{X} \boldsymbol{\beta}\|^2_2 + \lambda \|\boldsymbol{\beta}\|_1}

Der entscheidende Unterschied: \textbf{L1-Penalty} $\|\boldsymbol{\beta}\|_1 = \sum_j |\beta_j|$ statt L2-Penalty.

\textbf{Equivalente Konstrained Form:}
\formel{\min_{\boldsymbol{\beta}} \|\mathbf{y} - \mathbf{X} \boldsymbol{\beta}\|^2_2 \quad \text{unter} \quad \sum_{j=1}^p |\beta_j| \leq t}}

\definition{\textbf{Kritischer Unterschied Ridge vs. Lasso:} Die L1-Penalty hat eine wichtige Eigenschaft:
\begin{itemize}[leftmargin=1.5em]
    \item \textbf{Ridge (L2):} Schrumpft alle Koeffizienten gegen Null, setzt sie aber (fast) nie exakt zu Null.
    \item \textbf{Lasso (L1):} \textbf{Setzt einige Koeffizienten exakt zu Null!} Dies führt zu \textbf{Variablen-Auswahl} (feature selection).
\end{itemize}

\textbf{Geometrische Intuition:} 
\begin{itemize}[leftmargin=1.5em]
    \item Ridge-Constraint: Kreis $\sum \beta_j^2 \leq t$ (glatte, runde Form).
    \item Lasso-Constraint: Diamant/Rhombus $\sum |\beta_j| \leq t$ (spitze Ecken).
\end{itemize}

Die Ecken des Diamanten liegen auf den Achsen, daher ist die Wahrscheinlichkeit hoch, dass die Optimallösung an einer Ecke liegt → Koeffizient $= 0$.}

\definition{\textbf{Lasso -- Lösung und Variablen-Auswahl:} Im Gegensatz zu Ridge hat Lasso \textbf{keine geschlossene Lösung}. Die Lösung muss numerisch berechnet werden, z.\,B.\ via:
\begin{itemize}[leftmargin=1.5em]
    \item \textbf{Coordinate Descent:} Schneller Algorithmus, speziell für Lasso optimiert.
    \item \textbf{Proximal Gradient Descent:} Allgemeineres Verfahren.
    \item \textbf{ADMM (Alternating Direction Method of Multipliers):} Robust und skalierbar.
\end{itemize}

\textbf{Variablen-Auswahl-Effekt:} Wenn Lasso Koeffizienten auf exakt Null setzt, werden diese Variablen aus dem Modell entfernt. Dies ist besonders wertvoll wenn:
\begin{itemize}[leftmargin=1.5em]
    \item Viele Prädiktoren irrelevant sind (sparse data).
    \item Interpretierbarkeit wichtig ist (kleinere, sparsamere Modelle).
\end{itemize}}

\definition{\textbf{Wahl von $\lambda$ im Lasso:} Wie bei Ridge via Cross-Validation. \textbf{Besonderheit:} Es gibt oft einen Bereich von $\lambda$-Werten, die dieselbe Teilmenge von Variablen auswählen (die gleichen Koeffiziente zu Null setzen).}

\subsubsection{Ridge vs. Lasso vs. Elastic Net}

\textbf{Vergleich Ridge und Lasso:}
\begin{table}[h]
    \centering
    \small
    \begin{tabular}{|l|l|l|}
        \hline
        \textbf{Aspekt} & \textbf{Ridge} & \textbf{Lasso} \\
        \hline
        Penalty & $\lambda \sum \beta_j^2$ (L2) & $\lambda \sum |\beta_j|$ (L1) \\
        \hline
        Schätzung & Geschlossene Form & Numerisch (Coord. Descent) \\
        \hline
        Koeffiziente = 0? & Selten (fast nie) & Oft (Variablen-Auswahl) \\
        \hline
        Multikollinearität & Gut gehandhabt & Wählt eine Variable aus \\
        \hline
        Interpretierbarkeit & Alle Variablen im Modell & Sparse Modell \\
        \hline
        High $p$, low $n$ & Funktioniert & Funktioniert besser \\
        \hline
    \end{tabular}
\end{table}

\definition{\textbf{Elastic Net:} Ein Kompromiss zwischen Ridge und Lasso, der beide Penalties kombiniert.

\textbf{Objective Function:}
\formel{L_{\text{elastic}}(\boldsymbol{\beta}, \lambda, \alpha) = \sum_{i=1}^n (y_i - \mathbf{x}_i^T \boldsymbol{\beta})^2 + \lambda \left[ \alpha \sum_{j=1}^p |\beta_j| + (1-\alpha) \sum_{j=1}^p \beta_j^2 \right]}

wobei $\alpha \in [0, 1]$ ist der Mixing-Parameter:
\begin{itemize}[leftmargin=1.5em]
    \item $\alpha = 0$: Pure Ridge Regression.
    \item $\alpha = 1$: Pure Lasso Regression.
    \item $0 < \alpha < 1$: Hybrid Modell mit beiden Penalties.
\end{itemize}

\textbf{Vorteile von Elastic Net:}
\begin{itemize}[leftmargin=1.5em]
    \item Kombiniert Variable-Auswahl von Lasso mit Stabilität von Ridge.
    \item Wenn Prädiktoren stark korreliert sind, Lasso wählt oft nur eine Variable aus; Elastic Net wählt mehrere aber mit reduzierten Koeffizienten.
    \item Flexibler als beide Einzelmethoden.
\end{itemize}}

\subsubsection{Sparsity und Degrees of Freedom}

\definition{\textbf{Sparsity:} Die Eigenschaft, dass viele Koeffizienten exakt Null sind (nur wenige nicht-Null Koeffizienten).

\textbf{Effective Degrees of Freedom (df):} Eine Maßzahl für die Modellkomplexität:
\begin{itemize}[leftmargin=1.5em]
    \item In OLS: $\text{df} = p$ (Anzahl der Prädiktoren).
    \item In Lasso: $\text{df}$ = Anzahl der nicht-Null Koeffizienten (variiert mit $\lambda$).
    \item In Ridge: $\text{df} = \text{trace}(\mathbf{X}(\mathbf{X}^T \mathbf{X} + \lambda \mathbf{I})^{-1} \mathbf{X}^T)$ (irgendwas zwischen 0 und $p$).
\end{itemize}

Dies ist relevant für Modellvergleiche und AIC/BIC.}

\definition{\textbf{Pfad der Lösungen (Solution Path):} Oft visualisiert man wie Koeffizienten sich ändern wenn $\lambda$ variiert:
\begin{itemize}[leftmargin=1.5em]
    \item \textbf{Ridge-Pfad:} Alle Koeffizienten schrumpfen kontinuierlich gegen Null (keine Diskontinuitäten).
    \item \textbf{Lasso-Pfad:} Koeffizienten springen auf Null, wenn $\lambda$ einen kritischen Punkt überschreitet (Diskontinuitäten, ``sparse zones'').
    \item Ein solcher Plot (oder Tabelle) gibt Einsicht welche Variablen zuerst ausgewählt werden und bei welchem $\lambda$.
\end{itemize}}

\textbf{Praktische Empfehlungen:}
\begin{itemize}[leftmargin=1.5em]
    \item \textbf{Ridge:} Wenn alle Prädiktoren potenziell wichtig sind, oder hochgradig korreliert (z.\,B.\ Spektroskopie, Gen-Expression).
    \item \textbf{Lasso:} Wenn Variablen-Auswahl gewünscht ist oder $p \gg n$ (viel mehr Prädiktoren als Beobachtungen).
    \item \textbf{Elastic Net:} Kompromiss, wenn Unsicherheit über die beste Methode besteht, oder Daten sowohl Korrelation als auch Sparsity haben.
    \item \textbf{Cross-Validation:} Immer nutzen um $\lambda$ und $\alpha$ zu wählen!
\end{itemize}

\newpage

\subsection{Entscheidungsbäume (Decision Trees) und Random Forests}

\subsubsection{Decision Trees (Entscheidungsbäume)}

\definition{\textbf{Decision Tree:} Ein superwachter Lernalgorithmus für \textbf{Klassifikation} und \textbf{Regression}, der Daten rekursiv in immer homogenere Teilmengen aufteilt, um eine baumähnliche Entscheidungsstruktur zu erzeugen.

\textbf{Grundstruktur:}
\begin{itemize}[leftmargin=1.5em]
    \item \textbf{Knoten (Nodes):} Entscheidungspunkte, an denen eine Bedingung auf einer Variablen überprüft wird (z.\,B.\ $x_j \leq c$).
    \item \textbf{Äste (Branches):} Kanten zwischen Knoten, die den zwei Ausgängen einer Bedingung entsprechen (ja/nein, true/false).
    \item \textbf{Blattknoten (Leaf Nodes):} Endknoten ohne weitere Verzweigungen. In Klassifikation: Enthält eine Klasse; in Regression: Enthält einen Vorhersagewert.
    \item \textbf{Wurzelbene (Root):} Der oberste Knoten, mit dem die Aufteilung beginnt.
\end{itemize}

\textbf{Vorhersage:} Ein neuer Datenpunkt fließt von der Wurzel durch den Baum, indem er bei jedem Knoten der entsprechenden Bedingung folgt, bis er ein Blatt erreicht.}

\definition{\textbf{Baumaufbau: Recursive Partitioning}

\textbf{Algorithmus (Top-Down):}
\begin{enumerate}[leftmargin=2em]
    \item \textbf{Starte:} Mit allen Trainingsdaten an der Wurzel.
    \item \textbf{Finde beste Aufteilung:} Wähle die Variable $x_j$ und den Schwellwert $c$, der nach der Aufteilung die \textbf{reinsten} (homogensten) Kindknoten erzeugt.
    \item \textbf{Teile:} Unterteile Daten in zwei Gruppen: $x_j \leq c$ (linker Ast) und $x_j > c$ (rechter Ast).
    \item \textbf{Rekursion:} Wiederhole Schritte 2--3 für jeden Kindknoten (solange Stoppkriterien nicht erreicht).
    \item \textbf{Basisfall:} Knoten werden zu Blättern, wenn:
    \begin{itemize}[leftmargin=1.5em]
        \item Alle Daten derselbe Klasse (100\% rein).
        \item Minimale Anzahl Beobachtungen pro Knoten erreicht.
        \item Maximale Baumtiefe erreicht.
        \item Keine weitere Verbesserung möglich.
    \end{itemize}
\end{enumerate}}

\definition{\textbf{Reine-Kriterien (Impurity Measures):}

\textbf{In der Klassifikation:}

\textbf{1. Gini-Index:} Wahrscheinlichkeit, dass ein zufällig gewählter Punkt falsch klassifiziert wird, wenn man die Klasse zufällig nach der Klassenverteilung wählt.
\formel{\text{Gini} = 1 - \sum_{k=1}^K p_k^2}
wobei $p_k$ ist der Anteil der Klasse $k$ im Knoten. Gini = 0 wenn rein (ein Klasse), Gini = $\frac{K-1}{K}$ wenn uniform verteilt.

\textbf{2. Entropie (Information Gain):} Maß für Unordnung basierend auf Informationstheorie.
\formel{\text{Entropy} = -\sum_{k=1}^K p_k \log_2(p_k)}
Wird oft mit dem \textbf{Information Gain} kombiniert:
\formel{\text{IG}(\text{Split}) = \text{Entropy}(\text{Parent}) - \sum_{i} \frac{n_i}{n} \text{Entropy}(\text{Child}_i)}
Der Information Gain misst die Reduzierung der Entropie durch eine Aufteilung.

\textbf{In der Regression:}

\textbf{Fehlerreduktion (Least Squares):} Wähle die Aufteilung, die die \textbf{Residualsumme der Quadrate (RSS)} maximiert reduziert.
\formel{\text{RSS} = \sum_i (y_i - \hat{y})^2}
Nach einer Aufteilung werden die RSS in beiden Kindknoten berechnet, und die beste Aufteilung minimiert die kombinierte RSS.}

\definition{\textbf{Overfitting und Pruning:}

\textbf{Problem:} Entscheidungsbäume können sehr tief werden und perfekt die Trainingsdaten fitten, generalisieren aber schlecht auf neue Daten (Overfitting).

\textbf{Lösungsansätze:}
\begin{itemize}[leftmargin=1.5em]
    \item \textbf{Early Stopping:} Begrenzen der Baumtiefe, minimale Anzahl Beobachtungen pro Knoten, minimales Reine-Kriterium.
    \item \textbf{Post-Pruning (Nachträgliches Beschneiden):} Trainiere einen vollen Baum, dann entferne Knoten von unten nach oben, wenn dies die Validierungs-Performance verbessert (z.\,B.\ via \textbf{Reduced Error Pruning}).
    \item \textbf{Cost Complexity Pruning:} Minimiere Kosten-Funktion, die Trainings-Fehler plus Komplexitäts-Strafe umfasst.
\end{itemize}}

\textbf{Vorteile von Decision Trees:}
\begin{itemize}[leftmargin=1.5em]
    \item \textbf{Interpretierbarkeit:} Die Entscheidungslogik ist transparent (einfach zu visualisieren und zu erklären).
    \item \textbf{Keine Skalierung nötig:} Funktioniert mit rohen Daten, unempfindlich gegenüber Monoton-Transformationen.
    \item \textbf{Gemischte Datentypen:} Kann numerische und kategoriale Prädiktoren direkt verarbeiten.
    \item \textbf{Nichtlinearität:} Erfasst Interaktionen und komplexe Muster automatisch.
\end{itemize}

\textbf{Nachteile von Decision Trees:}
\begin{itemize}[leftmargin=1.5em]
    \item \textbf{Instabilität:} Kleine Änderungen in den Trainingsdaten können zu völlig anderen Bäumen führen.
    \item \textbf{Overfitting:} Neigung zu großen Bäumen ohne Regularisierung.
    \item \textbf{Bias bei unbalancierten Klassen:} Tendenziell zu Mehrheitsklasse vorgespannt.
    \item \textbf{Greedy-Algorithmus:} Macht lokal optimale Entscheidungen (nicht global optimal).
\end{itemize}

\subsubsection{Random Forests}

\definition{\textbf{Random Forest:} Ein \textbf{Ensemble-Lernalgorithmus}, der viele Entscheidungsbäume trainiert und ihre Vorhersagen kombiniert, um die Genauigkeit und Stabilität zu verbessern.

\textbf{Kernidee:} 
\begin{itemize}[leftmargin=1.5em]
    \item ``Weisheit der Massen'': Viele unabhängige (oder gering korrelierte) Bäume treffen bessere Entscheidungen als ein einzelner Baum.
    \item \textbf{Bootstrap-Sampling:} Trainiere jeden Baum auf einer zufälligen Stichprobe (mit Zurücklegen) der Trainingsdaten.
    \item \textbf{Random Feature Selection:} Bei jedem Knoten wähle aus einem zufälligen Untermengen von Prädiktoren die beste Aufteilung.
\end{itemize}}

\definition{\textbf{Random Forest Algorithmus:}

\begin{enumerate}[leftmargin=2em]
    \item \textbf{Schleife von $t = 1$ bis $B$ (Anzahl Bäume):}
    \begin{itemize}
        \item Ziehe eine Bootstrap-Stichprobe der Größe $n$ (mit Zurücklegen) aus den Trainingsdaten.
        \item Trainiere einen vollständigen (ungekürzten) Decision Tree auf dieser Stichprobe.
        \item Verwende bei jedem Knoten nur $m < p$ zufällig ausgewählte Prädiktoren bei der Suche nach der besten Aufteilung (üblicherweise $m \approx \sqrt{p}$ für Klassifikation, $m \approx p/3$ für Regression).
    \end{itemize}
    
    \item \textbf{Vorhersage:}
    \begin{itemize}
        \item \textbf{Klassifikation:} Nimm die Mehrheits-Stimme aller $B$ Bäume: $\hat{y} = \text{Modus}(\hat{y}_1, \ldots, \hat{y}_B)$.
        \item \textbf{Regression:} Nimm den Durchschnitt: $\hat{y} = \frac{1}{B} \sum_{t=1}^B \hat{y}_t$.
    \end{itemize}
\end{enumerate}}

\definition{\textbf{Out-of-Bag (OOB) Error Estimation:}

\textbf{Vorteil:} Da Bootstrap-Stichproben verwendet werden, sind ungefähr $e^{-1} \approx 37\%$ der Trainingsdaten nicht in jeder Stichprobe enthalten (``out-of-bag''). Diese können zur Validierung verwendet werden, ohne einen separaten Validierungssatz zu benötigen.

\textbf{Verfahren:}
\begin{enumerate}[leftmargin=1.5em]
    \item Für jede Beobachtung: Identifiziere alle Bäume, in deren Bootstrap-Stichprobe sie \textbf{nicht} enthalten war.
    \item Vorhersage: Gemittelt (oder Mehrheits-Stimme) dieser Bäume für diese Beobachtung.
    \item \textbf{OOB Error}: Der Fehler dieser ``Out-of-Bag''-Vorhersagen ist eine gute Schätzung der Testfehler.
\end{enumerate}}

\definition{\textbf{Feature Importance (Variablenbedeutung):}

\textbf{Mean Decrease in Impurity (MDI):} Für jeden Prädiktor, summiere über alle Bäume und alle Knoten, um wie viel der Impurity (Gini oder RSS) bei Aufteilungen in diesem Prädiktor reduziert wurde.

\formel{\text{Importance}(X_j) = \frac{1}{B} \sum_{t=1}^B \sum_{\text{nodes: split on } X_j} \text{(Impurity reduction)}}

Höhere Werte zeigen, dass die Variable häufig und früh (near root) zu großen Reduktionen führt, und daher important ist.

\textbf{Mean Decrease in Accuracy (MDA) / Permutation Importance:} Eine Alternative:
\begin{enumerate}[leftmargin=1.5em]
    \item Berechne OOB Error (oder Validierungs-Error) mit originalem Datensatz.
    \item Permutiere die Werte einer Variable $X_j$ zufällig, berechne OOB Error neu.
    \item Die Differenz zeigt, wie wichtig die Variable ist (größere Verschlechterung $\Rightarrow$ wichtigere Variable).
\end{enumerate}}

\textbf{Vorteile von Random Forests:}
\begin{itemize}[leftmargin=1.5em]
    \item \textbf{Robustheit:} Bootstrap + Ensemble reduzieren Varianz und Overfitting.
    \item \textbf{Hohe Accuracy:} Oft eine der besten off-the-shelf Performance.
    \item \textbf{Robustness gegen Outlier:} Einzelne extreme Werte beeinflussen weniger als bei anderen Methoden.
    \item \textbf{Feature Importance:} Automatische Quantifizierung welche Variablen wichtig sind.
    \item \textbf{Parallelsierbar:} Bäume können unabhängig trainiert werden.
    \item \textbf{Keine Normalisierung nötig:} Wie Decision Trees, unempfindlich gegen Skalierung.
\end{itemize}

\textbf{Nachteile von Random Forests:}
\begin{itemize}[leftmargin=1.5em]
    \item \textbf{Interpretierbarkeit:} Schwer zu verstehen warum die Vorhersage gemacht wurde (``Black Box'').
    \item \textbf{Speicher und Rechenzeit:} Viele Bäume benötigen viel RAM und CPU.
    \item \textbf{Hyperparameter-Tuning:} Zu optimierende Parameter wie $B$ (Anzahl Bäume), $m$ (Features pro Split), Baumtiefe.
\end{itemize}

\definition{\textbf{Vergleich: Decision Trees vs.\ Random Forests}}
\begin{table}[h]
    \centering
    \small
    \begin{tabular}{|l|l|l|}
        \hline
        \textbf{Aspekt} & \textbf{Decision Tree} & \textbf{Random Forest} \\
        \hline
        Modell & Ein Baum & Ensemble von Bäumen \\
        \hline
        Varianz & Hoch & Niedrig (durch Ensemble) \\
        \hline
        Bias & Moderat & Ähnlich wie Baum \\
        \hline
        Interpretierbarkeit & Hoch & Niedrig \\
        \hline
        Performance & Moderat & Höher (meist) \\
        \hline
        Rechenzeit & Niedrig & Höher \\
        \hline
        Robust gegen Outlier & Schwach & Stark \\
        \hline
        Feature Importance & Einfach & Möglich (MDI, Permutation) \\
        \hline
    \end{tabular}
\end{table}

\definition{\textbf{Weitere Ensemble-Methoden:}

\begin{itemize}[leftmargin=1.5em]
    \item \textbf{Bagging (Bootstrap Aggregating):} Verallgemeinerung von Random Forests auf beliebige Base-Learner (nicht nur Bäume). Trainiere viele Modelle auf Bootstrap-Stichproben und aggregiere.
    
    \item \textbf{Boosting (z.\,B.\ AdaBoost, Gradient Boosting):} Sequenzielle Methode: Trainiere Bäume nacheinander, wobei jeder neue Baum sich auf die Fehler des vorigen konzentriert. Wichtung von Trainingsdaten wird angepasst. Oft bessere Performance als Bagging, aber anfälliger für Overfitting.
    
    \item \textbf{Gradient Boosting Machines (GBM):} Verallgemeinertes Boosting, das Gradienten-Abstieg verwendet. Führt zu sehr guten praktischen Ergebnissen, aber mehr Hyperparameter-Tuning.
    
    \item \textbf{Extreme Gradient Boosting (XGBoost):} Optimierte Implementierung von GBM mit zusätzlichen Regularisierungen. Sehr populär in Machine Learning Wettbewerben.
\end{itemize}}

\subsection{Linear Discriminant Analysis (LDA) und Quadratic Discriminant Analysis (QDA)}

\subsubsection{Klassifikation via Wahrscheinlichkeitsdichten}

\definition{\textbf{Klassifikations-Setting:}
\begin{itemize}[leftmargin=1.5em]
    \item \textbf{Gegeben:} Datenmatrix $\mathbf{X}$ der Dimension $n \times p$ (n Beobachtungen, p Prädiktoren) und kategoriale Response $y \in \{1, \ldots, K\}$ mit K Klassen.
    \item \textbf{Annahme:} Beobachtungen in Klasse $k$ entstehen aus einer Population $P_k$ mit Wahrscheinlichkeitsdichte $f_k(\mathbf{x})$.
    \item \textbf{Alternativ:} $f_k(\mathbf{x}) = f(\mathbf{x} | C = k)$ — die bedingte Dichte von $\mathbf{x}$ gegeben Klasse $k$.
    \item \textbf{Ziel:} Klassifikationsregel finden, die eine neue Beobachtung $\mathbf{x}_0$ einer der K Klassen zuordnet.
\end{itemize}}

\definition{\textbf{Maximum Likelihood (ML) Klassifikation:} Die einfachste Klassifikationsregel:
\formel{\hat{k} = \arg\max_{k=1,\ldots,K} f_k(\mathbf{x})}

Wir ordnen eine Beobachtung $\mathbf{x}$ derjenigen Klasse $k$ zu, für die die Wahrscheinlichkeitsdichte $f_k(\mathbf{x})$ maximal ist.

\textbf{Entscheidungsgrenze:} Die Werte $\mathbf{x}$, für die $f_1(\mathbf{x}) = f_2(\mathbf{x})$ (im binären Fall), bilden die Entscheidungsgrenze zwischen den Klassen.}

\definition{\textbf{Bayes Klassifikation mit Prior-Wahrscheinlichkeiten:} In der Praxis sollten wir auch die \textbf{Prior-Wahrscheinlichkeiten} $p_k = P(C = k)$ beachten (die Häufigkeit, mit der Klasse $k$ in der Population auftritt).

Die \textbf{Bayes-Regel}:
\formel{\hat{k} = \arg\max_{k=1,\ldots,K} f_k(\mathbf{x}) \cdot p_k}

Wir ordnen die Beobachtung $\mathbf{x}$ derjenigen Klasse $k$ zu, für welche das Produkt aus Likelihood $f_k(\mathbf{x})$ und Prior $p_k$ maximal ist.

\textbf{Intuition:} Auch wenn $f_2(\mathbf{x})$ größer als $f_1(\mathbf{x})$ ist, kann Klasse 1 die bessere Wahl sein, wenn Klasse 1 viel häufiger in der Population auftritt ($p_1 \gg p_2$).}

\subsubsection{Linear Discriminant Analysis (LDA)}

\definition{\textbf{LDA Annahmen:} 
\begin{itemize}[leftmargin=1.5em]
    \item \textbf{Normalverteilung:} Innerhalb jeder Klasse $k$ sind die Prädiktoren $\mathbf{x}$ multivariat normalverteilt: $\mathbf{x} | C = k \sim N_p(\boldsymbol{\mu}_k, \boldsymbol{\Sigma})$.
    \item \textbf{Gleiche Kovarianzmatrizen:} Alle K Klassen teilen dieselbe Kovarianzmatrix $\boldsymbol{\Sigma}$ (identische Streuungsstruktur).
    \item Dies führt zu einer \textbf{linearen} Entscheidungsgrenze.
\end{itemize}}

\definition{\textbf{LDA im univariaten Fall (p=1):}

Mit nur einem Prädiktor $x$ und gemeinsamer Varianz $\sigma^2$ in allen Klassen ist die bedingte Dichte:
\formel{f_k(x) = \frac{1}{\sqrt{2\pi\sigma^2}} \exp\left(-\frac{1}{2\sigma^2}(x - \mu_k)^2\right)}

Nach der Bayes-Regel klassifizieren wir $x$ zu Klasse $k$, für welche die \textbf{Diskriminanzfunktion}
\formel{\delta_k(x) = x \cdot \frac{\mu_k}{\sigma^2} - \frac{\mu_k^2}{2\sigma^2} + \log(p_k)}
maximal ist. Diese ist \textbf{linear in x}.

\textbf{Einfacher Fall (gleiche Prior):} Wenn $p_k = 1/K$ für alle $k$, vereinfacht sich die Regel zu: wähle die Klasse mit dem Mittelwert $\mu_k$, der $x$ am nächsten liegt.}

\definition{\textbf{LDA im multivariaten Fall (p > 1):}

Mit p Prädiktoren und gemeinsamer Kovarianzmatrix $\boldsymbol{\Sigma}$ (Größe $p \times p$) ist die Diskriminanzfunktion für Klasse $k$:
\formel{\delta_k(\mathbf{x}) = \mathbf{x}^T \boldsymbol{\Sigma}^{-1} \boldsymbol{\mu}_k - \frac{1}{2} \boldsymbol{\mu}_k^T \boldsymbol{\Sigma}^{-1} \boldsymbol{\mu}_k + \log(p_k)}

Dies ist eine \textbf{affine Funktion} (linear + Konstante) von $\mathbf{x}$.

\textbf{Klassifikationsregel:} 
\formel{\hat{k} = \arg\max_{k=1,\ldots,K} \delta_k(\mathbf{x})}

Entscheidungsgrenzen zwischen Klassen $k$ und $k'$ entstehen, wo $\delta_k(\mathbf{x}) = \delta_{k'}(\mathbf{x})$, und sind \textbf{lineare Hyperebenen} (Geraden in 2D, Ebenen in 3D, etc.).}

\definition{\textbf{Parameterschätzung in LDA:}

Aus Trainingsdaten schätzen wir:
\begin{itemize}[leftmargin=1.5em]
    \item \textbf{Klassenmittelwerte:} $\hat{\boldsymbol{\mu}}_k = \frac{1}{n_k} \sum_{i: y_i = k} \mathbf{x}_i$ (Mittelwert aller Beobachtungen in Klasse $k$).
    \item \textbf{Prior-Wahrscheinlichkeiten:} $\hat{p}_k = \frac{n_k}{n}$ (Anteil der Klasse $k$).
    \item \textbf{Gepoolte Kovarianzmatrix:} 
    \formel{\hat{\boldsymbol{\Sigma}}_{\text{pooled}} = \frac{1}{n - K} \sum_{k=1}^K (n_k - 1) \hat{\boldsymbol{\Sigma}}_k}
    wobei $\hat{\boldsymbol{\Sigma}}_k$ die Stichproben-Kovarianzmatrix innerhalb Klasse $k$ ist.
\end{itemize}}

\textbf{Vorteile von LDA:}
\begin{itemize}[leftmargin=1.5em]
    \item \textbf{Einfach:} Geschlossene Lösung, stabil zu berechnen.
    \item \textbf{Interpretierbar:} Lineare Entscheidungsgrenzen sind leicht zu verstehen.
    \item \textbf{Effizient:} Schnell selbst für hochdimensionale Daten.
    \item \textbf{Gut bei kleinen Stichproben:} Benötigt relativ wenig Daten für gute Schätzung.
\end{itemize}

\textbf{Nachteile von LDA:}
\begin{itemize}[leftmargin=1.5em]
    \item \textbf{Normalverteilungsannahme:} Kann bei nicht-normalen Daten schlecht funktionieren.
    \item \textbf{Gleiche Kovarianzmatrizen:} Annahme ist oft zu restriktiv; in Realität haben Klassen oft unterschiedliche Spreads.
    \item \textbf{Lineare Grenzen:} Nicht flexibel genug für komplexe, nichtlineare Entscheidungsgrenzen.
\end{itemize}

\subsubsection{Quadratic Discriminant Analysis (QDA)}

\definition{\textbf{QDA Annahmen:}
\begin{itemize}[leftmargin=1.5em]
    \item \textbf{Normalverteilung:} Wie in LDA, $\mathbf{x} | C = k \sim N_p(\boldsymbol{\mu}_k, \boldsymbol{\Sigma}_k)$.
    \item \textbf{Unterschiedliche Kovarianzmatrizen:} Jede Klasse $k$ hat ihre eigene Kovarianzmatrix $\boldsymbol{\Sigma}_k$ (unterschiedliche Streuungsstrukturen pro Klasse).
    \item Dies führt zu einer \textbf{quadratischen} Entscheidungsgrenze.
\end{itemize}}

\definition{\textbf{QDA Diskriminanzfunktion:}

Die Diskriminanzfunktion für Klasse $k$ wird:
\formel{\delta_k(\mathbf{x}) = -\frac{1}{2} \log|\boldsymbol{\Sigma}_k| - \frac{1}{2}(\mathbf{x} - \boldsymbol{\mu}_k)^T \boldsymbol{\Sigma}_k^{-1}(\mathbf{x} - \boldsymbol{\mu}_k) + \log(p_k)}

\textbf{Wichtig:} Der zweite Term enthält jetzt $\mathbf{x}^T \boldsymbol{\Sigma}_k^{-1} \mathbf{x}$ (quadratisch in $\mathbf{x}$), nicht wie in LDA die globale gepoolte Form. Dies erzeugt \textbf{quadratische Entscheidungsgrenzen}.

Klassifikationsregel:
\formel{\hat{k} = \arg\max_{k=1,\ldots,K} \delta_k(\mathbf{x})}

Entscheidungsgrenzen zwischen zwei Klassen sind \textbf{Quadriken} (Parabeln in 2D, Paraboloide in 3D, Hyperflächen zweiten Grades in höheren Dimensionen).}

\definition{\textbf{Parameterschätzung in QDA:}

\begin{itemize}[leftmargin=1.5em]
    \item \textbf{Klassenmittelwerte und Prior:} Wie in LDA, $\hat{\boldsymbol{\mu}}_k$ und $\hat{p}_k$.
    \item \textbf{Klassenspezifische Kovarianzmatrizen:} 
    \formel{\hat{\boldsymbol{\Sigma}}_k = \frac{1}{n_k - 1} \sum_{i: y_i = k} (\mathbf{x}_i - \hat{\boldsymbol{\mu}}_k)(\mathbf{x}_i - \hat{\boldsymbol{\mu}}_k)^T}
    (Stichproben-Kovarianzmatrix innerhalb Klasse $k$, ohne Pooling).
\end{itemize}}

\definition{\textbf{Vergleich: LDA vs. QDA}}

\begin{table}[h]
    \centering
    \small
    \begin{tabular}{|l|l|l|}
        \hline
        \textbf{Aspekt} & \textbf{LDA} & \textbf{QDA} \\
        \hline
        Annahme Kovarianz & Gleich $\boldsymbol{\Sigma}$ & Unterschiedlich $\boldsymbol{\Sigma}_k$ \\
        \hline
        Entscheidungsgrenze & Linear (Hyperplane) & Quadratisch (Quadrik) \\
        \hline
        Flexibilität & Niedrig & Höher \\
        \hline
        Parameter zu schätzen & Wenige & Viele (pro Klasse) \\
        \hline
        Stichproben erforderlich & Weniger & Mehr \\
        \hline
        Bias/Varianz & Höher Bias, Niedrig Varianz & Niedrig Bias, Höher Varianz \\
        \hline
        Wann besser? & Ähnliche Klassen-Spreads & Unterschiedliche Spreads \\
        \hline
    \end{tabular}
\end{table}

\textbf{Vor- und Nachteile von QDA:}

\textbf{Vorteile:}
\begin{itemize}[leftmargin=1.5em]
    \item \textbf{Flexibler:} Erlaubt unterschiedliche Kovarianzen pro Klasse — realistischere Modellierung.
    \item \textbf{Besser bei heterogenen Klassen:} Wenn Klassen unterschiedliche Streuungen/Formen haben.
\end{itemize}

\textbf{Nachteile:}
\begin{itemize}[leftmargin=1.5em]
    \item \textbf{Viele Parameter:} Bei $p$ Prädiktoren und $K$ Klassen sind insgesamt $K \cdot p(p+1)/2$ Kovarianzparameter zu schätzen — Overfitting-Risiko!
    \item \textbf{Kleine Stichproben:} Für kleine $n$ relativ zu $p$ und $K$ schwierig, stabil zu schätzen (schlecht gestelltes Problem).
    \item \textbf{Rechenkomplexität:} Höher als LDA (Invertieren mehrerer Matrizen).
\end{itemize}

\subsubsection{Varianten und verwandte Methoden}

\definition{\textbf{Fisher's Linear Discriminant:} Ein anderer Zugang (speziell für K=2, binäre Klassifikation) findet die \textbf{optimale Projektionsrichtung} $\mathbf{w}$, die die Klassen maximal trennt (maximale Separabilität).

\textbf{Ziel:} Maximiere das Verhältnis (Between-Class zu Within-Class Variabilität):
\formel{J(\mathbf{w}) = \frac{(\bar{\mathbf{x}}_1 - \bar{\mathbf{x}}_2)^2}{\sigma_1^2 + \sigma_2^2}}

Die Lösung ist proportional zu:
\formel{\mathbf{w} \propto \boldsymbol{\Sigma}_{\text{pooled}}^{-1}(\bar{\mathbf{x}}_1 - \bar{\mathbf{x}}_2)}

\textbf{Unterschied zu LDA:} Fisher braucht keine Normalverteilungsannahme und ist konzeptionell auf Dimensionsreduktion ausgerichtet (Projektion auf diese eine Richtung), nicht primär auf Klassifikation.}

\definition{\textbf{Regularized Discriminant Analysis (RDA):} Ein Mittelweg zwischen LDA und QDA durch Regularisierung:
\formel{\hat{\boldsymbol{\Sigma}}_k(\lambda) = \lambda \hat{\boldsymbol{\Sigma}}_{\text{pooled}} + (1-\lambda) \hat{\boldsymbol{\Sigma}}_k}

Der Parameter $\lambda \in [0, 1]$ interpoliert:
\begin{itemize}[leftmargin=1.5em]
    \item $\lambda = 1$: Vollständiges LDA (gepoolte Kovarianz, wenige Parameter).
    \item $\lambda = 0$: Vollständiges QDA (klassenspezifische Kovarianzen, viele Parameter).
    \item $0 < \lambda < 1$: Regularisierter Kompromiss (reduziertes Overfitting-Risiko, erhöhte Flexibilität).
\end{itemize}

Der optimale $\lambda$ wird typischerweise via Cross-Validation bestimmt.}

\definition{\textbf{Mixture Discriminant Analysis (MDA):} Erweitert LDA/QDA, indem jede Klasse selbst als \textbf{Mischung} (Mixture) mehrerer normalverteilter Komponenten modelliert wird:

\formel{f_k(\mathbf{x}) = \sum_{m=1}^{M_k} \pi_{km} \cdot N_p(\boldsymbol{\mu}_{km}, \boldsymbol{\Sigma}_{km})}

Dies ermöglicht komplexere, multimodale Klassenverteilungen statt nur eine Normalverteilung pro Klasse.}

\textbf{Praktische Anwendungen — wann welche Methode:}
\begin{itemize}[leftmargin=1.5em]
    \item Wähle \textbf{LDA}, wenn:
    \begin{itemize}
        \item Klassen ähnliche Kovarianzstrukturen haben.
        \item Daten einigermaßen normalverteilt sind.
        \item Kleine bis moderate Stichprobengröße.
        \item Interpretierbarkeit wichtig ist.
    \end{itemize}
    
    \item Wähle \textbf{QDA}, wenn:
    \begin{itemize}
        \item Klassen unterschiedliche Kovarianzen haben (erkennbar durch explorative Analyse).
        \item Genug Trainingsdaten vorhanden ($n \gg p \cdot K$).
        \item Flexibler ist besser als einfach.
    \end{itemize}
    
    \item Wähle \textbf{RDA}, wenn:
    \begin{itemize}
        \item Unsicher zwischen LDA und QDA.
        \item Overfitting befürchtet.
        \item Cross-Validation durchgeführt werden kann.
    \end{itemize}
    
    \item Ziehe \textbf{andere Klassifikatoren} in Betracht, wenn:
    \begin{itemize}
        \item Normalverteilungsannahmen stark verletzt sind.
        \item Nichtlineare Entscheidungsgrenzen nötig sind (z.\,B.\ SVM, KNN, Tree-based).
        \item Keine interpretierbare lineare/quadratische Regel nötig.
    \end{itemize}
\end{itemize}

\newpage

\subsection{k-Nearest Neighbors (kNN) — Klassifikation und Regression}

\definition{\textbf{kNN Grundkonzept:} Eine der \textbf{einfachsten Klassifikationsmethoden}, basierend auf räumlicher Nähe zwischen Objekten:}
\begin{itemize}[leftmargin=1.5em]
    \item \textbf{Lazy Learner / Memory-based:} Es gibt keinen eigentlichen ``Lern-Prozess'' während des Trainings — die Trainingsdaten werden einfach gespeichert.
    \item \textbf{Vorhersage-Phase:} Für eine neue Beobachtung $\mathbf{x}_0$ werden die $k$ \textbf{nächsten Nachbarn} in den Trainingsdaten gefunden (basierend auf einer Distanzmetrik).
    \item \textbf{Klassifikation:} Ein Mehrheitsabstimmung (Majority Vote) der Labels der $k$ Nachbarn bestimmt die Klasse von $\mathbf{x}_0$.
    \item \textbf{Regression:} Der Durchschnitt der Zielwerte der $k$ Nachbarn ist die Vorhersage.
\end{itemize}

\textbf{Algorithmus (Klassifikation):}
\begin{enumerate}[leftmargin=1.5em]
    \item Speichere alle $n$ Trainingsdaten $(\mathbf{x}_i, y_i)$.
    \item Für neue Beobachtung $\mathbf{x}_0$:
    \begin{itemize}
        \item Berechne die Distanz $d(\mathbf{x}_i, \mathbf{x}_0)$ zu allen Trainingsobjekten $i = 1, \ldots, n$.
        \item Identifiziere die $k$ Objekte mit den kleinsten Distanzen (die $k$ Nearest Neighbors).
        \item Klassifiziere $\mathbf{x}_0$ als diejenige Klasse, die unter diesen $k$ Nachbarn am häufigsten vorkommt (Mehrheitsabstimmung).
    \end{itemize}
\end{enumerate}

\subsubsection{Distanzmaße}

\definition{\textbf{Distanzfunktion:} Eine Funktion $d: \mathbb{R}^p \times \mathbb{R}^p \to [0, \infty)$ muss vier Bedingungen erfüllen (Metrik):}
\begin{enumerate}[leftmargin=1.5em]
    \item $d(\mathbf{x}, \mathbf{y}) \geq 0$ (Nicht-Negativität).
    \item $d(\mathbf{x}, \mathbf{y}) = 0 \Leftrightarrow \mathbf{x} = \mathbf{y}$ (Identität).
    \item $d(\mathbf{x}, \mathbf{y}) = d(\mathbf{y}, \mathbf{x})$ (Symmetrie).
    \item $d(\mathbf{x}, \mathbf{z}) \leq d(\mathbf{x}, \mathbf{y}) + d(\mathbf{y}, \mathbf{z})$ (Dreiecksungleichung).
\end{enumerate}

\definition{\textbf{Euklidische Distanz (am häufigsten):}
\formel{d_{\text{Euklidisch}}(\mathbf{x}_A, \mathbf{x}_B) = \sqrt{\sum_{j=1}^p (x_{A,j} - x_{B,j})^2} = \|\mathbf{x}_A - \mathbf{x}_B\|_2}

Dies ist die Standard-Luftliniendistanz im $\mathbb{R}^p$.}

\definition{\textbf{Manhattan-Distanz (L1-Norm):}
\formel{d_{\text{Manhattan}}(\mathbf{x}_A, \mathbf{x}_B) = \sum_{j=1}^p |x_{A,j} - x_{B,j}| = \|\mathbf{x}_A - \mathbf{x}_B\|_1}

Summe der absoluten Differenzen — wie man sich in einem Stadtgitter bewegen würde.}

\definition{\textbf{Minkowski-Distanz (Verallgemeinerung):}
\formel{d_p(\mathbf{x}_A, \mathbf{x}_B) = \left(\sum_{j=1}^p |x_{A,j} - x_{B,j}|^p\right)^{1/p} = \|\mathbf{x}_A - \mathbf{x}_B\|_p}

\begin{itemize}[leftmargin=1.5em]
    \item $p=1$: Manhattan-Distanz.
    \item $p=2$: Euklidische Distanz.
    \item $p=\infty$: Chebyshev-Distanz (Maximum der absoluten Differenzen).
\end{itemize}}

\definition{\textbf{Mahalanobis-Distanz (berücksichtigt Korrelation):}
\formel{d_{\text{Mahal}}(\mathbf{x}_A, \mathbf{x}_B) = \sqrt{(\mathbf{x}_A - \mathbf{x}_B)^T \boldsymbol{\Sigma}^{-1} (\mathbf{x}_A - \mathbf{x}_B)}}

wobei $\boldsymbol{\Sigma}$ die Kovarianzmatrix der Trainingsdaten ist. Dies berücksichtigt die Korrelation und Skalierung der Variablen — bessere Distanz als Euklidisch bei korrelierten Prädiktoren.}

\definition{\textbf{Tanimoto-Distanz (für binäre Vektoren):}

Für binäre Vektoren $\mathbf{x}_A, \mathbf{x}_B \in \{0,1\}^p$ (z.\,B.\ zur Charakterisierung chemischer Strukturen):

\textbf{Tanimoto-Ähnlichkeit:}
\formel{t_{AB} = \frac{\sum_{j=1}^p AND(x_{A,j}, x_{B,j})}{\sum_{j=1}^p OR(x_{A,j}, x_{B,j})} = \frac{\mathbf{x}_A^T \cdot \mathbf{x}_B}{\mathbf{x}_A^T \cdot \mathbf{1} + \mathbf{x}_B^T \cdot \mathbf{1} - \mathbf{x}_A^T \cdot \mathbf{x}_B}}

\textbf{Tanimoto-Distanz:}
\formel{d_{\text{Tanimoto}}(\mathbf{x}_A, \mathbf{x}_B) = 1 - t_{AB} = \frac{\sum_{j=1}^p XOR(x_{A,j}, x_{B,j})}{\sum_{j=1}^p OR(x_{A,j}, x_{B,j})}}

Zählt die Anzahl der Positionen, wo sich die Vektoren unterscheiden, geteilt durch die Anzahl der Positionen mit mind. einer 1.}

\textbf{Wichtig:} Die Wahl der Distanzmetrik kann großen Einfluss auf kNN-Ergebnisse haben. Standardmäßig wird Euklidisch verwendet, aber die Daten sollten ggfs. \textbf{normalisiert/standardisiert} werden (besonders wichtig!), da große Unterschiede in Skalierung die Distanz dominieren.

\subsubsection{Wahl von k und Entscheidungsgrenzen}

\definition{\textbf{Einfluss von k auf Entscheidungsgrenzen:}}
\begin{itemize}[leftmargin=1.5em]
    \item \textbf{k = 1:} Sehr \textbf{flexibel}, komplexe Entscheidungsgrenzen, raue/zackige Grenzen. \textbf{Hohes Overfitting-Risiko}.
    \item \textbf{Kleines k (z.\,B.\ k=3, 5):} Flexibel, folgt Trainingsdaten eng nach. Anfällig für Noise/Ausreißer.
    \item \textbf{Mittleres k (z.\,B.\ k=10--30):} Kompromiss zwischen Bias und Varianz.
    \item \textbf{Großes k (z.\,B.\ k=n):} Sehr \textbf{steif}, einfache, glatte Entscheidungsgrenzen. \textbf{Hohes Underfitting-Risiko}. Im Extremfall (k=n) wird alles der Mehrheitsklasse zugeordnet.
\end{itemize}

\textbf{Bias-Varianz Tradeoff bei kNN:}
\begin{itemize}[leftmargin=1.5em]
    \item Kleineres k: Höhere Varianz (Modell reagiert stark auf Trainingsdaten-Fluktuationen), niedrigerer Bias.
    \item Größeres k: Niedrigere Varianz (stabiler), höherer Bias (glattere, einfachere Grenzen).
\end{itemize}

\textbf{Wahl von k:}
\begin{itemize}[leftmargin=1.5em]
    \item \textbf{Cross-Validation:} Teste verschiedene $k$-Werte und wähle denjenigen mit bestem CV-Error.
    \item \textbf{Faustregel:} $k = \sqrt{n}$ oder $k = \sqrt{n}/2$, aber dies ist nur eine Orientierung.
    \item \textbf{Ungerade k:} Bei binärer Klassifikation hilft ungerader k um Unentschieden zu vermeiden.
\end{itemize}

\subsubsection{Vor- und Nachteile von kNN}

\textbf{Vorteile:}
\begin{itemize}[leftmargin=1.5em]
    \item \textbf{Sehr einfach:} Konzeptuell trivial zu verstehen.
    \item \textbf{Nonparametrisch:} Keine Verteilungsannahmen (anders als LDA/QDA).
    \item \textbf{Flexibel:} Kann komplexe, nichtlineare Entscheidungsgrenzen modellieren.
    \item \textbf{Für Klassifikation und Regression:} Funktioniert für beide Aufgaben.
    \item \textbf{Keine Trainingsphase:} Sofort einsatzbereit (Lazy Learning).
\end{itemize}

\textbf{Nachteile:}
\begin{itemize}[leftmargin=1.5em]
    \item \textbf{Speicher- und Rechenintensiv:} Alle Trainingsdaten müssen gespeichert werden; für jede neue Vorhersage müssen Distanzen zu allen Trainingsobjekten berechnet werden ($O(n \cdot p)$ pro Vorhersage).
    \item \textbf{Sensibel gegenüber Skalierung:} Variablen mit größerer Streuung dominieren die Distanzberechnung — \textbf{Normalisierung ist wesentlich}.
    \item \textbf{Ausreißer-Problem:} Ein einzelner Ausreißer nah bei einem Test-Punkt kann Vorhersage verderben.
    \item \textbf{Curse of Dimensionality:} In hochdimensionalen Räumen (großes p) werden Distanzen bedeutungslos; fast alle Punkte sind gleich weit entfernt. kNN funktioniert schlecht bei $p \gg n$.
    \item \textbf{Wahl von k unklar:} Keine offensichtliche Strategie, muss via CV getuned werden.
    \item \textbf{Unbalancierte Klassen:} Wenn eine Klasse viel häufiger ist, wird sie bevorzugt. Gewichtung der Nachbarn kann helfen.
\end{itemize}

\subsubsection{Varianten und Verbesserungen}

\definition{\textbf{Gewichtete kNN (Weighted kNN):} Statt gleiche Gewichte für alle $k$ Nachbarn, gewichte sie nach Distanz:
\formel{\hat{y} = \frac{\sum_{i \in \text{k-NN}} w_i y_i}{\sum_{i \in \text{k-NN}} w_i}}

Häufig: $w_i = 1/d(\mathbf{x}_i, \mathbf{x}_0)$ oder $w_i = \exp(-d(\mathbf{x}_i, \mathbf{x}_0)^2)$.

\textbf{Effekt:} Nähere Nachbarn haben mehr Einfluss — meist bessere Vorhersagen als einfaches kNN.}

\definition{\textbf{Lokale Adaptive Metriken:} Verwende unterschiedliche Metriken in verschiedenen Regionen des Merkmalsraums. Z.\,B.\ Mahalanobis-Distanz mit lokaler Kovarianzschätzung.}

\definition{\textbf{kNN mit Feature Selection:} Nicht alle Prädiktoren sind gleich wichtig. Mit Dimensionsreduktion (PCA, Filterung) oder \textbf{relevance weighting} können bessere Ergebnisse erzielt werden.}

\definition{\textbf{Approximate/Fast kNN:} Für große Datensätze: Space-Partitioning-Strukturen (KD-Trees, Ball-Trees) oder Locality-Sensitive Hashing (LSH) um die $k$ Nachbarn schneller zu finden, statt alle Distanzen zu berechnen.}

\textbf{Praktische Anwendungen:}
\begin{itemize}[leftmargin=1.5em]
    \item \textbf{Klassifizierung von Iris-Blüten, Handschriften, Bildern:} Wenn die Daten nicht zu hochdimensional sind.
    \item \textbf{Anomalieerkennung:} Objekte weit weg von allen anderen Objekten können Anomalien sein.
    \item \textbf{Recommender Systems:} ``Finde ähnliche Kunden/Produkte und empfehle, was diesen gefällt''.
    \item \textbf{Computer Vision:} Content-based Image Retrieval.
\end{itemize}

\textbf{Vergleich mit anderen Klassifikatoren:}
\begin{table}[h]
    \centering
    \small
    \begin{tabular}{|l|c|c|c|c|}
        \hline
        \textbf{Aspekt} & \textbf{kNN} & \textbf{LDA/QDA} & \textbf{Trees} & \textbf{SVM} \\
        \hline
        Einfachheit & Sehr einfach & Moderat & Einfach & Komplex \\
        \hline
        Komplexität & $O(n \cdot p)$ pro Test & Gering & Moderat & Hoch \\
        \hline
        Nonparametrisch & Ja & Nein & Ja & Nein \\
        \hline
        Interpretierbar & Nein & Ja & Ja & Nein \\
        \hline
        Curse of Dimensionality & Ja & Nein & Moderat & Moderat \\
        \hline
        Ideal für & Klein p, nonlinear & Normal-Daten & Mixed Types & Klassische ML \\
        \hline
    \end{tabular}
\end{table}






\subsection{Ähnlichkeitsmaße}
\newpage



\subsection{K-means, K-medoids und fuzzy-Clustering}

\definition{\textbf{K-means Clustering -- Grundkonzept:}}
K-means partitioniert einen Datensatz in $K$ disjunkte, nicht-leere Gruppen (Cluster), so dass Objekte innerhalb eines Clusters ähnlicher sind als Objekte zwischen Clustern. Jede Beobachtung gehört zu genau einem Cluster. Die Anzahl der Cluster $K$ muss im Voraus bekannt sein.

\textbf{Partitionsformalisierung:} Beobachtungen mit Indizes $1, \ldots, n$; Cluster $C_1, \ldots, C_K$ als Indexmengen:
\begin{itemize}[leftmargin=1.5em]
    \item Jede Beobachtung gehört zu mindestens einem Cluster: $C_1 \cup C_2 \cup \cdots \cup C_K = \{1, \ldots, n\}$
    \item Keine Beobachtung in mehreren Clustern: $C_i \cap C_j = \emptyset$ für $i \neq j$
\end{itemize}

\definition{\textbf{Optimierungskriterium:} Minimierung der Within-Cluster-Streuung (Within-Cluster-Sum-of-Squares, WCSS):} $$W = \sum_{k=1}^{K} \sum_{i \in C_k} \|\mathbf{x}_i - \mathbf{c}_k\|_2^2$$ wobei $\mathbf{c}_k = \frac{1}{n_k} \sum_{i \in C_k} \mathbf{x}_i$ das Zentroid der Clusters $k$ ist und $n_k = |C_k|$ die Clustergröße.

\definition{\textbf{K-means Algorithmus:}}
\begin{enumerate}[leftmargin=1.5em]
    \item \textbf{Initialisierung:} Wähle zufällig $K$ Cluster-Zentren (oder $K$ zufällige Beobachtungen).
    \item \textbf{Zuordnungsschritt:} Ordne jede Beobachtung $\mathbf{x}_i$ dem nächsten Zentroid zu: $i \in C_k$ wenn $\|\mathbf{x}_i - \mathbf{c}_k\| = \min_j \|\mathbf{x}_i - \mathbf{c}_j\|$.
    \item \textbf{Update-Schritt:} Berechne für jeden Cluster $k$ das neue Zentroid: $\mathbf{c}_k^{\text{neu}} = \frac{1}{n_k} \sum_{i \in C_k} \mathbf{x}_i$.
    \item \textbf{Konvergenzprüfung:} Wiederhole Schritte 2--3 bis keine Cluster-Änderungen mehr auftreten.
\end{enumerate}

\textbf{Wahl der Cluster-Anzahl $K$:}
\begin{itemize}[leftmargin=1.5em]
    \item \textbf{Elbow-Methode:} Plotte WCSS gegen $K$. Die Kurve ist monoton fallend; wähle $K$, bei dem die Kurve abknickt (``Ellbogen'').
    \item \textbf{Hopkins-Statistik:} Testet, ob der Datensatz überhaupt eine Clusterstruktur aufweist: $$H = \frac{\sum_{i=1}^{n^*} d_U(i)}{\sum_{i=1}^{n^*} d_U(i) + \sum_{i=1}^{n^*} d_W(i)}$$ wobei $d_U(i)$ die Distanz zufällig erzeugter Punkte zum nächsten echten Datenpunkt und $d_W(i)$ die Distanz echter Datenpunkte zum nächsten echten Nachbarn ist.
    \begin{itemize}[leftmargin=1.5em]     \item $H \approx 0{,}5$: Daten sind zufällig verteilt -- \textbf{keine
        Clusterstruktur vorhanden}; Clustering nicht sinnvoll.
        \item $H \to 1$: Starke Clusterstruktur -- Clustering sinnvoll.
        \item $H < 0{,}25$: Regelmäßige Gitterstruktur -- ebenfalls keine natürliche
        Clusterung.
    \end{itemize}
\end{itemize}

\definition{\textbf{K-medoids Clustering:}}
Ähnlich K-means, aber Cluster-Zentren müssen \emph{echte Beobachtungen} sein (nicht
der Durchschnitt). Im Update-Schritt wird für jeden Cluster $k$ die Medoid-Beobachtung
gesucht:
$$\mathbf{m}_k = \arg\min_{\mathbf{x}_i \in C_k} \sum_{\mathbf{x}_j \in C_k}
\|\mathbf{x}_j - \mathbf{x}_i\|^2$$ 
\textbf{Vorteile:} Robuster gegenüber Ausreißern; kann beliebige Distanzmaße verwenden
(nicht nur Euklidisch); Zentren sind interpretierbar (echte Objekte).

\definition{\textbf{Fuzzy Clustering (Soft Clustering):}}
Im Gegensatz zu K-means (hard clustering) kann jede Beobachtung mit Zugehörigkeits\-graden
zu mehreren Clustern gehören.

\textbf{Membership-Koeffizient:} $u_{ij} \in [0, 1]$ gibt an, mit welchem Grad
Beobachtung $i$ zu Cluster $j$ gehört, mit Normierungsbedingung:
$$\sum_{j=1}^{K} u_{ij} = 1 \quad \text{für alle } i$$

\definition{\textbf{Fuzzy C-Means (FCM) Algorithmus:}} Minimiert die gewichtete Within-Cluster-Streuung: $$J_m = \sum_{i=1}^{n} \sum_{j=1}^{K} u_{ij}^m \|\mathbf{x}_i - \mathbf{c}_j\|^2$$ wobei $m > 1$ der \textbf{Fuzzifikationsparameter} ist (typisch $m = 2$). Größeres $m$ führt zu weicheren (unschärferen) Zugehörigkeiten.

\textbf{Update-Formeln (alternierend bis Konvergenz):}

\textbf{Membership-Update:} $$u_{ij} = \frac{1}{\displaystyle\sum_{l=1}^{K} \left(\frac{\|\mathbf{x}_i - \mathbf{c}_j\|} {\|\mathbf{x}_i - \mathbf{c}_l\|}\right)^{\!\frac{2}{m-1}}}$$

\textit{Interpretation:} Je näher Beobachtung $i$ am Zentroid $j$ liegt relativ zu allen anderen Zentroiden, desto höher ist $u_{ij}$.

\textbf{Zentroid-Update:} $$\mathbf{c}_j = \frac{\sum_{i=1}^{n} u_{ij}^m\, \mathbf{x}_i} {\sum_{i=1}^{n} u_{ij}^m}$$ \textit{Interpretation:} Jedes Zentroid ist ein gewichtetes Mittel aller Beobachtungen, wobei die Gewichte durch die Zugehörigkeitsgrade (hoch potenziert) bestimmt werden. \textbf{Finale Klassenzuweisung (optional):} Nach Konvergenz kann aus den Membership-Koeffizienten eine harte Partition abgeleitet werden: $i \in C_k$ wenn $u_{ik} = \max_j u_{ij}$.

\textbf{Vorteile gegenüber K-means:} Bildet Unsicherheit bei Zuordnungen ab; nützlich bei Grenzfällen zwischen Clustern; liefert vollständige Zugehörigkeitsverteilung pro Beobachtung.

\textbf{Vergleich der drei Verfahren:}
\begin{table}[h]
    \centering
    \small
    \begin{tabular}{|l|l|l|l|}
        \hline
        \textbf{Aspekt} & \textbf{K-means} & \textbf{K-medoids} &
        \textbf{Fuzzy C-Means} \\
        \hline
        Zuordnung & Hart (exklusiv) & Hart (exklusiv) & Weich (graduell) \\
        \hline
        Zentren & Durchschnitt & Echte Beobachtung & Durchschnitt (gewichtet) \\
        \hline
        Ausreißer-Robustheit & Gering & Hoch & Mittel \\
        \hline
        Distanzmaß & Euklidisch & Beliebig & Euklidisch (standard) \\
        \hline
        Zusatz-Parameter & -- & -- & Fuzzifikator $m$ \\
        \hline
        Interpretierbarkeit & Hoch & Sehr hoch & Mittel \\
        \hline
    \end{tabular}
\end{table}

\subsection{Model-based Clustering}

\definition{\textbf{Grundkonzept:}}
Im Gegensatz zu K-means/K-medoids/Fuzzy Clustering (modellfreie Verfahren) wird bei model-based Clustering angenommen, dass die Daten aus einer Mischung von $K \geq 2$ Wahrscheinlichkeitsverteilungen stammen (üblicherweise multivariate Normalverteilungen).

\textbf{Annahmen:} Jeder Cluster folgt einer eigenen Wahrscheinlichkeitsverteilung mit eigenem Mittelwertvektor und Kovarianzmatrix.

\definition{\textbf{Mixture Model (Mischmodell):}}
Die Dichtefunktion der Gesamtpopulation wird als gewichtete Summe von $K$ Komponentendichten modelliert:
$$f(\mathbf{x}) = \sum_{k=1}^{K} \pi_k f_k(\mathbf{x} \mid \boldsymbol{\theta}_k)$$
wobei:
\begin{itemize}[leftmargin=1.5em]
    \item $\pi_k$ das Gewicht (die Wahrscheinlichkeit) der $k$-ten Komponente mit $\sum_{k=1}^{K} \pi_k = 1$
    \item $f_k(\mathbf{x} \mid \boldsymbol{\theta}_k)$ die Dichte der $k$-ten Komponente mit Parametern $\boldsymbol{\theta}_k$ (z.B. Mittelwert $\boldsymbol{\mu}_k$ und Kovarianzmatrix $\boldsymbol{\Sigma}_k$)
\end{itemize}

\definition{\textbf{Univariate Normalverteilung $N(\mu, \sigma^2)$:}}
Dichtefunktion:
$$f(x) = \frac{1}{\sigma\sqrt{2\pi}} \exp\left(-\frac{(x-\mu)^2}{2\sigma^2}\right)$$
parametrisiert durch Mittelwert $\mu$ und Varianz $\sigma^2$. Häufigste Annahme in statistischen Modellen.

\definition{\textbf{Multivariate Normalverteilung $N(\boldsymbol{\mu}, \boldsymbol{\Sigma})$:}}
Dichtefunktion:
$$f(\mathbf{x}) = \frac{1}{(2\pi)^{p/2} |\boldsymbol{\Sigma}|^{1/2}} \exp\left(-\frac{1}{2}(\mathbf{x}-\boldsymbol{\mu})^T \boldsymbol{\Sigma}^{-1} (\mathbf{x}-\boldsymbol{\mu})\right)$$
parametrisiert durch Mittelwertvektor $\boldsymbol{\mu} \in \mathbb{R}^p$ und $p \times p$ Kovarianzmatrix $\boldsymbol{\Sigma}$.

\definition{\textbf{Expectation-Maximization (EM) Algorithmus:}}
Iteratives Verfahren zur Schätzung von Mischmodell-Parametern:
\begin{enumerate}[leftmargin=1.5em]
    \item \textbf{Initialisierung:} Starte mit initialen Parametern (z.B. aus K-means).
    \item \textbf{E-Schritt (Expectation):} Berechne für jede Beobachtung $\mathbf{x}_i$ die posteriore Wahrscheinlichkeit, dass sie zu Cluster $k$ gehört (ähnlich fuzzy membership):
    $$\tau_{ik} = \frac{\pi_k f_k(\mathbf{x}_i \mid \boldsymbol{\theta}_k)}{\sum_{j=1}^{K} \pi_j f_j(\mathbf{x}_i \mid \boldsymbol{\theta}_j)}$$
    \item \textbf{M-Schritt (Maximization):} Schätze Parameter neu basierend auf $\tau_{ik}$:
    \begin{itemize}
        \item Gewichte: $\pi_k = \frac{1}{n}\sum_{i=1}^{n} \tau_{ik}$
        \item Mittelwert: $\boldsymbol{\mu}_k = \frac{\sum_{i=1}^{n} \tau_{ik} \mathbf{x}_i}{\sum_{i=1}^{n} \tau_{ik}}$
        \item Kovarianzmatrix: $\boldsymbol{\Sigma}_k = \frac{\sum_{i=1}^{n} \tau_{ik} (\mathbf{x}_i - \boldsymbol{\mu}_k)(\mathbf{x}_i - \boldsymbol{\mu}_k)^T}{\sum_{i=1}^{n} \tau_{ik}}$
    \end{itemize}
    \item \textbf{Konvergenzprüfung:} Wiederholen bis Log-Likelihood konvergiert.
\end{enumerate}

\textbf{Vorteile:} Probabilistisches Framework; liefert Wahrscheinlichkeiten (nicht nur harte Zuordnungen); erlaubt statistische Inferenz; Modellwahl durch Likelihood-Kriterien (AIC, BIC) möglich.

\textbf{Nachteile:} Computational aufwendiger als K-means; Annahme von Normalverteilung nicht immer realistisch; für hochdimensionale Daten kann Kovarianzmatrix instabil sein.

\newpage

\subsection{Hierarchisches Clustering}

\definition{\textbf{Grundkonzept:}}
Hierarchische Methoden erzeugen eine Hierarchie von Partitionen (Clustersolutionen), statt einer einzelnen festen Partition wie K-means. Zwei grundlegende Ansätze:
\begin{itemize}[leftmargin=1.5em]
    \item \textbf{Agglomerativ:} Starten mit $n$ Clustern (jede Beobachtung einzeln), Zusammenfügen der nächsten Cluster iterativ bis ein einzelner Cluster übrig bleibt.
    \item \textbf{Divisiv:} Starten mit einem Cluster (alle Beobachtungen), Aufteilung iterativ bis $n$ einzelne Cluster. (Seltener, rechenintensiver.)
\end{itemize}

\definition{\textbf{Abstände zwischen Clustern (Linkage-Methoden):}}
Um zu entscheiden, welche Cluster zu fusionieren sind, muss der Abstand zwischen zwei Clustern definiert werden:

\begin{enumerate}[leftmargin=1.5em]
    \item \textbf{Complete Linkage (Furthest Neighbor):} Abstand zwischen den zwei am weitesten entfernten Punkten:
    $$d(C_i, C_j) = \max_{\mathbf{x}_a \in C_i, \mathbf{x}_b \in C_j} d(\mathbf{x}_a, \mathbf{x}_b)$$
    Vorteil: Vermeidet Kettenerzeugnis, produziert globularere Cluster.
    
    \item \textbf{Single Linkage (Nearest Neighbor):} Abstand zwischen den zwei nächsten Punkten:
    $$d(C_i, C_j) = \min_{\mathbf{x}_a \in C_i, \mathbf{x}_b \in C_j} d(\mathbf{x}_a, \mathbf{x}_b)$$
    Vorteil: Einfach; Nachteil: Anfällig für Kettenerzeugnis (``chaining'').
    
    \item \textbf{Average Linkage (UPGMA):} Durchschnittlicher Abstand zwischen allen Punkt-Paaren:
    $$d(C_i, C_j) = \frac{1}{n_i n_j} \sum_{\mathbf{x}_a \in C_i} \sum_{\mathbf{x}_b \in C_j} d(\mathbf{x}_a, \mathbf{x}_b)$$
    Kompromiss zwischen Single und Complete Linkage.
    
    \item \textbf{Centroid Method:} Abstand zwischen Cluster-Zentroiden:
    $$d(C_i, C_j) = \|\mathbf{c}_i - \mathbf{c}_j\|$$
    wobei $\mathbf{c}_k = \frac{1}{n_k} \sum_{\mathbf{x} \in C_k} \mathbf{x}$
    
    \item \textbf{Ward's Method:} Minimiert die Within-Cluster-Varianz bei Fusion:
    $$d(C_i, C_j) = \frac{n_i n_j}{n_i + n_j} \|\mathbf{c}_i - \mathbf{c}_j\|^2$$
    Ähnlich zu K-means Optimierungskriterium; oft beste praktische Ergebnisse.
\end{enumerate}

\definition{\textbf{Dendrogramm (Dendrogram):}}
Visualisierung der hierarchischen Struktur. Ein Baum-Diagramm mit:
\begin{itemize}[leftmargin=1.5em]
    \item Leaves (Blätter): Einzelne Beobachtungen
    \item Branches (Äste): Zusammengefügte Cluster
    \item Höhe der Verbindung: Distanz zwischen den Clustern bei Fusion
\end{itemize}

Durch horizontale Schnitte durch das Dendrogramm erhält man verschiedene Cluster-Lösungen für unterschiedliche Anzahlen $K$ von Clustern.

\definition{\textbf{Agglomerationskoeffizient (ac):}}
Maß für die Clusterungsfähigkeit eines Datensatzes. Für jede Beobachtung $i$ berechne:
$$m(i) = \frac{d_{i,\text{first}}}{d_{i,\text{last}}}$$
wobei $d_{i,\text{first}}$ die Distanz beim ersten Zusammenfügen von $i$ und $d_{i,\text{last}}$ beim letzten (Gesamt-)Zusammenfügen ist. Der Agglomerationskoeffizient ist das Mittel:
$$ac = \frac{1}{n} \sum_{i=1}^{n} (1 - m(i))$$

Werte nahe 1 deuten auf starke Clusterstruktur hin.

\textbf{Vor- und Nachteile:}
\begin{itemize}[leftmargin=1.5em]
    \item \textbf{Vorteil:} Flexibel; liefert Dendrogramm zur visuellen Exploration; keine Vorfestlegung von $K$ notwendig; verschiedene Linkage-Methoden ermöglichen unterschiedliche Cluster-Formen.
    \item \textbf{Nachteil:} Computationally teuer für große Datenmengen ($O(n^2)$ oder schlimmer); Fusionsentscheidungen sind final (keine Neuoptimierung); kann durch lokale Minima beeinflusst werden.
\end{itemize}

\newpage

\subsection{Bedingte Wahrscheinlichkeit und Bayes-Theorem}

\definition{\textbf{Bedingte Wahrscheinlichkeit (Konditionale Wahrscheinlichkeit):}}
Die Wahrscheinlichkeit des Eintretens eines Ereignisses $A$ unter der Bedingung, dass ein anderes Ereignis $B$ bereits eingetreten ist:
$$P(A \mid B) = \frac{P(A \cap B)}{P(B)}, \quad P(B) > 0$$

\textbf{Interpretation:} Wenn $B$ eingetreten ist, reduzieren sich die möglichen Ausgänge auf die Ergebnisse in $B$. Die bedingte Wahrscheinlichkeit ist eine Neueinschätzung von $P(A)$ unter dieser Information.

\definition{\textbf{Produktregel (Kettenregel):}}
Aus der Definition folgt unmittelbar:
$$P(A \cap B) = P(A \mid B) \cdot P(B) = P(B \mid A) \cdot P(A)$$

Für mehr als zwei Ereignisse:
$$P(A \cap B \cap C) = P(A) \cdot P(B \mid A) \cdot P(C \mid A \cap B)$$

\textbf{Anwendung:} Multivariate Verteilungen können als Produkt von bedingten Wahrscheinlichkeiten zerlegt werden. Die Faktorisierung hängt von der gewählten Reihenfolge ab und ermöglicht es, Domänenwissen über kausale Abhängigkeiten einzubeziehen.

\definition{\textbf{Satz der Totalen Wahrscheinlichkeit (Law of Total Probability):}}
Wenn $B_1, B_2, \ldots, B_K$ eine Partition des Stichprobenraums $\Omega$ bilden (d.h. $B_i \cap B_j = \emptyset$ und $\bigcup_{k} B_k = \Omega$), dann:
$$P(A) = \sum_{k=1}^{K} P(A \mid B_k) \cdot P(B_k)$$

\textbf{Interpretation:} Die Gesamtwahrscheinlichkeit von $A$ wird als gewichtete Summe über alle möglichen Bedingungen $B_k$ dargestellt.

\definition{\textbf{Satz von Bayes (Bayes' Theorem):}}
$$P(A \mid B) = \frac{P(B \mid A) \cdot P(A)}{P(B)}$$

wobei:
\begin{itemize}[leftmargin=1.5em]
    \item $P(A \mid B)$ die \textbf{Posteriori-Wahrscheinlichkeit} (posterior probability) — die aktualisierte Wahrscheinlichkeit nach Beobachtung von $B$
    \item $P(B \mid A)$ die \textbf{Likelihood} — die Wahrscheinlichkeit der Beobachtung $B$, wenn $A$ wahr ist
    \item $P(A)$ die \textbf{Prior-Wahrscheinlichkeit} (a-priori) — das Vorwissen über $A$ vor Beobachtung
    \item $P(B)$ die \textbf{Evidenz (Marginal Likelihood)} — die Gesamtwahrscheinlichkeit der Beobachtung
\end{itemize}

\textbf{Anwendungsform mit Normalisierungskonstante:}
$$P(A \mid B) = \frac{P(B \mid A) \cdot P(A)}{\sum_{k} P(B \mid A_k) \cdot P(A_k)}$$

\textbf{Interpretation:} Der Satz von Bayes ermöglicht es, Schlussfolgerungen "rückwärts" zu ziehen. Man geht von bekanntem $P(B \mid A)$ aus und berechnet die umgekehrte bedingte Wahrscheinlichkeit $P(A \mid B)$.

\beispiel{\textbf{Medizinischer Test:} Eine Krankheit tritt in 1\% der Population auf. Ein Test ist 95\% sensitiv (d.h. $P(+\mid \text{Krankheit}) = 0.95$) und 90\% spezifisch (d.h. $P(-\mid \text{keine Krankheit}) = 0.90$).

Frage: Wie groß ist $P(\text{Krankheit} \mid +)$?
\begin{align*}
P(\text{K} \mid +) &= \frac{P(+ \mid \text{K}) \cdot P(\text{K})}{P(+)} \\
P(+) &= P(+ \mid \text{K}) \cdot P(\text{K}) + P(+ \mid \text{kK}) \cdot P(\text{kK}) \\
&= 0.95 \cdot 0.01 + 0.10 \cdot 0.99 = 0.0095 + 0.099 = 0.1085 \\
P(\text{K} \mid +) &= \frac{0.95 \cdot 0.01}{0.1085} \approx 0.0876 \approx 8.76\%
\end{align*}

Obwohl der Test gut erscheint, ist die Wahrscheinlichkeit einer tatsächlichen Krankheit bei positivem Test weniger als 9\%! Dies ist die Base-Rate Fallacy.}

\newpage

\subsection{Probabilistische Graphische Modelle}

\definition{\textbf{Graphische Modelle (Graphical Models):}}
Repräsentation von Wahrscheinlichkeitsverteilungen mittels gerichteter oder ungerichteter Graphen, wobei:
\begin{itemize}[leftmargin=1.5em]
    \item \textbf{Knoten (Nodes):} Zufallsvariablen oder Parameter
    \item \textbf{Kanten (Edges):} Abhängigkeitsbeziehungen zwischen Variablen
\end{itemize}

\definition{\textbf{Bayes-Netz (Bayesian Network):}}
Ein gerichteter, azyklischer Graph (DAG: Directed Acyclic Graph), bei dem eine Kante $A \to B$ ausdrückt, dass $A$ ein direkter "Eltern"-Knoten von $B$ ist.

\textbf{Faktorisierung:} Eine multivariate Verteilung faktorisiert über die Graphstruktur:
$$P(X_1, \ldots, X_n) = \prod_{i=1}^{n} P(X_i \mid \text{Eltern}(X_i))$$

Dabei ist $\text{Eltern}(X_i)$ die Menge der direkten Vorgänger von $X_i$ im Graphen.

\beispiel{\textbf{Bedingte Unabhängigkeit:} In einem Bayes-Netz $X \to Y \to Z$:
\begin{itemize}[leftmargin=1.5em]
    \item $P(X, Y, Z) = P(X) \cdot P(Y \mid X) \cdot P(Z \mid Y)$
    \item Bedingte Unabhängigkeit: $X \perp Z \mid Y$ (X und Z sind unabhängig, gegeben Y)
    \item Dies reduziert die Parametrisierung erheblich!
\end{itemize}
}

\definition{\textbf{Bedingte Unabhängigkeit und d-Separation:}}
In einem Bayes-Netz sind zwei Variablen $X$ und $Z$ bedingt unabhängig gegeben $Y$ ($X \perp Z \mid Y$), wenn:
\begin{enumerate}[leftmargin=1.5em]
    \item \textbf{Seriell:} $X \to Y \to Z$: Information fließt von $X$ zu $Z$ über $Y$. Wenn $Y$ beobachtet wird, wird der Fluss blockiert.
    \item \textbf{Divergent:} $X \leftarrow Y \to Z$: Ein gemeinsamer Eltern-Knoten $Y$. Wenn $Y$ beobachtet wird, werden $X$ und $Z$ unabhängig.
    \item \textbf{Konvergent:} $X \to Y \leftarrow Z$: Ein gemeinsamer Kind-Knoten $Y$. Wenn $Y$ nicht beobachtet wird, sind $X$ und $Z$ unabhängig. Wenn $Y$ oder einer seiner Nachfahren beobachtet wird, werden $X$ und $Z$ abhängig ("explaining away").
\end{enumerate}

\definition{\textbf{Naive-Bayes-Klassifikator:}}
Ein einfaches probabilistisches Klassifikationsmodell, basierend auf bedingter Unabhängigkeit aller Features gegeben die Zielvariable.

\textbf{Annahme:} $P(X_1, \ldots, X_p \mid C) = \prod_{j=1}^{p} P(X_j \mid C)$, wobei $C$ die Klasse und $X_1, \ldots, X_p$ die Features sind.

\textbf{Klassifikationsregel (Maximierung der Posteriori):}
$$\hat{C} = \arg\max_c P(C = c \mid X_1, \ldots, X_p) = \arg\max_c P(C = c) \prod_{j=1}^{p} P(X_j \mid C = c)$$

\textbf{Praktische Implementierung:} Meist wird im Log-Raum gerechnet, um Underflow zu vermeiden:
$$\hat{C} = \arg\max_c \left[ \log P(C = c) + \sum_{j=1}^{p} \log P(X_j \mid C = c) \right]$$

\definition{\textbf{Laplace-Smoothing (Add-One Smoothing):}}
Zur Vermeidung von Null-Wahrscheinlichkeiten (wenn eine Feature-Klasse-Kombination in den Trainingsdaten nicht vorkommt):
$$P(X_j = x \mid C = c) = \frac{\text{count}(X_j = x, C = c) + \alpha}{\text{count}(C = c) + \alpha \cdot m}$$

wobei $\alpha$ der Smoothing-Parameter (typisch $\alpha = 1$) und $m$ die Anzahl möglicher Werte von $X_j$ ist.

\textbf{Effekt:} Pseudo-Zählungen werden hinzugefügt; dies sorgt dafür, dass selbst unsichtbare Kombinationen in Trainingsdaten positive Wahrscheinlichkeiten erhalten und verhindert das "Zero-Problem".

\textbf{Vorteile des Naive-Bayes-Klassifikators:}
\begin{itemize}[leftmargin=1.5em]
    \item Einfach und schnell zu trainieren und zu verwenden
    \item Funktioniert überraschend gut auch bei moderaten Verletzungen der Unabhängigkeitsannahme
    \item Benötigt relativ wenig Trainingsdaten
    \item Natürlich probabilistisch; liefert Klassifikations-Sicherheiten
\end{itemize}

\textbf{Nachteile:}
\begin{itemize}[leftmargin=1.5em]
    \item Naive Unabhängigkeitsannahme ist oft unrealistisch
    \item Kann zu schlechten Kalibrierungen führen (Wahrscheinlichkeiten nicht vertrauenswürdig)
    \item Feature-Engineering und Feature-Interaktionen werden ignoriert
    \item Nicht für hochgradig korrelierte Features geeignet
\end{itemize}

\newpage

\subsection{Bayes-Updates: Prior, Likelihood, Evidenz, Posterior}

\definition{\textbf{Bayesianische Inferenz -- Grundkonzept:}}
Im Gegensatz zur Frequentist'schen Statistik werden unbekannte Größen (Parameter, latente Variablen) als Zufallsvariablen behandelt und durch Wahrscheinlichkeitsverteilungen beschrieben. Der Satz von Bayes als Update-Regel bildet das Herzstück der Bayesianischen Statistik.

\definition{\textbf{Prior-Verteilung $p(\theta)$:}}
Die a-priori-Wahrscheinlichkeitsverteilung, die Vorkenntnisse über den Parameter $\theta$ vor Beobachtung von Daten ausdrückt. Quellen:
\begin{itemize}[leftmargin=1.5em]
    \item Domänenwissen und Expert-Meinung
    \item Literaturwerte
    \item Frühere ähnliche Studien
    \item "Uninformative" Priors (z.B. gleichmäßig) bei fehlendem Vorwissen
\end{itemize}

\textbf{Kritik:} Die Wahl des Priors ist subjektiv, kann aber durch Sensitivitätsanalyse überprüft werden.

\definition{\textbf{Likelihood-Funktion $p(\mathcal{D} \mid \theta)$ oder $L(\theta)$:}}
Die Wahrscheinlichkeit der beobachteten Daten $\mathcal{D}$ gegeben einen spezifischen Wert des Parameters $\theta$. Für unabhängige Beobachtungen:
$$L(\theta) = p(\mathcal{D} \mid \theta) = \prod_{i=1}^{n} p(\mathbf{x}_i \mid \theta)$$

\textbf{Interpretation:} Die Likelihood misst, wie gut die beobachteten Daten durch den Parameter $\theta$ erklärt werden.

\definition{\textbf{Evidence (Marginal Likelihood) $p(\mathcal{D})$:}}
Die Gesamtwahrscheinlichkeit der Daten unter allen möglichen Parameterwerten, marginalisiert über den Parameter:
$$p(\mathcal{D}) = \int p(\mathcal{D} \mid \theta) p(\theta) \, d\theta$$

\textbf{Rolle:} Normalisierungskonstante; wichtig für Modellvergleich (Bayes Factor) und für die Berechnung der korrekten Posteriori-Wahrscheinlichkeiten.

\definition{\textbf{Posterior-Verteilung $p(\theta \mid \mathcal{D})$:}}
Die aktualisierte Wahrscheinlichkeitsverteilung des Parameters nach Beobachtung der Daten. Der Satz von Bayes liefert:
$$p(\theta \mid \mathcal{D}) = \frac{p(\mathcal{D} \mid \theta) \cdot p(\theta)}{p(\mathcal{D})} = \frac{\text{Likelihood} \cdot \text{Prior}}{\text{Evidence}}$$

\textbf{Merkregel:}
$$\boxed{\text{Posterior} = \frac{\text{Likelihood} \times \text{Prior}}{\text{Evidence}}}$$

\textbf{Interpretation:} Der Posterior kombiniert Vorwissen (Prior) mit neuer Evidenz (Likelihood), um eine Posteriori-Verteilung zu erhalten, die beide Informationen balanciert.

\definition{\textbf{Sequential Bayesian Update:}}
Ein großer Vorteil des Bayesianischen Ansatzes ist die Möglichkeit, mit neuen Daten zu aktualisieren:
$$p(\theta \mid \mathcal{D}_1, \mathcal{D}_2) = \frac{p(\mathcal{D}_2 \mid \theta) \cdot p(\theta \mid \mathcal{D}_1)}{p(\mathcal{D}_2 \mid \mathcal{D}_1)}$$

Der \textbf{alte Posterior wird zum neuen Prior}. Dies erlaubt kontinuierliche Lernprozesse und ist fundamental für Online-Learning und sequentielle Entscheidungsfindung.

\beispiel{\textbf{Münzwurf-Beispiel:} Schätze die Erfolgswahrscheinlichkeit $p$ einer Münze.

\textbf{Prior:} Uninformativ, z.B. $p \sim \text{Beta}(1, 1)$ (gleichmäßig auf $[0, 1]$)

\textbf{Beobachtungen:} 5 Würfe, 3 Erfolge

\textbf{Likelihood:} $L(p) = p^3 (1-p)^2$

\textbf{Posterior (analytisch):} $p(p \mid \text{data}) \propto p^3 (1-p)^2 \cdot 1 \sim \text{Beta}(4, 3)$

\textbf{Posterior-Erwartungswert:} $\mathbb{E}[p \mid \text{data}] = \frac{4}{4+3} = \frac{4}{7} \approx 0.571$

Im Vergleich: Frequentist'sche MLE würde $\hat{p} = 3/5 = 0.6$ geben. Der Bayes-Schätzer ist durch den uninformativen Prior leicht zu diesem Wert gezogen.}

\definition{\textbf{Posteriorpunktschätzung:}}
Unterschiedliche Schätzverfahren je nach Zielkriterium:
\begin{itemize}[leftmargin=1.5em]
    \item \textbf{Maximum A Posteriori (MAP):} $\hat{\theta}_{\text{MAP}} = \arg\max_\theta p(\theta \mid \mathcal{D})$ — der Modus des Posteriors
    \item \textbf{Posterior Mean:} $\hat{\theta}_{\text{Mean}} = \mathbb{E}[\theta \mid \mathcal{D}] = \int \theta \cdot p(\theta \mid \mathcal{D}) \, d\theta$ — der Erwartungswert
    \item \textbf{Posterior Median:} $\hat{\theta}_{\text{Median}}$ — der Median des Posteriors
\end{itemize}

\definition{\textbf{Posteriorwahrscheinlichkeitsintervalle (Credible Intervals):}}
Im Gegensatz zu Frequentist'schen Konfidenzintervallen können Bayesianische Credible Intervals direkt probabilistisch interpretiert werden.

\textbf{Gleichschwanz-Credible Interval (ETI):} Das zentrale $100(1-\alpha)\%$-Intervall mit gleichen Schwänzen.

\textbf{Höchste Posteriordichte Intervall (HDI):} Der zusammenhängende Bereich der höchsten Posteriordichte, der $100(1-\alpha)\%$ der Wahrscheinlichkeit enthält.

\textbf{Interpretation:} Mit $100(1-\alpha)\%$ Kredibilität liegt der wahre Parameterwert im Credible Interval — eine intuitivere Aussage als das Frequentist'sche Konfidenzintervall.

\newpage

\subsection{Konjugierte Priors}

\definition{\textbf{Konjugierter Prior:}}
Ein Prior $p(\theta)$ ist konjugiert zur Likelihood $p(\mathcal{D} \mid \theta)$, wenn der resultierende Posterior $p(\theta \mid \mathcal{D})$ die gleiche funktionale Form wie der Prior hat (aus der gleichen Verteilungsfamilie).

\textbf{Vorteil:} Erlaubt analytische Lösungen; keine numerische Integration notwendig; Posterior hat oft geschlossene Form mit einfachen Update-Regeln.

\definition{\textbf{Beta-Binomial Konjugation:}}
Likelihood: Bernoulli oder Binomial mit Erfolgswahrscheinlichkeit $p$

Prior: $\theta \sim \text{Beta}(\alpha, \beta)$ mit Dichtefunktion $p(\theta) \propto \theta^{\alpha-1} (1-\theta)^{\beta-1}$

Posterior: Nach $k$ Erfolgen in $n$ Versuchen:
$$p(\theta \mid \text{data}) \sim \text{Beta}(\alpha + k, \beta + n - k)$$

\textbf{Update-Regel:} $\alpha \to \alpha + k$ (Erfolge dazuaddieren), $\beta \to \beta + (n-k)$ (Misserfolge dazuaddieren).

\textbf{Spezialfall:} Uninformativ Prior $\text{Beta}(1, 1)$ (uniform); mit Laplace Smoothing: $\text{Beta}(\alpha_0, \beta_0)$ mit $\alpha_0, \beta_0 > 1$.

\definition{\textbf{Normal-Normal Konjugation (mit bekannter Varianz):}}
Likelihood: Normalverteilung $N(\theta, \sigma^2)$ mit bekannter Varianz $\sigma^2$

Prior: $\theta \sim N(\mu_0, \tau_0^2)$

Posterior: $\theta \mid \mathcal{D} \sim N(\mu_n, \tau_n^2)$

\textbf{Update-Formeln:}
\begin{align*}
\tau_n^{-2} &= \tau_0^{-2} + n\sigma^{-2} \quad \text{(Genauigkeiten addieren sich)} \\
\mu_n &= \tau_n^2 \left( \tau_0^{-2} \mu_0 + n\sigma^{-2} \bar{x} \right)
\end{align*}

\definition{\textbf{Dirichlet-Multinomial Konjugation:}}
Likelihood: Multinomial (Verallgemeinerung der Binomial) mit Wahrscheinlichkeitsvektor $\boldsymbol{\theta} = (\theta_1, \ldots, \theta_k)$

Prior: $\boldsymbol{\theta} \sim \text{Dir}(\alpha_1, \ldots, \alpha_k)$

Posterior: Nach Beobachtung von $n_1, \ldots, n_k$ Counts:
$$\boldsymbol{\theta} \mid \text{data} \sim \text{Dir}(\alpha_1 + n_1, \ldots, \alpha_k + n_k)$$

\textbf{Anwendung:} Klassifikation mit mehreren Klassen; Wort-Häufigkeiten in Textanalyse (Naive Bayes).

\textbf{Übersicht Konjugierter Priors:}
\begin{table}[h]
    \centering
    \small
    \begin{tabular}{|l|l|l|}
        \hline
        \textbf{Likelihood} & \textbf{Konjugierter Prior} & \textbf{Posterior-Familie} \\
        \hline
        Bernoulli ($p$) & Beta & Beta \\
        \hline
        Binomial ($p$) & Beta & Beta \\
        \hline
        Multinomial ($\boldsymbol{p}$) & Dirichlet & Dirichlet \\
        \hline
        Normal ($\mu$, $\sigma^2$ bekannt) & Normal & Normal \\
        \hline
        Normal ($\sigma^2$, $\mu$ bekannt) & Inverse-Gamma & Inverse-Gamma \\
        \hline
        Poisson ($\lambda$) & Gamma & Gamma \\
        \hline
        Exponential ($\lambda$) & Gamma & Gamma \\
        \hline
    \end{tabular}
\end{table}

\definition{\textbf{Detailliertes Beispiel: Normal-Normal Updates (Gauß-Gauß Modell)}}

Betrachte eine eindimensionale Normalverteilung mit bekannter Varianz $\sigma^2 = 1$ und unbekanntem Lageparameter $\mu$:
$$p(x \mid \mu) = \frac{1}{\sqrt{2\pi}} e^{-\frac{1}{2}(x-\mu)^2}$$

Annahme: Der Prior ist eine Normalverteilung mit Mittelwert $\mu_0 = 0$ und Varianz $\tilde{\sigma}_0^2 = 1$:
$$p(\mu) = \frac{1}{\sqrt{2\pi}} e^{-\frac{1}{2}\mu^2}$$

\textbf{Interpretation des Priors als Pseudo-Datenpunkt:} Der Prior $N(\mu_0, \tilde{\sigma}_0^2)$ wirkt wie ein einzelner fiktiver Datenpunkt $x_0 = \mu_0$ mit Varianz $\sigma^2 = 1$.

\textbf{Erster Datenpunkt $x_1 = 1$:} Mit Bayes' Theorem:
$$p(\mu \mid D_1) \propto p(x_1 \mid \mu) \cdot p(\mu) = e^{-\frac{1}{2}(1-\mu)^2} \cdot e^{-\frac{1}{2}\mu^2}$$

Im Exponenten:
$$-\frac{1}{2}(1-\mu)^2 - \frac{1}{2}\mu^2 = -\frac{1}{2}(1 - 2\mu + \mu^2 + \mu^2) = -(\mu^2 - \mu + \frac{1}{2})$$

Quadratisches Ergänzen: $\mu^2 - \mu = (\mu - \frac{1}{2})^2 - \frac{1}{4}$

$$\Rightarrow p(\mu \mid D_1) \propto e^{-(\mu - \frac{1}{2})^2 + \frac{1}{4}} = e^{-(\mu - \frac{1}{2})^2} \cdot \text{const}$$

Also: $p(\mu \mid D_1) = N(\mu_1, \tilde{\sigma}_1^2)$ mit $\mu_1 = \frac{1}{2}$ und $\tilde{\sigma}_1^2 = \frac{1}{2}$

\textbf{Interpretation:} Der neue Mittelwert liegt zwischen dem Prior-Mittelwert ($\mu_0 = 0$) und dem MLE ($\mu_{\text{MLE}} = 1$). Die Varianz hat sich reduziert: Zuversicht in der Schätzung nimmt zu.

\textbf{Zweiter Datenpunkt $x_2 = -\frac{1}{2}$:} Der alte Posterior wird zum neuen Prior:
$$p(\mu \mid D_{1,2}) \propto p(x_2 \mid \mu) \cdot p(\mu \mid D_1)$$

Nach analoger Berechnung:
$$p(\mu \mid D_{1,2}) = N(\mu_2, \tilde{\sigma}_2^2) \quad \text{mit} \quad \mu_2 = \frac{1}{6}, \quad \tilde{\sigma}_2^2 = \frac{1}{3}$$

\textbf{Allgemeine Formel für $n$ Datenpunkte:}

Mit $n$ Datenpunkten $D_{1,\ldots,n} = \{x_1, \ldots, x_n\}$ und Prior $N(\mu_0, \tilde{\sigma}_0^2)$:

$$p(\mu \mid D_{1,\ldots,n}) = N(\mu_n, \tilde{\sigma}_n^2)$$

mit:
$$\tilde{\sigma}_n^2 = \frac{1}{\frac{1}{\tilde{\sigma}_0^2} + \frac{n}{\sigma^2}} \quad \text{und} \quad \mu_n = \tilde{\sigma}_n^2 \left( \frac{\mu_0}{\tilde{\sigma}_0^2} + \frac{1}{\sigma^2}\sum_{k=1}^{n} x_k \right)$$

\textbf{Alternative Darstellung (Präzisions-Gewichte):} Mit Präzisionen $\tau = \sigma^{-2}$:

$$\mu_n = \frac{\frac{1}{\tilde{\sigma}_0^2} \mu_0 + n \cdot \tau \cdot \bar{x}}{\frac{1}{\tilde{\sigma}_0^2} + n \cdot \tau}$$

Der Posterior ist ein gewichteter Durchschnitt von Prior-Mittelwert und Stichprobenmittelwert $\bar{x}$, gewichtet nach ihren Präzisionen.

\textbf{Varianz-Hyperparameter:} Der Einfluss des Priors kann durch Hyperparameter $\sigma_0^2 = \frac{1}{\alpha_0}$ eingestellt werden. Ein kleines $\sigma_0^2$ (großes $\alpha_0$) bedeutet starken Prior-Einfluss (vertrauensvoll), ein großes $\sigma_0^2$ (kleines $\alpha_0$) bedeutet schwachen Prior-Einfluss (uninformativ).

\textbf{Sequentielle Updates:} Crucially: Die Reihenfolge der Datenpunkte ist unerheblich, und Updates können sequentiell (Punkt für Punkt) oder auf einmal durchgeführt werden — das Ergebnis ist identisch (Batch vs. Online Learning).

\definition{\textbf{Normal-Inverse-Gamma (NIW) für Gauß-Verteilung mit unbekannten $\mu$ und $\sigma^2$:}}

Wenn auch die Varianz $\sigma^2$ unbekannt ist, wird ein konjugierter Prior benötigt, der sowohl über $\mu$ als auch über $\sigma^2$ faktorisiert.

Prior:
$$p(\mu, \sigma^2) = \underbrace{N(\mu \mid \mu_0, \frac{\sigma^2}{\kappa_0})}_{\text{Normal in } \mu} \times \underbrace{\text{Inv-Gamma}(\sigma^2 \mid \alpha_0, \beta_0)}_{\text{Inverse-Gamma in } \sigma^2}$$

wobei:
\begin{itemize}[leftmargin=1.5em]
    \item $\mu_0$: Prior-Mittelwert
    \item $\kappa_0$: Pseudo-Stichproben-Größe (Präzisions-Gewicht des Priors auf $\mu$)
    \item $\alpha_0, \beta_0$: Shape- und Rate-Parameter der Inverse-Gamma für $\sigma^2$
\end{itemize}

Posterior nach $n$ Beobachtungen:
$$p(\mu, \sigma^2 \mid D) = N(\mu \mid \mu_n, \frac{\sigma^2}{\kappa_n}) \times \text{Inv-Gamma}(\sigma^2 \mid \alpha_n, \beta_n)$$

mit Update-Regeln:
\begin{align*}
\kappa_n &= \kappa_0 + n \\
\mu_n &= \frac{\kappa_0 \mu_0 + n \bar{x}}{\kappa_n} \\
\alpha_n &= \alpha_0 + \frac{n}{2} \\
\beta_n &= \beta_0 + \frac{1}{2}\sum_{i=1}^{n}(x_i - \bar{x})^2 + \frac{n\kappa_0}{2\kappa_n}(\bar{x} - \mu_0)^2
\end{align*}

\textbf{Anwendung:} In der Praxis sehr häufig, wenn man ein Normalverteilungsmodell an Daten anpassen möchte, ohne Annahmen über die Varianz zu treffen.

\newpage




\subsection{Survivorship Bias und kognitive Verzerrungen}

\definition{\textbf{Survivorship Bias (Überlebendenverzerrung):}}
Eine statistische oder kognitive Verzerrung, die auftritt, wenn eine Analyse nur auf den ``Überlebenden'' oder Erfolgreichen eines Prozesses basiert, während die ``Nicht-Überlebenden'' oder Misserfolge systematisch übersehen werden. Dies führt zu fehlerhaften Schlussfolgerungen über den Prozess selbst.

\textbf{Mathematische Perspektive:} Wenn $D_{\text{obs}}$ die beobachteten (überlebenden) Datenpunkte sind und $D_{\text{missing}}$ die nicht-beobachteten (eliminieren) Datenpunkte, dann:
$$\text{Bias} = \mathbb{E}[\text{Statistik} \mid D_{\text{obs}}] \neq \mathbb{E}[\text{Statistik} \mid D_{\text{obs}} \cup D_{\text{missing}}]$$

\beispiel{\textbf{Flugzeug-Panzerfrage (2. Weltkrieg):}}

\textbf{Problem:} US-Militärs analysieren die zurückgekehrten Flugzeuge aus Kampfeinsätzen und beobachten gehäufte Treffer in bestimmten Sektoren (z.B. Tragflächen, Rumpf). Sie wollen diese Bereiche zusätzlich panzern, um künftige Flugzeuge widerstandsfähiger zu machen.

\textbf{Naïve Analyse:} ``Die meisten Treffer sind im Rumpf — panzere den Rumpf mehr!''

\textbf{Survivorship Bias:} Die ausgegebenen Flugzeuge sind bereits die ``Überlebenden''. Flugzeuge, die in kritischen Bereichen wie dem Cockpit, Motor oder Hydraulik getroffen wurden, sind \textbf{nicht} zurückgekommen — sie fehlen in der Analyse!

\textbf{Richtige Schlussfolgerung (Abraham Wald):} Die Bereiche \textbf{ohne} gehäufte Treffer (Cockpit, Motoren) sind wahrscheinlich die kritischen Punkte. Diese Flugzeuge sind nicht heimgekehrt, weil sie dort getroffen wurden. Daher sollten \textit{diese} Bereiche gepanzert werden.

\textbf{Lesson:} Überlebendenverzerrung ist eine systematische Lücke in der Datengrundlage, nicht nur ein statistisches Problem.

\beispiel{\textbf{Das Helm-Paradoxon (1. Weltkrieg):}}

\textbf{Beobachtung:} Soldaten erhielten Stahlhelme statt Stoff/Leder-Kopfbedeckungen. Nach Einführung der Metallhelme stieg die registrierte Anzahl von Kopfverletzungen in Feldspitälern um etwa das Fünffache.

\textbf{Naive Interpretation:} ``Stahlhelme sind schädlich! Sie verursachen mehr Kopfverletzungen!''

\textbf{Richtige Erklärung:} 
\begin{itemize}[leftmargin=1.5em]
    \item \textbf{Vor Helmausgabe:} Direkter Splitterbeschuß aus Schützengräben $\to$ tödliche Kopfverletzungen, Soldaten sterben sofort \textbf{nicht im Feldspital erfasst}
    \item \textbf{Nach Helmausgabe:} Gleiche Splitter treffen, aber Helm verhindert Tod; Soldaten kommen ins Feldspital mit Kopfverletzungen $\to$ \textbf{diese werden jetzt registriert}
\end{itemize}

Die Metallhelme waren nicht schädlich — sie erhöhten nur die Überlebensquote. Die scheinbare Zunahme von Kopfverletzungen war eine Messung der \textbf{Überlebenden}, nicht der Gesamtverletzten.

\definition{\textbf{Availability Bias (Verfügbarkeitsheuristik):}}
Eine verwandte Verzerrung, bei der das Gehirn/die Analyse leicht zugängliche oder auffällige Informationen übergewichtet, statt eine repräsentative Stichprobe zu verwenden.

\textbf{Beispiele:}
\begin{itemize}[leftmargin=1.5em]
    \item Schätze die Häufigkeit von Flugzeugunglücken vs. Autounfällen: Flugzeugunglücke sind präsenter in den Medien $\to$ Autounfälle sind tatsächlich häufiger
    \item Nur Kundenbewertungen von Amazon betrachtende statt alle Kunden: Bewertungen sind systematisch von stark-zufriedenen oder stark-unzufriedenen Kunden
    \item Social-Media-Analyse: Nur aktive Posts sehen, nicht stumme Mehrheit
\end{itemize}

\definition{\textbf{Confirmation Bias (Bestätigungsverzerrung):}}
Eine psychologische und statistische Bias, bei der nur Daten erfasst oder bevorzugt werden, die existierende Überzeugungen oder Erwartungen bestätigen. Dies erzeugt positive Rückkopplung und selbsterfüllende Prophezeiungen.

\beispiel{\textbf{Algorithmic Bias in der Polizeieinsatzplanung:}}

\textbf{Ausgangslage:} Ein datengetriebenes Modell soll effektiven Polizeieinsatz zur Verhinderung schwerer Verbrechen planen.

\textbf{Schritt 1 (Berechnung):} Verbrechensdichte wird basierend auf leichten und schweren Verbrechen berechnet. In historischen Daten ist Verbrechensdichte in bestimmten Stadtvierteln höher (z.B. wirtschaftlich benachteiligt).

\textbf{Schritt 2 (Verstärkte Einsätze):} Das Modell weist verstärkte Polizeipräsenz in Bereichen höherer Verbrechensdichte an.

\textbf{Schritt 3 (Zusätzliche Erfassung):} Durch mehr Polizeipräsenz vor Ort werden \textbf{zusätzlich} leichte Verbrechen entdeckt und registriert (z.B. Drogendelikte, Ordnungswidrigkeiten), nicht nur schwere Verbrechen.

\textbf{Positive Rückkopplung:} Die verstärkte Erfassung leichter Verbrechen erhöht die berechnete Verbrechensdichte in diesen Vierteln weiter $\to$ noch mehr Polizeieinsatz $\to$ noch mehr erfasste Verbrechen $\to$ Konzentration auf wenige Viertel verstärkt sich.

\textbf{Problem:} Das Modell wurde ursprünglich mit verzerrten Daten trainiert; nun verstärkt sich diese Verzerrung durch die Rückkopplung. Das System lernt, die historische Verzerrung zu replizieren und zu verstärken.

\textbf{Implikation:} Confirmation Bias ist nicht nur eine kognitiv-psychologische Verzerrung, sondern ein \textbf{strukturelles Risiko} in datengetriebenen Systemen, die sich selbst füttern (Feedback loops).

\definition{\textbf{Allgemeine Erkennungsmerkmale und Gegenmaßnahmen:}}

\textbf{Identifikation:}
\begin{itemize}[leftmargin=1.5em]
    \item \textbf{Überprüfung der Datengrundlage:} Welche Daten sind vorhanden, und warum? Welche Daten fehlen systematisch?
    \item \textbf{Selection Mechanism:} Wer oder was entscheidet, welche Fälle beobachtet werden?
    \item \textbf{Rückkopplungsschleifen:} Können Modellvorhersagen die Eingangsdaten beeinflussen?
\end{itemize}

\textbf{Gegenmaßnahmen:}
\begin{itemize}[leftmargin=1.5em]
    \item \textbf{Missing Data Analysis:} Statistiken über fehlende Datenpunkte sammeln
    \item \textbf{Sensitivity Analysis:} Wie robust sind Schlussfolgerungen unter verschiedenen Annahmen über fehlende Daten?
    \item \textbf{Domain Expertise:} Außerstatistische Quellen nutzen (Fachliteratur, Experteninterview)
    \item \textbf{Blind Randomization:} Bei möglichst vielen Schritten Zufälligkeit einführen
    \item \textbf{Feedback Loop Monitoring:} Feedback-Schleifen in ML-Systemen aktiv überwachen und begrenzen
    \item \textbf{Fairness Constraints:} Bei Einsatz von Modellen in gesellschaftskritischen Bereichen (Polizei, Kreditvergabe, Medizin) aktive Fairness-Constraints einführen
\end{itemize}

\newpage

\subsection{Wahrheitsmatrix und abgeleitete Kenngrößen}

\definition{\textbf{Confusion Matrix (Wahrheitsmatrix, Kontingenztabelle):}}
Eine tabellarische Darstellung der Leistung eines Klassifikationsmodells. Für binäre Klassifikation hat sie folgende Struktur:

\begin{center}
\begin{tabular}{|l|c|c|}
    \hline
    \textbf{} & \textbf{Actual: Negative (N)} & \textbf{Actual: Positive (P)} \\
    \hline
    \textbf{Predicted: Negative} & True Negative (TN) & False Negative (FN) \\
    \hline
    \textbf{Predicted: Positive} & False Positive (FP) & True Positive (TP) \\
    \hline
\end{tabular}
\end{center}

\textbf{Interpretationen:}
\begin{itemize}[leftmargin=1.5em]
    \item \textbf{TP (True Positive):} Modell sagt Positive voraus, Realität ist Positive (korrekt)
    \item \textbf{TN (True Negative):} Modell sagt Negative voraus, Realität ist Negative (korrekt)
    \item \textbf{FP (False Positive):} Modell sagt Positive voraus, Realität ist Negative (Fehler: ``Fehlalarm'')
    \item \textbf{FN (False Negative):} Modell sagt Negative voraus, Realität ist Positive (Fehler: ``übersehen'')
\end{itemize}

\definition{\textbf{Accuracy (Genauigkeit, Success Rate):}}
Der Anteil korrekter Vorhersagen an allen Vorhersagen:
$$\text{Accuracy} = \frac{\text{TP} + \text{TN}}{\text{TP} + \text{TN} + \text{FP} + \text{FN}}$$

\textbf{Wertebereich:} $[0, 1]$ oder $[0\%, 100\%]$. Höher ist besser.

\textbf{Limitation:} Bei unbalanciertem Datensatz (z.B. 95\% Negative, 5\% Positive) kann ein Modell, das alles als Negativ vorhersagt, eine hohe Accuracy haben, obwohl es unbrauchbar ist.

\definition{\textbf{Error Rate (Fehlerquote, Misclassification Rate):}}
Der Anteil falscher Vorhersagen:
$$\text{Error Rate} = \frac{\text{FP} + \text{FN}}{\text{TP} + \text{TN} + \text{FP} + \text{FN}} = 1 - \text{Accuracy}$$

\definition{\textbf{Sensitivity (Empfindlichkeit, Recall, True Positive Rate -- TPR):}}
Der Anteil der tatsächlich positiven Fälle, die korrekt identifiziert wurden. Misst die Fähigkeit, Positive zu erkennen:
$$\text{Sensitivity} = \text{TPR} = \frac{\text{TP}}{\text{TP} + \text{FN}} = \frac{\text{TP}}{P}$$

wobei $P = \text{TP} + \text{FN}$ die Gesamtzahl der tatsächlich positiven Fälle ist.

\textbf{Interpretation:} Wenn ein Email-Filter Sensitivity = 0.95 hat, werden 95\% aller Spam-Emails korrekt erkannt. 5\% gelangen in den Posteingang (False Negatives).

\textbf{Synonyme:} Recall, Hit Rate, Detection Rate, True Positive Rate (TPR).

\definition{\textbf{Specificity (Spezifität, True Negative Rate -- TNR):}}
Der Anteil der tatsächlich negativen Fälle, die korrekt identifiziert wurden. Misst die Fähigkeit, Negative zu erkennen:
$$\text{Specificity} = \text{TNR} = \frac{\text{TN}}{\text{TN} + \text{FP}} = \frac{\text{TN}}{N}$$

wobei $N = \text{TN} + \text{FP}$ die Gesamtzahl der tatsächlich negativen Fälle ist.

\textbf{Interpretation:} Wenn ein Email-Filter Specificity = 0.99 hat, werden 99\% aller legitimen Emails korrekt als legitim erkannt. 1\% werden fälschlicherweise als Spam markiert (False Positives).

\textbf{Synonyme:} Selectivity, True Negative Rate (TNR).

\definition{\textbf{Trade-off zwischen Sensitivity und Specificity:}}
Diese beiden Maße sind üblicherweise negativ korreliert. Durch Anheben des Klassifikations-Schwellenwerts (Threshold/Cutoff):
\begin{itemize}[leftmargin=1.5em]
    \item Höherer Threshold $\to$ weniger Positive vorhergesagt $\to$ weniger TP, aber auch weniger FP $\to$ höhere Specificity, niedrigere Sensitivity
    \item Niedrigerer Threshold $\to$ mehr Positive vorhergesagt $\to$ mehr TP, aber auch mehr FP $\to$ höhere Sensitivity, niedrigere Specificity
\end{itemize}

Die optimale Wahl hängt vom Kontext ab: Bei Krebserkennung wird höhere Sensitivity bevorzugt (lieber ein False Positive als einen Fall übersehen); bei Spam-Filterung ist höhere Specificity wünschenswert (legitime Emails sollten nicht gelöscht werden).

\definition{\textbf{Precision (Präzision, Positive Predictive Value -- PPV):}}
Wenn das Modell eine Positive vorhersagt, wie wahrscheinlich ist diese Vorhersage korrekt:
$$\text{Precision} = \text{PPV} = \frac{\text{TP}}{\text{TP} + \text{FP}}$$

\textbf{Interpretation:} Ein Email-Filter mit Precision = 0.96 bedeutet: 96\% der Emails, die als Spam markiert wurden, sind tatsächlich Spam. 4\% sind legitime Emails (False Positives).

\textbf{Unterschied zu Sensitivity:}
\begin{itemize}[leftmargin=1.5em]
    \item \textbf{Precision:} Von allen vorhergesagten Positives, wie viele sind korrekt? (Focus: Korrektheit der Positive-Vorhersagen)
    \item \textbf{Sensitivity/Recall:} Von allen echten Positives, wie viele werden erkannt? (Focus: Vollständigkeit)
\end{itemize}

\definition{\textbf{Recall (Rückruf):}}
Ein Synonym für Sensitivity. Misst, wie vollständig die positiven Fälle erfasst wurden:
$$\text{Recall} = \text{Sensitivity} = \frac{\text{TP}}{\text{TP} + \text{FN}}$$

\definition{\textbf{Trade-off zwischen Precision und Recall:}}
Diese sind ebenfalls negativ korreliert. Durch Herabsetzen des Threshold:
\begin{itemize}[leftmargin=1.5em]
    \item Mehr positive Vorhersagen $\to$ höheres Recall (mehr echte Positive gefunden), aber niedrigere Precision (mehr False Positives)
    \item Weniger positive Vorhersagen $\to$ höhere Precision (weniger False Positives), aber niedriges Recall (echte Positive übersehen)
\end{itemize}

\textbf{Anwendungsbeispiel:} In Informationsbeschaffung (Information Retrieval):
\begin{itemize}[leftmargin=1.5em]
    \item \textbf{High Recall, Low Precision:} Suchmaschine gibt viele Ergebnisse zurück (findet viele relevante Seiten), aber auch viele irrelevante. Nützlich bei Durchsuchen großer Dokumentmengen.
    \item \textbf{High Precision, Low Recall:} Suchmaschine gibt nur hochrelevante Ergebnisse zurück, könnte aber viele verpassen. Nützlich bei Suche nach seltenen/spezifischen Dokumenten.
\end{itemize}

\definition{\textbf{F1-Score (F-Measure, F-Measure):}}
Die harmonische Mittelwert von Precision und Recall, gewichtet mit gleichen Anteilen:
$$F_1 = 2 \cdot \frac{\text{Precision} \cdot \text{Recall}}{\text{Precision} + \text{Recall}} = \frac{2 \cdot \text{TP}}{2 \cdot \text{TP} + \text{FP} + \text{FN}}$$

\textbf{Eigenschaften:}
\begin{itemize}[leftmargin=1.5em]
    \item Wertebereich: $[0, 1]$; höher ist besser
    \item $F_1 = 1$ nur wenn Precision und Recall beide 1 sind (perfekt)
    \item $F_1 = 0$ wenn TP = 0
    \item Besser als einfaches Mittel $(P + R)/2$ bei unbalancierten Klassen (harmonische Mittel bestraft extrem niedrige Werte stärker)
\end{itemize}

\textbf{Generalisierung:} Der gewichtete F-Score $F_\beta$ betont Recall stärker ($\beta > 1$) oder Precision stärker ($\beta < 1$):
$$F_\beta = (1 + \beta^2) \cdot \frac{\text{Precision} \cdot \text{Recall}}{\beta^2 \cdot \text{Precision} + \text{Recall}}$$

\definition{\textbf{False Positive Rate (FPR) und False Negative Rate (FNR):}}
\begin{itemize}[leftmargin=1.5em]
    \item \textbf{FPR (False Positive Rate):} Anteil falscher Positive an allen Negativen:
    $$\text{FPR} = \frac{\text{FP}}{\text{TN} + \text{FP}} = \frac{\text{FP}}{N} = 1 - \text{Specificity}$$
    
    \item \textbf{FNR (False Negative Rate):} Anteil falscher Negative an allen Positiven:
    $$\text{FNR} = \frac{\text{FN}}{\text{TP} + \text{FN}} = \frac{\text{FN}}{P} = 1 - \text{Sensitivity}$$
\end{itemize}

\definition{\textbf{ROC Curve (Receiver Operating Characteristic):}}
Eine grafische Darstellung der Modellleistung bei unterschiedlichen Klassifikations-Schwellenwerten.

\textbf{Achsen:}
\begin{itemize}[leftmargin=1.5em]
    \item \textbf{X-Achse:} False Positive Rate (FPR) = $1 - \text{Specificity}$
    \item \textbf{Y-Achse:} True Positive Rate (TPR) = Sensitivity
\end{itemize}

\textbf{Punkte auf der Kurve:} Jeder Punkt repräsentiert eine andere Threshold-Einstellung:
\begin{itemize}[leftmargin=1.5em]
    \item \textbf{Linke untere Ecke (0, 0):} Threshold sehr hoch, fast keine Positives vorhergesagt (viele FN, wenige FP)
    \item \textbf{Rechte obere Ecke (1, 1):} Threshold sehr niedrig, fast alle Negatives vorhergesagt (wenige FN, viele FP)
    \item \textbf{Diagonale (0, 0) zu (1, 1):} Zufälliger Klassifizierer (50/50 Vorhersagen)
    \item \textbf{Linke obere Ecke (0, 1):} Idealer Klassifizierer (hoher TPR, niedriger FPR)
\end{itemize}

\textbf{Interpretation:} Je weiter die Kurve nach oben-links gewölbt ist, desto besser das Modell. Ein Modell unter der Diagonale ist schlechter als Zufall und kann durch Invertieren der Vorhersagen verbessert werden.

\definition{\textbf{AUC (Area Under the Curve):}}
Die Fläche unter der ROC-Kurve. Ein Aggregat-Maß der Modellleistung über alle Threshold-Werte.

\textbf{Wertebereich:} $[0, 1]$

\textbf{Interpretation:}
\begin{itemize}[leftmargin=1.5em]
    \item \textbf{AUC = 0.5:} Zufälliger Klassifizierer (Kurve folgt der Diagonale)
    \item \textbf{AUC = 1.0:} Perfekter Klassifizierer
    \item \textbf{AUC = 0.9--1.0:} Exzellent
    \item \textbf{AUC = 0.8--0.9:} Gut
    \item \textbf{AUC = 0.7--0.8:} Befriedigend
    \item \textbf{AUC < 0.7:} Schwach/unzureichend
\end{itemize}

\textbf{Wahrscheinlichkeitsinterpretation:} AUC ist die Wahrscheinlichkeit, dass das Modell einen zufällig ausgewählten positiven Beispiel höher rangiert als einen zufällig ausgewählten negativen Beispiel.

\definition{\textbf{Balanced Accuracy (Ausgeglichene Genauigkeit):}}
Der Durchschnitt von Sensitivity und Specificity:
$$\text{Balanced Accuracy} = \frac{\text{Sensitivity} + \text{Specificity}}{2} = \frac{\text{TPR} + \text{TNR}}{2}$$

\textbf{Verwendung:} Bei stark unbalancierten Datensätzen ist Balanced Accuracy ein besseres Maß als einfache Accuracy, da sie beide Klassen gleich gewichtet.

\textbf{Zusammenfassung der Performance-Metriken:}

\begin{table}[h]
    \centering
    \small
    \begin{tabular}{|l|l|l|l|}
        \hline
        \textbf{Metrik} & \textbf{Formel} & \textbf{Fokus} & \textbf{Best For} \\
        \hline
        Accuracy & $\frac{\text{TP+TN}}{n}$ & Gesamt-Korrektheit & Balancierte Daten \\
        \hline
        Sensitivity/Recall & $\frac{\text{TP}}{\text{TP+FN}}$ & Positive finden & Medizin, Sicherheit \\
        \hline
        Specificity & $\frac{\text{TN}}{\text{TN+FP}}$ & Negative korrekt erkennen & Spam-Filter \\
        \hline
        Precision & $\frac{\text{TP}}{\text{TP+FP}}$ & Zuverlässigkeit Positive & Klassifizierung kritisch \\
        \hline
        F1-Score & $\frac{2 \cdot P \cdot R}{P+R}$ & Balance P \& R & Unbalancierte Daten \\
        \hline
        AUC & Area under ROC & Ranking-Qualität & Allgemeiner Vergleich \\
        \hline
        Balanced Accuracy & $\frac{\text{TPR+TNR}}{2}$ & Beide Klassen & Unbalancierte Daten \\
        \hline
    \end{tabular}
\end{table}

\newpage

\subsection{Support Vector Machines}

\definition{\textbf{Support Vector Machines (SVM):}}
Ein Klassifikationsalgorithmus, der eine optimale Trennungsebene (Hyperplane) zwischen zwei oder mehr Klassen findet. Ziel ist es, den Abstand zur nächsten Trainingsinstanz zu maximieren (Maximum Margin Principle).

\definition{\textbf{Hyperplane:}}
Eine Fläche im p-dimensionalen Raum, definiert durch die Gleichung:
$$H: w^T x + b = 0$$

wobei:
\begin{itemize}[leftmargin=1.5em]
    \item $w = (w_1, w_2, \ldots, w_p)$ der Normalvektor ist (bestimmt die Orientierung)
    \item $b$ der Bias-Term ist (bestimmt die Position)
    \item $x = (x_1, x_2, \ldots, x_p)$ ein Datenpunkt ist
\end{itemize}

\textbf{Interpretationen:}
\begin{itemize}[leftmargin=1.5em]
    \item In $\mathbb{R}^2$ ist eine Hyperplane eine Linie
    \item In $\mathbb{R}^3$ ist eine Hyperplane eine Ebene
    \item Im höherdimensionalen Raum verallgemeinert sich das Konzept
\end{itemize}

\definition{\textbf{Maximum Margin Classifier:}}
Ein Klassifikator, der die Hyperplane mit dem größten Abstand (Margin) zu den nächsten Trainingsbeobachtungen findet. Dies ist nur möglich, wenn die Klassen linear separierbar sind.

\textbf{Optimierungsproblem:}
\begin{itemize}[leftmargin=1.5em]
    \item Minimiere: $\frac{1}{2}||w||^2$ (oder äquivalent: maximiere $\frac{2}{||w||}$)
    \item Unter den Bedingungen: $y_i (w^T x_i + b) \geq 1$ für alle Trainingsbeobachtungen $i$
\end{itemize}

wobei $y_i \in \{-1, +1\}$ die Klassenbezeichnungen sind.

\definition{\textbf{Decision Rule (Klassifizierungsregel):}}
Für eine neue Beobachtung $x_0$ wird die Klassifizierung vorgenommen über:
$$\text{Prediction} = \text{sign}(w^T x_0 + b)$$

\textbf{Interpretation:}
\begin{itemize}[leftmargin=1.5em]
    \item Falls $w^T x_0 + b > 0$: Vorhersage ist Klasse $+1$
    \item Falls $w^T x_0 + b < 0$: Vorhersage ist Klasse $-1$
    \item Die Entfernung von der Hyperplane ist proportional zu $|w^T x_0 + b|$
\end{itemize}

\definition{\textbf{Support Vectors:}}
Trainingsbeobachtungen, die auf den Margins (den parallelen Hyperplanes $w^T x + b = 1$ oder $w^T x + b = -1$) liegen oder näher daran sind. Diese Punkte bestimmen die Position und Orientierung der optimalen Hyperplane.

\textbf{Eigenschaft:} Beobachtungen, die weit von der Hyperplane entfernt liegen, beeinflussen die Lösung nicht.

\definition{\textbf{Lagrange-Multipliers und Lagrange-Funktion:}}
Zur Lösung des Optimierungsproblems wird die Lagrange-Funktion verwendet:
$$L(w, b, \alpha) = \frac{1}{2}||w||^2 - \sum_{i=1}^{n} \alpha_i [y_i (w^T x_i + b) - 1]$$

wobei $\alpha_i \geq 0$ die Lagrange-Multipliers sind.

\textbf{Aus den Ableitungen ergeben sich zwei wichtige Erkenntnisse:}
\begin{itemize}[leftmargin=1.5em]
    \item $w = \sum_{i=1}^{n} \alpha_i y_i x_i$: Der Gewichtsvektor ist eine Linearkombination der Trainingsbeobachtungen
    \item $\sum_{i=1}^{n} \alpha_i y_i = 0$: Die Lagrange-Multipliers müssen diese Bedingung erfüllen
\end{itemize}

\definition{\textbf{Karush-Kuhn-Tucker (KKT) Bedingungen:}}
Notwendige Bedingungen für die Optimalität der Lösung:
$$\alpha_i \cdot [y_i (w^T x_i + b) - 1] = 0 \quad \text{für alle } i$$

\textbf{Interpretation:}
\begin{itemize}[leftmargin=1.5em]
    \item Falls $\alpha_i > 0$: Die Beobachtung $i$ liegt genau auf einem der Margins ($y_i (w^T x_i + b) = 1$) → Support Vector
    \item Falls $\alpha_i = 0$: Die Beobachtung $i$ liegt nicht auf dem Margin → hat keinen Einfluss auf die Lösung
\end{itemize}

\definition{\textbf{Support Vector Classifiers (SVC):}}
Eine Verallgemeinerung der Maximum Margin Classifiers für den Fall, dass die Klassen \textbf{nicht} linear separierbar sind oder eine robustere Klassifizierung gewünscht ist.

\textbf{Kernidee:} Wir erlauben, dass einige Trainingsbeobachtungen in die Margin eindringen oder sogar falsch klassifiziert werden.

\definition{\textbf{Slack Variables (Schlupfvariablen):}}
Variablen $\xi_i \geq 0$ für jede Trainingsbeobachtung $i$, die ausdrücken, wie sehr eine Beobachtung gegen die Margin-Bedingungen verstößt:
$$\xi_i = 0: \quad \text{Beobachtung liegt auf der korrekten Seite des Margins}$$
$$0 < \xi_i \leq 1: \quad \text{Beobachtung dringt in den Margin ein, wird aber korrekt klassifiziert}$$
$$\xi_i > 1: \quad \text{Beobachtung wird falsch klassifiziert}$$

\definition{\textbf{Support Vector Classifier Optimierungsproblem:}}
\begin{itemize}[leftmargin=1.5em]
    \item Minimiere: $\frac{1}{2}||w||^2 + C \sum_{i=1}^{n} \xi_i$
    \item Unter den Bedingungen: 
    \begin{itemize}[leftmargin=1.5em]
        \item $y_i (w^T x_i + b) \geq 1 - \xi_i$ für alle $i$
        \item $\xi_i \geq 0$ für alle $i$
    \end{itemize}
\end{itemize}

\textbf{Komponenten:}
\begin{itemize}[leftmargin=1.5em]
    \item $\frac{1}{2}||w||^2$: Margin-Term (versucht, die Margin breit zu halten)
    \item $C \sum_{i=1}^{n} \xi_i$: Strafterm (bestraft Verletzungen der Constraints)
    \item $C$: Regulalisierungsparameter (Tuning-Parameter)
\end{itemize}

\definition{\textbf{Regulalisierungsparameter C:}}
Ein Hyperparameter, der den Trade-off zwischen Margin-Maximierung und Klassifizierungsfehler kontrolliert:
\begin{itemize}[leftmargin=1.5em]
    \item \textbf{Großes C:} Hohe Intoleranz gegenüber Verletzungen → enges Margin → kann zu Overfitting führen (hohe Varianz)
    \item \textbf{Kleines C:} Hohe Toleranz gegenüber Verletzungen → breites Margin → mehr Bias, weniger Varianz
    \item \textbf{C = 0:} Entspricht dem Maximum Margin Classifier (nur für linear separierbare Daten)
    \item \textbf{C $\to \infty$:} Alle Slack Variables müssen 0 sein (perfekte Klassifizierung erzwungen)
\end{itemize}

\textbf{Wahl von C:} Üblicherweise via Kreuzvalidation bestimmt.

\definition{\textbf{Kernel Trick (Kernel-Methode):}}
Eine Technik, um nichtlineare Klassifizierungsgrenzen zu erzeugen, ohne explizit in höherdimensionale Räume zu transformieren.

\textbf{Grundidee:}
\begin{enumerate}[leftmargin=1.5em]
    \item Statt die Merkmale $x_1, \ldots, x_p$ direkt zu verwenden, transformieren wir sie in einen höherdimensionalen Raum: $\phi(x)$
    \item In diesem erweiterten Raum ist die Trennung möglich linear
    \item Die Lösung hängt nur von Skalarprodukte $\phi(x_i) \cdot \phi(x_j)$ ab, nicht von der expliziten Form von $\phi$
    \item Die Kernel-Funktion $K(x_i, x_j)$ berechnet diese Skalarprodukte effizient: $K(x_i, x_j) = \phi(x_i) \cdot \phi(x_j)$
\end{enumerate}

\definition{\textbf{Gängige Kernel-Funktionen:}}

\textbf{1. Linear Kernel:}
$$K(x_i, x_j) = x_i^T x_j$$
Erzeugt eine lineare Klassifizierungsgrenze (Hyperplane). Äquivalent zu SVM ohne Kernel-Transformation.

\textbf{2. Polynomial Kernel:}
$$K(x_i, x_j) = (1 + x_i^T x_j)^d$$
wobei $d$ der Grad des Polynoms ist. Erzeugt polynomiale Klassifizierungsgrenzen.

\textbf{Beispiel:} Für $d = 2$ und 2D-Punkte $x, z \in \mathbb{R}^2$:
$$(1 + x^T z)^2 = 1 + 2x_1z_1 + 2x_2z_2 + x_1^2 z_1^2 + 2x_1z_1 x_2z_2 + x_2^2 z_2^2$$
Dies entspricht einer Transformation in den 6-dimensionalen Raum: $\phi(x) = (1, \sqrt{2}x_1, \sqrt{2}x_2, x_1^2, \sqrt{2}x_1x_2, x_2^2)^T$

\textbf{3. Radial Basis Function (RBF) Kernel:}
$$K(x_i, x_j) = \exp\left(-\gamma ||x_i - x_j||^2\right)$$
wobei $\gamma > 0$ ein Hyperparameter ist. Erzeugt komplexe, nichtlineare Klassifizierungsgrenzen und ist einer der beliebtesten Kernels.

\textbf{Interpretation:}
\begin{itemize}[leftmargin=1.5em]
    \item Misst die Ähnlichkeit zwischen zwei Punkten basierend auf ihrer euklidischen Distanz
    \item Kleine $\gamma$: Breitere Ähnlichkeitsfunktion (glatte Grenze)
    \item Große $\gamma$: Engere Ähnlichkeitsfunktion (sehr flexible, lokale Grenze)
\end{itemize}

\textbf{4. Sigmoid Kernel (Tanh Kernel):}
$$K(x_i, x_j) = \tanh(\kappa x_i^T x_j + \theta)$$
wobei $\kappa$ und $\theta$ Parameter sind. Ähnelt einer neural network activation function.

\definition{\textbf{SVM für Multi-class Klassifikation:}}
Die ursprüngliche SVM ist für binäre Klassifikation entwickelt. Für mehr als 2 Klassen gibt es zwei Hauptansätze:

\textbf{1. One-versus-One Approach (OvO, All-Pairs):}
\begin{itemize}[leftmargin=1.5em]
    \item Für $K$ Klassen werden $\binom{K}{2} = \frac{K(K-1)}{2}$ binäre SVMs trainiert (eine für jedes Paar von Klassen)
    \item Beispiel: Für 3 Klassen werden 3 SVMs trainiert (1 vs. 2, 1 vs. 3, 2 vs. 3)
    \item Bei Vorhersage: Jede SVM macht eine Vorhersage, und die Klasse mit den meisten Stimmen gewinnt (Voting)
    \item Speichereffizient für große Datenmengen, aber kann bei vielen Klassen rechenaufwändig sein
\end{itemize}

\textbf{2. One-versus-All Approach (OvA, One-versus-Rest):}
\begin{itemize}[leftmargin=1.5em]
    \item Für $K$ Klassen werden $K$ binäre SVMs trainiert (Klasse $i$ vs. alle anderen Klassen)
    \item Beispiel: Für 3 Klassen werden 3 SVMs trainiert: (1 vs. 2,3), (2 vs. 1,3), (3 vs. 1,2)
    \item Bei Vorhersage: Alle $K$ SVMs bewerten den Testpunkt, und die Klasse mit dem höchsten Konfidenz-Score wird ausgewählt
    \item Weniger Modelle zu trainieren, aber kann zu Unausgewogenheit bei der Klassifizierung führen
\end{itemize}

\definition{\textbf{Vorteile von SVM:}}
\begin{itemize}[leftmargin=1.5em]
    \item \textbf{Hohe Dimensionalität:} Funktioniert gut, auch wenn die Anzahl der Features die Anzahl der Beobachtungen übersteigt
    \item \textbf{Speichereffizient:} Nur die Support Vectors werden benötigt zur Speicherung und Vorhersage
    \item \textbf{Flexible Klassifizierungsgrenzen:} Durch verschiedene Kernels können komplexe Grenzen modelliert werden
    \item \textbf{Theoretisch fundiert:} Starke mathematische Grundlagen und Generalisierungsgüte (strukturelle Risikominimierung)
    \item \textbf{Robust:} Weniger anfällig für Overfitting durch strukturelle Regularisierung
\end{itemize}

\definition{\textbf{Nachteile von SVM:}}
\begin{itemize}[leftmargin=1.5em]
    \item \textbf{Rechenaufwand:} Training kann bei großen Datensätzen ($n > 50000$) sehr langsam sein (O($n^3$) Komplexität)
    \item \textbf{Hyperparameter-Tuning:} Wahl von $C$ und Kernel-Parametern erfordert sorgfältige Kreuzvalidation
    \item \textbf{Skalierung erforderlich:} Input-Features sollten üblicherweise normalisiert oder standardisiert werden
    \item \textbf{Interpretation schwierig:} Die Entscheidungsgrenzen sind schwer zu interpretieren, besonders mit komplexen Kernels
    \item \textbf{Klassenunausgeglichenheit:} SVM kann mit stark unausgeglichenen Klassen Schwierigkeiten haben (kann durch Gewichtung von $C$ adressiert werden)
\end{itemize}

\definition{\textbf{Vergleich: SVM vs. andere Klassifizierer}}

\begin{table}[h]
    \centering
    \small
    \begin{tabular}{|l|l|l|l|}
        \hline
        \textbf{Aspekt} & \textbf{SVM} & \textbf{Logistische Regression} & \textbf{Random Forest} \\
        \hline
        Hochdimensionalität & Sehr gut & Gut & Mäßig \\
        \hline
        Kernels möglich & Ja & Nein & Nein \\
        \hline
        Rechenaufwand & Hoch & Niedrig & Mittel \\
        \hline
        Hyperparameter & Mittel & Niedrig & Hoch \\
        \hline
        Interpretierbarkeit & Niedrig & Hoch & Mittel \\
        \hline
        Feature Engineering & Wichtig & Wichtig & Weniger wichtig \\
        \hline
        Multi-class & Extra Aufwand & Native & Native \\
        \hline
    \end{tabular}
\end{table}

\newpage

\subsection{Validierung, Underfitting und Overfitting}

\definition{\textbf{Generalisierung (Generalization):}}
Die Fähigkeit eines Modells, auf neuen, ungesehenen Daten gute Vorhersagen zu treffen. Dies ist das Endziel des maschinellen Lernens.

\textbf{Zentrale Problematik:} Wir trainieren ein Modell auf Trainingsdaten, möchten aber, dass es auf neuen Testdaten gut funktioniert.

\definition{\textbf{Training Error vs. Test Error (Generalisierungsfehler):}}

\textbf{Training Error (Trainingsfehler):}
$$\text{Train Error} = \frac{1}{n_{\text{train}}} \sum_{i=1}^{n_{\text{train}}} L(y_i, \hat{y}_i)$$

Der Fehler des Modells auf den Trainingsdaten, die zum Trainieren verwendet wurden.

\textbf{Test Error (Testfehler, Generalisierungsfehler):}
$$\text{Test Error} = \frac{1}{n_{\text{test}}} \sum_{i=1}^{n_{\text{test}}} L(y_i, \hat{y}_i)$$

Der Fehler des Modells auf neuen, ungesehenen Testdaten.

\textbf{Erkenntnis:} Training Error ist normalerweise niedriger als Test Error, da das Modell die Trainingsdaten "gesehen" hat.

\definition{\textbf{Underfitting (Unteranpassung, High Bias):}}
Das Modell ist zu simpel und kann die zugrunde liegende Struktur der Daten nicht erfassen.

\textbf{Charakteristika:}
\begin{itemize}[leftmargin=1.5em]
    \item Hoher Trainingsfehler
    \item Hoher Testfehler
    \item Das Modell macht systematische Fehler
    \item Training Error $\approx$ Test Error (beide hoch)
\end{itemize}

\textbf{Beispiele:}
\begin{itemize}[leftmargin=1.5em]
    \item Eine lineare Regression auf nicht-lineare Daten
    \item Ein Entscheidungsbaum mit maximaler Tiefe 1 (nur ein Split)
    \item Ein zu simple logistische Regression bei komplexen Klassifizierungsproblemen
\end{itemize}

\textbf{Ursachen:}
\begin{itemize}[leftmargin=1.5em]
    \item Modell hat zu wenig Kapazität (zu wenig Parameter)
    \item Unzureichende Komplexität des Modells
    \item Zu viel Regularisierung (z.B. großes $\lambda$ in Ridge Regression)
\end{itemize}

\textbf{Behebung (Lösungsstrategien):}
\begin{itemize}[leftmargin=1.5em]
    \item Komplexeres Modell wählen (z.B. von linear zu polynomial)
    \item Mehr oder bessere Features hinzufügen
    \item Regularisierungsstärke reduzieren (kleineres $\lambda$)
    \item Training länger durchführen (mehr Iterationen)
\end{itemize}

\definition{\textbf{Overfitting (Überanpassung, High Variance):}}
Das Modell ist zu komplex und passt sich an die Trainingsdaten ``zu gut'' an, einschließlich des Rauschens.

\textbf{Charakteristika:}
\begin{itemize}[leftmargin=1.5em]
    \item Sehr niedriger Trainingsfehler
    \item Hoher Testfehler
    \item Das Modell hat die Trainingsdaten memoriert, nicht generalisiert
    \item Training Error $\ll$ Test Error (großer Unterschied)
\end{itemize}

\textbf{Beispiele:}
\begin{itemize}[leftmargin=1.5em]
    \item Ein Entscheidungsbaum, der bis zur maximalen Tiefe wächst und jeden Trainingspoint auf seiner eigenen Blatt-Node hat (100\% Training Accuracy)
    \item Ein Polynom-Modell 10. Grades für 12 Trainingspunkte
    \item Ein neuronales Netzwerk mit zu vielen Hidden Units ohne Regularisierung
\end{itemize}

\textbf{Ursachen:}
\begin{itemize}[leftmargin=1.5em]
    \item Modell hat zu viel Kapazität (zu viele Parameter)
    \item Zu wenig Trainingsdaten relativ zur Modellkomplexität
    \item Zu lange Training (bei iterativen Verfahren)
    \item Unzureichende oder fehlende Regularisierung
\end{itemize}

\textbf{Behebung (Lösungsstrategien):}
\begin{itemize}[leftmargin=1.5em]
    \item Einfacheres Modell wählen (z.B. Baum mit maximaler Tiefe begrenzen)
    \item Mehr Trainingsdaten sammeln
    \item Regularisierung hinzufügen oder erhöhen (größeres $\lambda$)
    \item Early Stopping verwenden
    \item Cross-Validation zur Modellauswahl nutzen
\end{itemize}

\definition{\textbf{Bias-Variance Trade-off:}}
Ein fundamentales Konzept im Machine Learning, das die Balance zwischen Underfitting und Overfitting beschreibt.

\textbf{Bias (Verzerrung):}
$$\text{Bias} = E[\hat{f}(x)] - f(x)$$

Der systematische Fehler eines Modells. Misst, wie sehr das Modell die wahre Funktion konsistent verfehlt.
\begin{itemize}[leftmargin=1.5em]
    \item \textbf{Hoher Bias:} Modell ist zu simpel → Underfitting
    \item \textbf{Niedriger Bias:} Modell kann komplexe Muster erfassen
\end{itemize}

\textbf{Variance (Varianz):}
$$\text{Variance} = E[(\hat{f}(x) - E[\hat{f}(x)])^2]$$

Die Empfindlichkeit des Modells gegenüber Schwankungen in den Trainingsdaten. Misst, wie sehr sich das Modell ändert, wenn man die Trainingsdaten änderung.
\begin{itemize}[leftmargin=1.5em]
    \item \textbf{Hohe Varianz:} Modell ist sehr flexibel → Overfitting
    \item \textbf{Niedrige Varianz:} Modell ist stabil
\end{itemize}

\textbf{Bias-Variance Decomposition:}
$$\text{MSE} = \text{Bias}^2 + \text{Variance} + \text{Irreducible Error}$$

Der erwartete Test-Fehler zerfällt in drei Komponenten:
\begin{itemize}[leftmargin=1.5em]
    \item \textbf{Bias$^2$:} Systematischer Fehler des Modells
    \item \textbf{Variance:} Zufälliger Fehler aufgrund von Instabilität
    \item \textbf{Irreducible Error:} Rauschen in den Daten, das nicht durch ein Modell reduziert werden kann
\end{itemize}

\textbf{Trade-off-Visualisierung:}
\begin{itemize}[leftmargin=1.5em]
    \item \textbf{Sehr einfaches Modell:} Hoher Bias, niedrige Varianz → Underfitting
    \item \textbf{Mittleres Modell:} Optimaler Punkt (niedrigste Test Error)
    \item \textbf{Sehr komplexes Modell:} Niedriger Bias, hohe Varianz → Overfitting
\end{itemize}

\textbf{Praktische Konsequenz:} Das beste Modell ist nicht das komplexeste! Man muss einen Kompromiss finden.

\definition{\textbf{Validierung:}}
Der Prozess, die Leistung eines Modells auf Daten zu bewerten, die es nicht beim Training gesehen hat.

\textbf{Warum Validierung?}
\begin{itemize}[leftmargin=1.5em]
    \item Training Error gibt keine Aussage über Generalisierung
    \item Wir brauchen unabhängige Daten, um das echte Generalisierungspotential zu schätzen
    \item Validierung hilft, Underfitting vs. Overfitting zu diagnostizieren
\end{itemize}

\definition{\textbf{Hold-out Validation (Train-Test Split):}}
Die Daten werden in zwei disjunkte Teile aufgeteilt:
\begin{itemize}[leftmargin=1.5em]
    \item \textbf{Trainingsset:} Üblicherweise 60-80\% der Daten, zum Trainieren des Modells
    \item \textbf{Testset:} Üblicherweise 20-40\% der Daten, zum Evaluieren der Generalisierung
\end{itemize}

\textbf{Vorteile:}
\begin{itemize}[leftmargin=1.5em]
    \item Einfach zu implementieren
    \item Schnell recheneffizient
\end{itemize}

\textbf{Nachteile:}
\begin{itemize}[leftmargin=1.5em]
    \item Höhere Varianz der Schätzung (besonders bei kleinen Datensätzen)
    \item Wir "verschwenden" Daten, die nicht zum Training verwendet werden
\end{itemize}

\definition{\textbf{Cross-Validation (Kreuzvalidation):}}
Eine robustere Validierungsmethode, die die Daten mehrfach aufspaltet, um stabilere Leistungsschätzungen zu erhalten.

\textbf{K-Fold Cross-Validation:}
\begin{enumerate}[leftmargin=1.5em]
    \item Teile die Daten in $k$ Folds (Partitionen) auf (üblicherweise $k = 5$ oder $10$)
    \item Für jede Fold $i = 1, \ldots, k$:
    \begin{enumerate}[leftmargin=1.5em]
        \item Trainiere das Modell auf den $k-1$ anderen Folds
        \item Evaluiere auf Fold $i$
    \end{enumerate}
    \item Berechne den Durchschnittsfehler über alle $k$ Evaluierungen
\end{enumerate}

$$\text{CV Error} = \frac{1}{k} \sum_{i=1}^{k} \text{Error}_i$$

\textbf{Vorteile:}
\begin{itemize}[leftmargin=1.5em]
    \item Stabiler und aussagekräftiger als Hold-out Validation
    \item Nutzt alle Daten zum Training (nur in verschiedenen Kombinationen)
    \item Liefert Konfidenzintervalle/Standardabweichung
\end{itemize}

\textbf{Spezialfälle:}
\begin{itemize}[leftmargin=1.5em]
    \item \textbf{Leave-One-Out CV (LOOCV):} $k = n$ (Extremfall, sehr rechenintensiv)
    \item \textbf{Stratified K-Fold:} Für imbalancierte Klassifikation, behält Klassenverteilung in jedem Fold
    \item \textbf{Time Series CV:} Für Zeitreihen, respektiert zeitliche Ordnung
\end{itemize}

\definition{\textbf{Regularisierung (Regularization):}}
Techniken zur Reduzierung von Overfitting durch Vereinfachung oder Bestrafung des Modells.

\textbf{L1 Regularisierung (Lasso):}
$$\text{Cost} = \text{MSE} + \lambda \sum_{j=1}^{p} |w_j|$$

Bestraft die Summe der absoluten Koeffizienten. Kann manche Koeffizienten exakt auf 0 setzen (Feature Selection).

\textbf{L2 Regularisierung (Ridge):}
$$\text{Cost} = \text{MSE} + \lambda \sum_{j=1}^{p} w_j^2$$

Bestraft die Summe der quadrierten Koeffizienten. Verkleinert große Koeffizienten, aber setzt sie nicht auf 0.

\textbf{Elastic Net:}
$$\text{Cost} = \text{MSE} + \lambda_1 \sum_{j=1}^{p} |w_j| + \lambda_2 \sum_{j=1}^{p} w_j^2$$

Kombination von L1 und L2, vereinigt die Eigenschaften beider.

\textbf{Hyperparameter $\lambda$ (Regularisierungsstärke):}
\begin{itemize}[leftmargin=1.5em]
    \item \textbf{$\lambda = 0$:} Keine Regularisierung → höhere Varianz, Overfitting möglich
    \item \textbf{Großes $\lambda$:} Starke Regularisierung → höherer Bias, Underfitting möglich
    \item \textbf{Optimales $\lambda$:} Durch Cross-Validation bestimmt
\end{itemize}

\definition{\textbf{Early Stopping:}}
Ein Regularisierungstechnik für iterative Trainingsalgorithmen (z.B. Gradient Descent, Neuronale Netze).

\textbf{Idee:}
\begin{enumerate}[leftmargin=1.5em]
    \item Während des Trainings wird nach jeder Epoche/Iteration der Validierungsfehler berechnet
    \item Solange Validierungsfehler sinkt, weiter trainieren
    \item Sobald Validierungsfehler zu steigen beginnt (Overfitting-Zeichen), Training stoppen
\end{enumerate}

\textbf{Praktische Implementierung:}
\begin{itemize}[leftmargin=1.5em]
    \item Führe separate Validierungsset (z.B. 10-20\% der Trainingsdaten)
    \item Monitore Validierungsfehler
    \item Speichere das beste Modell (minimaler Validierungsfehler)
    \item Nach N Epochen ohne Verbesserung (Patience), stop
\end{itemize}

\textbf{Vorteile:}
\begin{itemize}[leftmargin=1.5em]
    \item Verhindert Overfitting bei iterativen Verfahren
    \item Automatische, datengetriebene Modellauswahl
    \item Keine zusätzliche Hyperparameter-Tuning erforderlich
\end{itemize}

\definition{\textbf{Diagnostische Plots: Learning Curves:}}
Plots von Training Error und Test Error über die Modellkomplexität (oder Trainingszeit), um Underfitting vs. Overfitting zu diagnostizieren.

\textbf{Szenarien:}

\textbf{1. Underfitting (beide Kurven hoch, nahe beieinander):}
\begin{itemize}[leftmargin=1.5em]
    \item Training Error hoch, Test Error hoch
    \item Kurven konvergieren zu hohem Fehler
    \item Lösung: Komplexeres Modell oder bessere Features
\end{itemize}

\textbf{2. Overfitting (großer Gap zwischen Kurven):}
\begin{itemize}[leftmargin=1.5em]
    \item Training Error niedrig, Test Error hoch
    \item Kurven trennen sich
    \item Lösung: Regularisierung, mehr Daten, einfacheres Modell
\end{itemize}

\textbf{3. Optimales Modell (Kurven nah beieinander und beide niedrig):}
\begin{itemize}[leftmargin=1.5em]
    \item Training Error und Test Error beide niedrig
    \item Kleine Lücke zwischen den Kurven
    \item Das Modell generalisiert gut
\end{itemize}

\definition{\textbf{Auswirkung von Datenmenge auf Overfitting:}}
Ein größerer Trainningsdatensatz kann Overfitting reduzieren:
\begin{itemize}[leftmargin=1.5em]
    \item Mit wenigen Daten: Kleines Modell ist besser (Overfitting leicht)
    \item Mit vielen Daten: Komplexere Modelle werden möglich (mehr Daten können die Komplexität rechtfertigen)
\end{itemize}

\textbf{Praktische Empfehlung:} ``Collect more data'' ist oft die beste Lösung für Overfitting, wenn möglich.

\definition{\textbf{Model Complexity vs. Generalisierung:}}

\begin{table}[h]
    \centering
    \small
    \begin{tabular}{|l|l|l|l|}
        \hline
        \textbf{Situation} & \textbf{Bias} & \textbf{Variance} & \textbf{Action} \\
        \hline
        Underfitting & Hoch & Niedrig & Modell komplexer machen \\
        \hline
        Optimal & Mittel & Mittel & Fertig! \\
        \hline
        Overfitting & Niedrig & Hoch & Regularisierung/Daten \\
        \hline
        High Bias + High Var & Hoch & Hoch & Datenqualität prüfen \\
        \hline
    \end{tabular}
\end{table}

\definition{\textbf{Zusammenfassung: Checkliste für gute Generalisierung:}}

\begin{itemize}[leftmargin=1.5em]
    \item  Trainings- und Testsets verwenden (nicht nur Trainingsdaten evaluieren)
    \item  Cross-Validation zur stabilen Leistungsschätzung durchführen
    \item  Learning Curves plotten zur Diagnose
    \item  Modellkomplexität anpassen
    \item  Regularisierung verwenden (L1, L2, oder Early Stopping)
    \item  Hyperpameter via Cross-Validation tunen
    \item  Mehr Trainingsdaten sammeln, wenn möglich
    \item  Feature Engineering / Feature Selection durchführen
\end{itemize}

\newpage

\subsection{Zeitreihenanalyse}

\definition{\textbf{Zeitreihe (Time Series):}}
Eine Sequenz von Beobachtungen $y_1, y_2, \ldots, y_T$ eines Merkmals, gemessen zu verschiedenen Zeitpunkten $t = 1, 2, \ldots, T$ in fester zeitlicher Ordnung.

\textbf{Charakteristiken:}
\begin{itemize}[leftmargin=1.5em]
    \item \textbf{Feste zeitliche Ordnung:} Die Reihenfolge ist essentiell und nicht vertauschbar
    \item \textbf{Abhängigkeit:} Beobachtungen sind nicht unabhängig, sondern zeitlich abhängig
    \item \textbf{Univariat:} Eine Variable pro Zeitpunkt (im Gegensatz zu multivariat)
    \item \textbf{Äquidistant:} Zeitpunkte sind üblicherweise diskret und in regelmäßigen Abständen (z.B. täglich, monatlich, jährlich)
\end{itemize}

\textbf{Typische Beispiele:}
\begin{itemize}[leftmargin=1.5em]
    \item Temperaturmessungen (stündlich, täglich)
    \item Aktienkurse
    \item Flugpassagierzahlen (monatlich)
    \item Verkehrsunfallzahlen
\end{itemize}

\definition{\textbf{Stochastischer Prozess:}}
Eine Folge $(Y_t)_{t \in T}$ von Zufallsvariablen $Y_t$, wobei $t$ aus einer geordneten Indexmenge $T$ (Zeit) stammt.

\textbf{Typen nach Zeitstruktur:}
\begin{itemize}[leftmargin=1.5em]
    \item \textbf{Zeitdiskret:} $T = \mathbb{N}$ oder $\mathbb{Z}$ (Standardfall in angewandter Zeitreihenanalyse)
    \item \textbf{Zeitstetig:} $T = [0, \infty)$ oder $\mathbb{R}$ (seltener in praktischen Anwendungen)
\end{itemize}

\textbf{Beziehung:} Eine Zeitreihe $(y_1, \ldots, y_T)$ ist eine Realisierung eines stochastischen Prozesses $(Y_t)$.

\definition{\textbf{Stationarität (Stationarity):}}
Eine Zeitreihe ist stationär, wenn sich ihre statistische Struktur über die Zeit nicht systematisch verändert.

\textbf{Mathematisch (strikte Stationarität):}
Die gemeinsame Verteilung von $(Y_t, Y_{t+1}, \ldots, Y_{t+k})$ ist identisch mit der Verteilung von $(Y_{t+h}, Y_{t+1+h}, \ldots, Y_{t+k+h})$ für alle $t, k, h$.

\textbf{Praktisch (schwache/zweite-Ordnung Stationarität):}
\begin{itemize}[leftmargin=1.5em]
    \item Konstanter Mittelwert: $E(Y_t) = \mu$ für alle $t$
    \item Konstante Varianz: $\text{Var}(Y_t) = \sigma^2$ für alle $t$
    \item Autokovarianzen hängen nur vom Lag ab: $\text{Cov}(Y_t, Y_{t+h})$ ist unabhängig von $t$
\end{itemize}

\textbf{Wichtigkeit:} Viele Zeitreihenmodelle (z.B. ARMA) setzen Stationarität voraus!

\definition{\textbf{Autokorrelationsfunktion (ACF):}}
Misst die Korrelation der Zeitreihe mit ihren verzögerten Versionen (Lags).

\textbf{Empirische Autokorrelation zum Lag $\omega$:}
$$r_\omega = \frac{\sum_{t=\omega+1}^{T} (y_t - \bar{y})(y_{t-\omega} - \bar{y})}{\sum_{t=1}^{T} (y_t - \bar{y})^2}, \quad -1 \leq r_\omega \leq 1$$

\textbf{Eigenschaften:}
\begin{itemize}[leftmargin=1.5em]
    \item $r_0 = 1$ (Korrelation mit sich selbst)
    \item $r_\omega = r_{-\omega}$ (symmetrisch)
    \item $|r_\omega|$ misst die Stärke der zeitlichen Abhängigkeit
    \item Vorzeichen gibt die Richtung an (positiv/negativ)
\end{itemize}

\textbf{Grafische Darstellung:} Das \textbf{Autokorrelogramm} ist ein Plot von $r_\omega$ gegen den Lag $\omega$, mit Vertrauensbändern (üblicherweise $\pm 1.96 / \sqrt{T}$).

\textbf{Interpretation:}
\begin{itemize}[leftmargin=1.5em]
    \item Schnell abfallende ACF: Kurzzeitabhängigkeit, stationäre Reihe
    \item Langsam abfallende ACF: Langzeitabhängigkeit, nicht-stationäre Reihe
    \item Saisonale Muster: Spitzen bei Lags entsprechend der Saisonperiode
\end{itemize}

\definition{\textbf{Partielle Autokorrelationsfunktion (PACF):}}
Die Korrelation zwischen $Y_t$ und $Y_{t-\omega}$, nachdem der Einfluss der dazwischenliegenden Lags entfernt wurde.

\textbf{Bedeutung:} PACF hilft bei der Identifikation der Ordnung eines autoregressiven Prozesses (wie viele Verzögerungen relevant sind).

\definition{\textbf{Komponenten einer Zeitreihe:}}
Eine Zeitreihe wird üblicherweise in folgende Hauptkomponenten zerlegt:

\textbf{1. Trend ($\mu_t$):}
Langfristige systematische Veränderung des mittleren Niveaus. Die Zeitreihe steigt oder fällt über längere Zeit.

\textbf{2. Saisonalität ($s_t$):}
Regelmäßige Schwankungen mit einer festen Periode $p$ (z.B. $p = 12$ für monatliche Daten, $p = 4$ für Quartalsdaten). Beispiele: höhere Eisverkäufe im Sommer, höhere Heizkosten im Winter.

\textbf{3. Restkomponente ($\varepsilon_t$):}
Nicht erklärbare irregulöre Schwankungen oder Rauschen. Annahme: $E(\varepsilon_t) = 0$, $\text{Var}(\varepsilon_t) = \sigma_\varepsilon^2$.

\definition{\textbf{Klassisches Komponentenmodell:}}

\textbf{Additives Modell:}
$$y_t = \mu_t + s_t + \varepsilon_t$$

Komponenten überlagern sich additiv. Verwendet wenn die Saisonalschwankungen konstant sind (absolut).

\textbf{Multiplikatives Modell:}
$$y_t = \mu_t \cdot s_t \cdot \varepsilon_t$$

Komponenten multiplizieren sich. Verwendet wenn die Saisonalschwankungen proportional zum Trend sind (relativ).

\textbf{Transformation:} Das multiplikative Modell kann durch Logarithmieren in ein additives transformiert werden:
$$\log(y_t) = \log(\mu_t) + \log(s_t) + \log(\varepsilon_t)$$

\definition{\textbf{Glättung (Smoothing): Simple Moving Average (SMA):}}
Ein Filter zur Elimination von irregulären Schwankungen und Schätzung des lokalen Trends.

\textbf{Formula für Fenster der Breite $2q+1$:}
$$\hat{\mu}_t = \frac{1}{2q+1} \sum_{i=-q}^{q} y_{t+i}, \quad t = q+1, \ldots, T-q$$

\textbf{Beispiel:} Mit $q = 2$ (Fenster der Breite 5):
$$\hat{\mu}_t = \frac{1}{5}(y_{t-2} + y_{t-1} + y_t + y_{t+1} + y_{t+2})$$

\textbf{Eigenschaften:}
\begin{itemize}[leftmargin=1.5em]
    \item Linearer Filter (lineare Transformation der Zeitreihe)
    \item Symmetrisches Zeitfenster
    \item Größeres $q$ $\Rightarrow$ stärkere Glättung, aber weniger Datenpunkte
    \item Grenzen: $q$ Beobachtungen am Anfang und Ende gehen verloren (NA)
\end{itemize}

\definition{\textbf{Trendbereinigung und Saisonbereinigung:}}

\textbf{Trendbereinigung:}
$$y_t^{\text{detrend}} = y_t - \hat{\mu}_t$$
Entfernt die Trendkomponente.

\textbf{Saisonbereinigung:}
$$y_t^{\text{deseason}} = y_t - (\hat{\mu}_t + \hat{s}_t)$$
Entfernt Trend und Saison, bleibt nur Restkomponente.

\textbf{Praktische Schritte zur Schätzung von Saisonkomponenten:}
\begin{enumerate}[leftmargin=1.5em]
    \item Schätze Trend $\hat{\mu}_t$ via SMA (mit Fensterbreite $p = $ Saisonperiode)
    \item Berechne trendbereinigt: $y_t - \hat{\mu}_t$
    \item Schätze Saisonkomponenten $\hat{s}_k$ als Durchschnitt über alle Jahre für Saison $k$
    \item Zentriere die Saisonkomponenten: $\hat{s}_k - \overline{\hat{s}}$
\end{enumerate}

\definition{\textbf{White Noise Prozess:}}
Ein stochastischer Prozess mit statistisch unabhängigen und identisch verteilten Zufallsvariablen:
$$Y_t \sim \text{IID}(\mu, \sigma^2)$$

\textbf{Eigenschaften:}
\begin{itemize}[leftmargin=1.5em]
    \item $E(Y_t) = \mu$ (konstant)
    \item $\text{Var}(Y_t) = \sigma^2$ (konstant)
    \item ACF: $r(\omega) = \begin{cases} 1 & \text{wenn } \omega = 0 \\ 0 & \text{wenn } \omega \neq 0 \end{cases}$
    \item Kein zeitlicher Zusammenhang
\end{itemize}

\textbf{Interpretation:} Reines Rauschen, keine vorhersagbare Struktur.

\definition{\textbf{Autoregressive Prozesse (AR):}}
Ein autoregressiver Prozess der Ordnung $p$ ist definiert als:
$$Y_t = \phi_1 Y_{t-1} + \phi_2 Y_{t-2} + \cdots + \phi_p Y_{t-p} + \varepsilon_t, \quad \text{AR}(p)$$

wobei $\varepsilon_t$ White Noise ist.

\textbf{Interpretation:} Der aktuelle Wert hängt von den $p$ vorherigen Werten und einem Rauschterm ab (wie eine Regression auf Verzögerungen).

\textbf{PACF-Eigenschaft:} Die PACF eines AR($p$)-Prozesses bricht nach Lag $p$ ab (wird 0).

\definition{\textbf{Moving Average Prozesse (MA):}}
Ein Moving Average Prozess der Ordnung $q$ ist definiert als:
$$Y_t = \varepsilon_t + \theta_1 \varepsilon_{t-1} + \theta_2 \varepsilon_{t-2} + \cdots + \theta_q \varepsilon_{t-q}, \quad \text{MA}(q)$$

wobei $\varepsilon_t$ White Noise ist.

\textbf{Interpretation:} Der aktuelle Wert ist eine gewichtete Summe aus aktuellen und vergangenen Rauschschocks.

\textbf{ACF-Eigenschaft:} Die ACF eines MA($q$)-Prozesses bricht nach Lag $q$ ab.

\definition{\textbf{ARMA Prozesse (Autoregressive Moving Average):}}
Eine Kombination von AR und MA:
$$Y_t = \phi_1 Y_{t-1} + \cdots + \phi_p Y_{t-p} + \varepsilon_t + \theta_1 \varepsilon_{t-1} + \cdots + \theta_q \varepsilon_{t-q}, \quad \text{ARMA}(p, q)$$

\textbf{Modellspezifikation:} Wird durch zwei Parameter definiert:
\begin{itemize}[leftmargin=1.5em]
    \item $p$: AR-Ordnung (bestimmt via PACF)
    \item $q$: MA-Ordnung (bestimmt via ACF)
\end{itemize}

\textbf{Voraussetzung:} Die Zeitreihe muss stationär sein.

\definition{\textbf{ARIMA Prozesse (Autoregressive Integrated Moving Average):}}
Für nicht-stationäre Zeitreihen mit Trend:
$$\text{ARIMA}(p, d, q): \quad \Delta^d Y_t = \phi_1 \Delta^d Y_{t-1} + \cdots + \phi_p \Delta^d Y_{t-p} + \varepsilon_t + \theta_1 \varepsilon_{t-1} + \cdots + \theta_q \varepsilon_{t-q}$$

wobei $\Delta^d Y_t$ die $d$-fach differenzierte Zeitreihe ist.

\textbf{Parameter:}
\begin{itemize}[leftmargin=1.5em]
    \item $p$: AR-Ordnung
    \item $d$: Differenzierungsordnung (üblicherweise 0, 1, oder 2)
    \item $q$: MA-Ordnung
\end{itemize}

\textbf{Interpretation:} Differenzieren macht nicht-stationäre Reihen stationär, danach wird ARMA angewendet.

\definition{\textbf{SARIMA Prozesse (Seasonal ARIMA):}}
Erweiterung für Zeitreihen mit saisonalen Mustern:
$$\text{SARIMA}(p, d, q) \times (P, D, Q)_s$$

wobei:
\begin{itemize}[leftmargin=1.5em]
    \item $(p, d, q)$: Nicht-saisonale Komponenten
    \item $(P, D, Q)$: Saisonale Komponenten
    \item $s$: Saisonperiode (z.B. 12 für monatliche Daten)
\end{itemize}

\textbf{Anwendung:} Für Daten mit regelmäßigen saisonalen Schwankungen (z.B. Eisverkäufe, Heizkosten).

\definition{\textbf{Modellidentifikation (Ordnungswahl):}}
Faustregel zur Bestimmung von $(p, q)$ aus Daten:

\textbf{ACF/PACF-Faustregel:}
\begin{itemize}[leftmargin=1.5em]
    \item \textbf{PACF schneidet bei Lag $p$ ab} $\Rightarrow$ AR($p$)-Prozess oder als $(p, d, q)$ in ARIMA
    \item \textbf{ACF schneidet bei Lag $q$ ab} $\Rightarrow$ MA($q$)-Prozess oder als $(p, d, q)$ in ARIMA
\end{itemize}

\textbf{Praxis:}
\begin{enumerate}[leftmargin=1.5em]
    \item Visualisiere die Zeitreihe → teste auf Trend und Saisonalität
    \item Differenziere wenn nötig ($d = 1$ oder $2$)
    \item Berechne ACF und PACF der differenzierten Reihe
    \item Nutze Faustregel zur Schätzung von $p$ und $q$
    \item Vergleiche Modelle via AIC/BIC-Kriterien
    \item Überprüfe Residuen auf White Noise
\end{enumerate}

\definition{\textbf{Residuenanalyse (Diagnostic Checking):}}
Nach der Modellanpassung müssen die Residuen überprüft werden:

\textbf{Anforderungen für gute Residuen:}
\begin{itemize}[leftmargin=1.5em]
    \item White Noise (keine zeitliche Abhängigkeit)
    \item Konstante Varianz (Homoskedastizität)
    \item Normalverteilung (besonders für Vorhersageintervalle)
    \item Keine systematischen Muster
\end{itemize}

\textbf{Diagnostische Tests:}
\begin{itemize}[leftmargin=1.5em]
    \item \textbf{ACF/PACF der Residuen:} Sollte White Noise ähneln (Werte nahe 0)
    \item \textbf{Ljung-Box-Test:} Testet auf serielle Korrelation (Nullhypothese: White Noise)
    \item \textbf{Shapiro-Wilk-Test:} Testet auf Normalverteilung
    \item \textbf{Time Plot:} Visualisierung auf Muster
\end{itemize}

\definition{\textbf{Zusammenfassung: Workflow der Zeitreihenanalyse:}}

\begin{enumerate}[leftmargin=1.5em]
    \item \textbf{Exploration:} Visualisiere Zeitreihe, identifiziere Trend, Saison, Ausreißer
    \item \textbf{Transformation:} Trend- und Saisonbereinigung falls nötig (SMA, Differenzieren)
    \item \textbf{Stationaritätsprüfung:} ACF/PACF, Unit Root Tests (ADF-Test)
    \item \textbf{Modellwahl:} ACF/PACF Analyse, Bestimmung $(p, d, q)$
    \item \textbf{Modellanpassung:} Schätzung der Parameter (Maximum Likelihood, etc.)
    \item \textbf{Diagnostik:} Residuenanalyse, Tests auf White Noise
    \item \textbf{Vorhersage:} Prognosen für zukünftige Werte mit Konfidenzintervallen
\end{enumerate}

\textbf{Software:} In R: `arima(x, order = c(p, d, q), seasonal = list(order = c(P, D, Q), period = s))`

\newpage

\subsection{Importance Sampling und MCMC}

\definition{\textbf{Monte Carlo Integration:}}
Numerische Methode zur Schätzung von Integralen (und Erwartungswerten) durch Stichprobenziehen.

\textbf{Grundidee:} Um einen Erwartungswert $E_p[f(x)] = \int f(x) p(x) dx$ zu schätzen:
\begin{enumerate}[leftmargin=1.5em]
    \item Ziehe Stichproben $x_1, x_2, \ldots, x_n$ aus der Verteilung $p(x)$
    \item Berechne $\hat{E}_p[f(x)] = \frac{1}{n} \sum_{i=1}^{n} f(x_i)$
\end{enumerate}

\textbf{Eigenschaft:} Nach dem Gesetz der großen Zahlen: $\hat{E}_p[f(x)] \to E_p[f(x)]$ wenn $n \to \infty$.

\textbf{Vorteil:} Funktioniert auch für hochdimensionale und komplexe Verteilungen.

\definition{\textbf{Sampling-Problem:}}
In vielen praktischen Anwendungen (z.B. Bayesianische Inferenz) ist es schwierig, direkt aus der Zielverteilung $p(x)$ zu samplen:

\textbf{Beispiel:} Posterior-Verteilung in Bayesianischen Modellen:
$$p(\theta | \text{data}) = \frac{p(\text{data} | \theta) p(\theta)}{p(\text{data})}$$

Die Normalisierungskonstante $p(\text{data})$ ist oft intractable.

\definition{\textbf{Importance Sampling:}}
Eine Methode zur Monte Carlo Schätzung von Erwartungswerten, wenn direkte Stichprobenziehung aus $p(x)$ nicht möglich ist.

\textbf{Grundidee:}
\begin{enumerate}[leftmargin=1.5em]
    \item Definiere eine leicht zu sample-nde \textbf{Proposal-Verteilung} $q(x)$ (auch ``Importance Density'' genannt)
    \item Ziehe Stichproben $x_1, \ldots, x_n$ aus $q(x)$ (nicht aus $p(x)$!)
    \item Gewichte die Stichproben um die Unterschiede zwischen $p$ und $q$ zu korrigieren
\end{enumerate}

\textbf{Mathematik:}
$$E_p[f(x)] = \int f(x) p(x) dx = \int f(x) \frac{p(x)}{q(x)} q(x) dx = E_q\left[f(x) \frac{p(x)}{q(x)}\right]$$

\textbf{Importance Weights (Importance Gewichte):}
$$w(x) = \frac{p(x)}{q(x)}$$

Messen, um wie viel wichtiger eine Stichprobe aus $q$ im Kontext von $p$ ist.

\definition{\textbf{Importance Sampling Algorithmus:}}

\begin{enumerate}[leftmargin=1.5em]
    \item Wähle Proposal-Verteilung $q(x)$
    \item Für $i = 1, \ldots, n$:
    \begin{enumerate}[leftmargin=1.5em]
        \item Ziehe $x_i \sim q(x)$
        \item Berechne Gewicht $w_i = \frac{p(x_i)}{q(x_i)}$
    \end{enumerate}
    \item Normalisiere Gewichte: $\tilde{w}_i = \frac{w_i}{\sum_{j=1}^{n} w_j}$
    \item Schätze Erwartungswert: $\hat{E}_p[f(x)] = \sum_{i=1}^{n} \tilde{w}_i f(x_i)$
\end{enumerate}

\textbf{Schätzung der Normalisierungskonstante:}
$$\hat{Z} = \frac{1}{n} \sum_{i=1}^{n} w_i = \frac{1}{n} \sum_{i=1}^{n} \frac{p(x_i)}{q(x_i)}$$

Dies ist nützlich für Modellvergleiche (Bayes Factors).

\definition{\textbf{Anforderungen an die Proposal-Verteilung $q(x)$:}}

\textbf{Effektive Proposal:}
\begin{itemize}[leftmargin=1.5em]
    \item $q(x) > 0$ überall wo $p(x) > 0$ (Support Bedingung: wichtig!)
    \item $q(x)$ sollte $p(x)$ gut approximieren (minimale Gewichtsvarianz)
    \item Leicht zu samplen und Wahrscheinlichkeiten zu berechnen
\end{itemize}

\textbf{Problem bei schlechter Wahl:}
\begin{itemize}[leftmargin=1.5em]
    \item Viele Stichproben haben sehr kleine oder sehr große Gewichte
    \item ``Degeneracy Problem'': Wenige Stichproben dominieren die Schätzung
    \item Hohe Varianz der Schätzung
\end{itemize}

\definition{\textbf{Effektive Stichprobengröße (Effective Sample Size -- ESS):}}
Ein Maß für die Anzahl der ``effektiven'' unabhängigen Stichproben:
$$\text{ESS} = \frac{\left(\sum_{i=1}^{n} w_i\right)^2}{\sum_{i=1}^{n} w_i^2}$$

\textbf{Interpretation:}
\begin{itemize}[leftmargin=1.5em]
    \item $\text{ESS} = n$: Perfekte Gewichte (uniforme Verteilung)
    \item $\text{ESS} \ll n$: Schlechte Proposal-Verteilung (Degeneracy)
    \item Ratio $\text{ESS}/n$ ist ein Effizienzmaß
\end{itemize}

\definition{\textbf{Limitation von Importance Sampling:}}
In hochdimensionalen Problemen funktioniert einfaches Importance Sampling oft schlecht:
\begin{itemize}[leftmargin=1.5em]
    \item Gewichte werden extrem variabel
    \item ESS kann dramatisch sinken
    \item Viele Stichproben haben verschwindend kleine Gewichte
\end{itemize}

\textbf{Lösung:} Markov Chain Monte Carlo (MCMC) Methoden.

\definition{\textbf{Markov Chains:}}
Eine Folge von Zufallsvariablen $(X_t)_{t=0,1,2,\ldots}$ mit der \textbf{Markov-Eigenschaft}: Die zukünftige Verteilung hängt nur vom aktuellen Zustand ab, nicht von der Geschichte:
$$P(X_{t+1} | X_t, X_{t-1}, \ldots, X_0) = P(X_{t+1} | X_t)$$

\textbf{Übergangswahrscheinlichkeit:}
$$T(x' | x) = P(X_{t+1} = x' | X_t = x)$$

Beschreibt die Wahrscheinlichkeit vom Zustand $x$ in den Zustand $x'$ zu gehen.

\definition{\textbf{Ergodische Markov Chains:}}
Eine Markov Chain ist ergodisch, wenn sie:
\begin{itemize}[leftmargin=1.5em]
    \item \textbf{Irreduzibel:} Jeden Zustand kann man von jedem anderen Zustand erreichen
    \item \textbf{Aperiodisch:} Keine periodischen Oszillationen
\end{itemize}

\textbf{Wichtigste Eigenschaft:} Eine ergodische Markov Chain konvergiert zu ihrer stationären Verteilung (unabhängig vom Startwert).

\definition{\textbf{Stationäre Verteilung:}}
Eine Verteilung $\pi(x)$ ist stationär für eine Markov Chain mit Übergängen $T(x' | x)$, wenn:
$$\pi(x') = \int \pi(x) T(x' | x) dx$$

\textbf{Interpretation:} Wenn die Chain lang genug läuft, ist die Randverteilung $\pi(x)$.

\definition{\textbf{MCMC (Markov Chain Monte Carlo):}}
Ein Algorithmus zur Stichprobenziehung aus einer komplexen Verteilung $p(x)$, indem man eine Markov Chain konstruiert, deren stationäre Verteilung $p(x)$ ist.

\textbf{Strategie:}
\begin{enumerate}[leftmargin=1.5em]
    \item Konstruiere Übergangswahrscheinlichkeiten $T(x' | x)$ so dass die stationäre Verteilung $p(x)$ ist
    \item Starte die Chain an einem beliebigen Punkt
    \item Lasse die Chain viele Iterationen laufen (bis Konvergenz)
    \item Die späteren Stichproben sind (ungefähr) aus $p(x)$
\end{enumerate}

\definition{\textbf{Detailed Balance Bedingung:}}
Eine hinreichende Bedingung für stationäre Verteilung $\pi(x)$:
$$\pi(x) T(x' | x) = \pi(x') T(x | x')$$

\textbf{Interpretation:} Die Wahrscheinlichkeit, von $x$ nach $x'$ zu gehen, multipliziert mit $\pi(x)$, gleicht die umgekehrte Richtung aus.

\textbf{Wichtig:} Viele populäre MCMC-Methoden nutzen diese Bedingung zur Konstruktion.

\definition{\textbf{Metropolis-Hastings Algorithmus:}}
Ein Algorithmus zur Konstruktion von MCMC-Übergängen mit Detail Balance.

\textbf{Algorithmus (pro Iteration $t$):}

\begin{enumerate}[leftmargin=1.5em]
    \item Starte mit Zustand $x_t$
    \item \textbf{Proposal:} Ziehe Kandidat $x' \sim q(x' | x_t)$ aus Proposal-Verteilung
    \item \textbf{Acceptance Ratio:} Berechne
    $$\alpha(x', x_t) = \min\left(1, \frac{p(x') q(x_t | x')}{p(x_t) q(x' | x_t)}\right)$$
    \item \textbf{Accept/Reject:} 
    \begin{itemize}[leftmargin=1.5em]
        \item Mit Wahrscheinlichkeit $\alpha$: Akzeptiere $x_{t+1} = x'$
        \item Mit Wahrscheinlichkeit $1 - \alpha$: Verwerfe $x_{t+1} = x_t$ (bleibe)
    \end{itemize}
\end{enumerate}

\textbf{Wichtige Spezialfälle:}

\textbf{1. Random Walk Metropolis:}
$$q(x' | x_t) = q(x' - x_t) \quad \text{(symmetrisch)}$$

Dann vereinfacht sich $\alpha = \min(1, \frac{p(x')}{p(x_t)})$.

\textbf{2. Independent Metropolis:}
$$q(x' | x_t) = q(x') \quad \text{(unabhängig von aktueller Position)}$$

\definition{\textbf{Gibbs Sampling:}}
Ein spezieller MCMC-Algorithmus für multivariate Verteilungen mit leicht zu sampling-nden Conditional Distributions.

\textbf{Idee:} Für $X = (X_1, \ldots, X_p)$ sample iterativ aus den Conditional Distributions $p(x_j | x_{-j})$ wobei $x_{-j}$ alle anderen Komponenten sind.

\textbf{Algorithmus (pro Iteration $t$):}

\begin{enumerate}[leftmargin=1.5em]
    \item Starte mit $(x_1^{(t)}, \ldots, x_p^{(t)})$
    \item Für $j = 1, \ldots, p$:
    \begin{enumerate}[leftmargin=1.5em]
        \item Sample $x_j^{(t+1)} \sim p(x_j | x_1^{(t+1)}, \ldots, x_{j-1}^{(t+1)}, x_{j+1}^{(t)}, \ldots, x_p^{(t)})$
    \end{enumerate}
\end{enumerate}

\textbf{Vorteil:} Höhere Akzeptanzrate als Random Walk Metropolis (keine Rejections).

\textbf{Gibbs ist ein Spezialfall von Metropolis-Hastings:} Mit $\alpha = 1$ immer.

\definition{\textbf{Burn-in Phase:}}
Die anfänglichen Iterationen der MCMC Chain, die verworfen werden.

\textbf{Grund:} Die Chain muss aus dem Startzustand zur Zielverteilung konvergieren. Anfängliche Stichproben sind nicht von $p(x)$ und verzerren die Schätzung.

\textbf{Praktische Richtlinie:}
\begin{itemize}[leftmargin=1.5em]
    \item Typischerweise: Burn-in = $\frac{1}{3}$ bis $\frac{1}{2}$ der Gesamtiterationen
    \item Größeres Burn-in für schlecht mixing chains
    \item Kann durch Trace Plots diagnostiziert werden
\end{itemize}

\definition{\textbf{Thinning (Ausdünnung):}}
Nur jede $k$-te Stichprobe aus der MCMC Chain behalten.

\textbf{Gründe:}
\begin{itemize}[leftmargin=1.5em]
    \item \textbf{Speichereffizienz:} Weniger Stichproben speichern
    \item \textbf{Autokorrelation reduzieren:} Aufeinanderfolgende Stichproben sind stark korreliert
\end{itemize}

\textbf{Nachteil:} Verliert Information, führt zu höherer Varianz.

\textbf{Moderne Ansicht:} Thinning ist oft nicht notwendig; besser, alle Stichproben zu behalten und Autokorrelation in der Varianzschätzung zu berücksichtigen.

\definition{\textbf{Konvergenzdiagnose:}}
Tests zur Überprüfung, ob die MCMC Chain zur stationären Verteilung konvergiert hat.

\textbf{1. Visual Inspection (Trace Plots):}
\begin{itemize}[leftmargin=1.5em]
    \item Plot der Kettensamples über Iterationen
    \item Gut konvergent: Schnelle Oszillationen um konstantes Niveau (stationär)
    \item Schlecht: Trends, Drifts, oder periodische Muster
\end{itemize}

\textbf{2. Autocorrelation Function (ACF):}
$$\rho_h = \frac{\text{Cov}(X_t, X_{t+h})}{\text{Var}(X_t)}$$

\textbf{Interpretation:} Schnell abfallende ACF deutet auf gutes Mixing hin (unabhängigere Stichproben).

\textbf{3. Gelman-Rubin Diagnostic ($\hat{R}$, Potential Scale Reduction Factor):}

Läuft mehrere unabhängige Chains mit verschiedenen Startwerten. Vergleicht Within-Chain und Between-Chain Varianz:
$$\hat{R} = \sqrt{\frac{(\hat{n}-1) W/\hat{n} + \hat{B}/\hat{n}}{W}}$$

wobei $W$ = Within-Chain Varianz, $\hat{B}$ = Between-Chain Varianz.

\textbf{Regel:} $\hat{R} < 1.1$ deutet auf Konvergenz hin. $\hat{R}$ nahe 1.0 ist ideal.

\textbf{4. Effective Sample Size (ESS):}
$$\text{ESS} = \frac{N}{1 + 2 \sum_{h=1}^{\infty} \rho_h}$$

Berücksichtigt Autokorrelation. Höhere ESS ist besser.

\definition{\textbf{Varianzschätzung in MCMC:}}
Da Stichproben autokorreliert sind, ist die naive Varianzschätzung zu optimistisch:
$$\widehat{\text{Var}}_{\text{naive}} = \frac{1}{N} \sum_{i=1}^{N} (x_i - \bar{x})^2$$

\textbf{Korrigierte Varianzschätzung (mit Autokorrelation):}
$$\widehat{\text{Var}}_{\text{korrigiert}} = \widehat{\text{Var}}_{\text{naive}} \cdot \left(1 + 2 \sum_{h=1}^{M} \rho_h \right)$$

wobei $M$ ein Cutoff-Lag ist (z.B. bis ACF signifikant abfällt).

\definition{\textbf{Ising-Modell (Anwendungsbeispiel):}}
Ein klassisches Modell aus der statistischen Physik zur Beschreibung von Magnetismus.

\textbf{Modell:}
\begin{itemize}[leftmargin=1.5em]
    \item Gitter von ``Spins'' $s_i \in \{-1, +1\}$
    \item Energie: $E(s) = -J \sum_{\langle i,j \rangle} s_i s_j - h \sum_i s_i$
    \begin{itemize}[leftmargin=1.5em]
        \item $J$: Kopplung zwischen benachbarten Spins
        \item $h$: Äußeres magnetisches Feld
    \end{itemize}
    \item Boltzmann-Verteilung: $p(s | \beta) = \frac{1}{Z(\beta)} \exp(-\beta E(s))$
    \begin{itemize}[leftmargin=1.5em]
        \item $\beta = 1/(k_B T)$: Inverse Temperatur
        \item $Z(\beta)$: Partition function (Normalisierungskonstante)
    \end{itemize}
\end{itemize}

\textbf{Anwendung von Gibbs Sampling:}
\begin{enumerate}[leftmargin=1.5em]
    \item Conditional Distribution für Spin $i$ ist analytisch verfügbar:
    $$p(s_i | s_{-i}) = \frac{1}{1 + \exp(-2\beta s_i (J \sum_{\langle j \rangle} s_j + h))}$$
    \item Iterativ für jeden Spin samplen
    \item Gleichgewichtsstichproben approx. aus $p(s | \beta)$
\end{enumerate}

\textbf{Interessante Physik:}
\begin{itemize}[leftmargin=1.5em]
    \item \textbf{Phasenübergang:} Bei kritischer Temperatur $T_c$ (in 2D: $\beta_c = \frac{1}{2}\ln(1+\sqrt{2})$) gibt es einen Phasenübergang
    \item Unter $T_c$: System ``magnetisiert'' sich (Spins richten sich aus)
    \item Über $T_c$: System ist ungeordnet (Spins zufällig)
\end{itemize}

\definition{\textbf{Vergleich: Importance Sampling vs. MCMC}}

\begin{table}[h]
    \centering
    \small
    \begin{tabular}{|l|l|l|}
        \hline
        \textbf{Aspekt} & \textbf{Importance Sampling} & \textbf{MCMC} \\
        \hline
        Direkter Sampling & Ja, aus $q$ & Nein, Markov Chain \\
        \hline
        Unabhängige Samples & Ja & Nein (autokorreliert) \\
        \hline
        Proposal-Wahl & Kritisch & Flexibler \\
        \hline
        Hochdimensional & Schwierig & Besser \\
        \hline
        Gewichte & Nötig & Nicht nötig \\
        \hline
        Konvergenz-Diagnose & Einfach (ESS) & Komplex (Gelman-Rubin, etc.) \\
        \hline
        Rechenzeit & Schnell & Langsamer (Burn-in etc.) \\
        \hline
        Batch-Processing & Gut geeignet & Sequentiell \\
        \hline
    \end{tabular}
\end{table}

\definition{\textbf{Praktische Empfehlungen:}}

\textbf{Wann Importance Sampling verwenden:}
\begin{itemize}[leftmargin=1.5em]
    \item Niedrige bis mittlere Dimensionen ($p < 20$)
    \item Wenn gute Proposal-Verteilung verfügbar
    \item Schnelle Schätzung der Normalisierungskonstante nötig
\end{itemize}

\textbf{Wann MCMC verwenden:}
\begin{itemize}[leftmargin=1.5em]
    \item Hohe Dimensionen ($p > 20$)
    \item Komplexe Posterior-Verteilungen in Bayesianischen Modellen
    \item Wenn Conditional Distributions leicht zu samplen sind (Gibbs)
\end{itemize}

\textbf{Hybrid-Ansätze:}
\begin{itemize}[leftmargin=1.5em]
    \item \textbf{Adaptive Importance Sampling:} Proposal-Verteilung wird iterativ angepasst
    \item \textbf{SMC (Sequential Monte Carlo):} Sequentielle Importance Sampling mit Resampling
    \item \textbf{Variational Inference:} Schnellere Alternative zu MCMC (aber Approximation)
\end{itemize}
\newpage